\documentclass[11pt,reqno]{amsart}
\usepackage[T1]{fontenc}
\usepackage{lmodern}
\usepackage{amsmath,amssymb,amsthm,mathtools}
\usepackage[margin=1.04in]{geometry}
\usepackage{microtype,booktabs,longtable,array,enumitem}
\usepackage{xcolor,listings}
\usepackage[colorlinks=true,linkcolor=blue!45!black,citecolor=blue!45!black,
  urlcolor=blue!45!black,breaklinks=true]{hyperref}
\numberwithin{equation}{section}
\newtheorem{theorem}{Theorem}[section]
\newtheorem{proposition}[theorem]{Proposition}
\newtheorem{lemma}[theorem]{Lemma}
\newtheorem{corollary}[theorem]{Corollary}
\theoremstyle{definition}
\newtheorem{definition}[theorem]{Definition}
\newtheorem{remark}[theorem]{Remark}
\newcommand{\Q}{\mathbb Q}
\newcommand{\Z}{\mathbb Z}
\newcommand{\R}{\mathbb R}
\newcommand{\C}{\mathbb C}
\newcommand{\F}{\mathbb F}
\newcommand{\OO}{\mathcal O}
\newcommand{\OK}{\mathcal O_K}
\newcommand{\Tr}{\operatorname{Tr}}
\newcommand{\Nm}{\operatorname{N}}
\newcommand{\eK}{\mathrm e_K}
\newcommand{\ee}{\mathrm e}
\newcommand{\TT}{\mathbb T}
\newcommand{\D}{\mathfrak D}
\newcommand{\A}{\mathcal A}
\newcommand{\Mstar}{\mathfrak M_*}
\newcommand{\mstar}{\mathfrak m_*}
\setcounter{tocdepth}{1}
\setlist{itemsep=3pt,topsep=5pt}
\lstset{language=Python,basicstyle=\ttfamily\scriptsize,
  columns=fullflexible,keepspaces=true,breaklines=true,
  showstringspaces=false,upquote=true,tabsize=4,
  frame=single,rulecolor=\color{black!20},framerule=.3pt,
  xleftmargin=3pt,xrightmargin=3pt,aboveskip=9pt,belowskip=9pt}
\emergencystretch=1.2em
\title[Uniform transfers for fifth through eighth powers]
{Higher powers in the Waring--Siegel problem:\\
uniform transfers for $k=5,6,7,8$}
\author{Research note for the Waring--Siegel project}
\date{24 November 2026}
\hypersetup{pdftitle={Higher powers in the Waring-Siegel problem: uniform transfers for k=5,6,7,8},
 pdfauthor={Waring-Siegel project}}
\begin{document}
\begin{abstract}
We prove $G_{\mathbb K}(8)\le38$ for every number field $K$, alongside the earlier
transfers $G_{\mathbb K}(5)\le17$, $G_{\mathbb K}(6)\le24$, and $G_{\mathbb K}(7)\le31$.
The decisive field moment has excess $0.117157055860<2/17$.
The analytic and local proofs and exact verification programs are included.
Only the summand counts are uniform over fields; the ninth-power transfer
remains unfinished.
\end{abstract}
\maketitle
\raggedbottom
{\small\tableofcontents}
\clearpage
\section{Conclusions and precise scope}
\label{sec:conclusions}

This expanded manuscript completes the eighth-power transfer and incorporates
the previously established fifth-, sixth-, and seventh-power arguments.
The numerical bounds in all four cases are independent of the number field.

\begin{theorem}\label{thm:main-transfer}
For every number field $K$,
\[
 \boxed{G_{\mathbb K}(5)\le17,\quad G_{\mathbb K}(6)\le24,\quad G_{\mathbb K}(7)\le31,\quad G_{\mathbb K}(8)\le38.}
\]
The number of summands is independent of the degree, signature, ideal class
group, and local ramification indices.  Constants in the norm bounds and
the sufficiently-large target threshold may depend on $K$.
\end{theorem}

We use the Waring--Siegel convention
\[
 J_{K,k}=\langle x^k:x\in\mathcal O_K\rangle_{\mathbb Z}.
\]
The assertion $G_{\mathbb K}(k)\le N$ means that there are constants $C,N_0$, depending
on $K,k,N$, such that every totally positive $\nu\in J_{K,k}$ with
$|N_{K/\mathbb Q}\nu|>N_0$ has a representation
\[
 \nu=x_1^k+\cdots+x_N^k,\qquad x_i\in\mathcal O_K,
 \qquad \sigma(x_i)\ge0\quad(\sigma:K\hookrightarrow\mathbb R),
\]
with
\[
 |N_{K/\mathbb Q}x_i|\le C|N_{K/\mathbb Q}\nu|^{1/k}.
\]
Zero summands may be appended.  In fields without real embeddings the sign
condition is empty.  The target is required to belong to the actual additive
order $J_{K,k}$; this is not a claim about all algebraic integers irrespective
of their local congruences.

\begin{center}
\begin{tabular}{@{}c c c p{6.5cm}@{}}
\toprule
$k$&Rational target&Primitive local bound&Outcome over arbitrary $K$\\
\midrule
5&17&10&Transfer proved; proof retained here.\\
6&24&12&Transfer proved; proof retained here.\\
7&31&14&Transfer proved; proof retained here.\\
8&38&35&Transfer completed in this continuation, including all three previously outstanding field arguments.\\
9&46&30&Unfinished. The simpler auxiliary route tested here does not cross the required threshold.\\
\bottomrule
\end{tabular}
\end{center}

The published rational comparison bounds are $17,24,31,39,47$, respectively
\cite{VW1995,VW1994,Wooley2016}.  The values $38$ and $46$ are the stronger
targets in the project's rational manuscript \cite{ProjectG89}.
No equality or historical priority claim is asserted.

\subsection{The new eighth-power input}

The additional field estimate is
\[
 \int_{\mathbb T_K}|g_8(\alpha)|^{34}\,d\alpha
       \ll X^{26+\Delta_{34}+\epsilon},\qquad
 \Delta_{34}=\frac{5857852793}{50000000000}=0.117157055860.
\]
Its exact margin is $4-34\Delta_{34}=0.016660100760>0$.
The eighteen degree-eight seed rows and twenty-eight degree-sixteen
coefficient seed rows are established over the field itself. Twelve
rounds of the fixed quadratic terminal policy then give the displayed
excess. The proof does not import a rational smooth-moment table as a
number-field theorem.

The new material begins in Section~\ref{sec:field-seed-moments};
the major-arc argument is Section~\ref{sec:aux-major-24}, the quadratic
terminal is Section~\ref{sec:eighth-quadratic-terminal}, and the global
conclusion is Section~\ref{sec:eighth-global}. The complete new exact
program is Appendix~\ref{app:eighth-field}. The earlier analytic and
local proofs, including the local bound $35$ for eighth powers, are
retained so that the present manuscript can be read independently of
the previous PDF.

\subsection{The earlier fifth-power input}


The fifth-power estimate is the joint moment
\begin{equation}\label{eq:quintic-overview}
 \int_{\mathbb T_K}|f_5(\alpha)|^2|g_5(\alpha)|^{14}\,d\alpha
          \ll X^{11.0764+\epsilon},
 \qquad X=P^{[K:\mathbb Q]}.
\end{equation}
Here $f_5$ is an ordinary fifth-power sum dominating the support of the
ideal-smooth sum $g_5$.  Its excess $191/2500<1/13$ gives a minor-arc
saving $17/35000$ in the seventeen-variable representation integral.
The auxiliary degree-ten input is also proved over the field itself:
\[
 \int_{\mathbb T_K}|g_{10}(\alpha)|^{44}\,d\alpha
                         \ll X^{34.227+\epsilon}.
\]
The two exact certificates have six and twenty rows, respectively.
They start from elementary moments and trivial bounds, so no rational
smooth-moment table is presumed valid over $K$.

The retained fifth-power argument is in Sections~\ref{sec:tenth-moment}--\ref{sec:k5-global}.
The required coefficient sampling is developed in
Section~\ref{sec:coefficient-sampling-common}.
The exact programs are in Appendices~\ref{app:quintic} and~\ref{app:decic}.
Sections~\ref{sec:restricted-field} and~\ref{sec:quad-sampling-field}
record the common field estimates now used in the completed eighth-power transfer.
The earlier proofs, local results, and verification programs are retained,
so this manuscript can be read without the previous PDF.

The sixth- and seventh-power inputs remain
\[
 \int|g_6|^{20}\ll X^{14.22+\epsilon},\qquad
 \int|g_7|^{28}\ll X^{21.11765+\epsilon}.
\]
Principal blocks of prime ideals give the efficient-differencing congruences
without a class-number-one assumption.  Small supports at all embeddings
provide the needed nonvanishing even at complex places.  On the major arcs,
a smooth root is replaced by a product of a fixed number of smaller ordinary
roots, supplying a direct Poisson main term.  Each such product contributes
one Waring summand.

\subsection{Theorems imported and arguments adapted}

Wooley's number-field Vinogradov mean-value theorem
\cite[Theorem~15.1 and Corollary~15.3]{Wooley2019} is used directly.
The coefficient-uniform affine curve theorem of
Binyamini--Cluckers--Kato \cite[Corollary~12]{BCK2025} and Cluckers'
local-field oscillatory-integral estimate \cite[Proposition~3.3]{Cluckers2011}
are likewise already stated over the required fields.

Efficient differencing and the auxiliary-moment processes
\cite{WooleyThesis,Wooley1992,VW1995,VW2000} are techniques whose needed
field versions are proved here.  In particular, the sharper fifth-power
coefficient estimate corresponding to \cite[Lemma~4.3]{VW1995} requires
four conductor identifications and a local congruence argument.
Its use in the cubic terminal also requires support completion with uniform
linear and degree-five twists.  These are parts of the proof, rather than
consequences of merely changing the notation in a rational theorem.

The restricted-height mean of \cite[Theorem~4.2]{BW2023} is also transferred
explicitly.  Here an ideal factor is allowed to be nonprincipal; the quotient
root lies in one of finitely many fixed integral ideals, with its corresponding
coefficient torus.  It supplies one of the inputs to the eighth-power theorem; the other
seed, terminal, local, and global arguments are given separately below.

The ordinary localization and pruning arguments are reproduced from the
earlier analytic work \cite{ProjectAnalytic}.  Their shifting method is
related to \cite[Theorem~2.1]{Wooley2012}, with its corrigendum
\cite{WooleyCorr}; the latter paper is not quoted as a ready-made
number-field localization theorem.
\section{The ordinary analytic estimates over \texorpdfstring{$K$}{K}}
\label{sec:ordinary-bundle}

Fix $k\ge3$ and a number field $K$ of degree $n$.
Write $V=K\otimes_{\mathbb Q}\mathbb R$, $\mathcal O=\mathcal O_K$,
and $\mathfrak D$ for the different.
At an archimedean place put $d_v=1$ or $2$, according as it is real or complex.
Use the trace character $e_K(z)=\exp(2\pi i\operatorname{Tr}z)$ and probability
Haar measure on $\mathbb T_K=V/\mathfrak D^{-1}$.
Normalize Lebesgue measure on $V$ so that $\mathcal O$ has covolume one,
and take its trace-dual measure on the coefficient space.
For a fixed smooth compactly supported weight put
\[
 X=P^n,\qquad f(\theta)=\sum_{x\in\mathcal O}w(x/P)e_K(\theta x^k).
\]
All assertions are uniform in bounded smooth families with common compact support.
For the major-arc assertions, take product weights supported away from zero
at every embedding.  Such weights suffice for the representation integrals.

For $\gamma\in K/\mathfrak D^{-1}$ define
\[
 \mathfrak a_\gamma=\{a\in\mathcal O:a\gamma\in\mathfrak D^{-1}\},
 \quad Q_\gamma=N\mathfrak a_\gamma,
 \quad A_P(\beta)=\prod_{v\mid\infty}\max\{1,P^k|\beta_v|\}^{d_v}.
\]
Let $\mathfrak M_*$ be the union of
$Q_\gamma A_P(\theta-\gamma)\le C_KX$, let
$\mathfrak m_*=\mathbb T_K\setminus\mathfrak M_*$, and let
$\mathfrak N_c$ be the union with right side $X^c$, where $0<c<1$.
The fixed constant $C_K>0$ can be chosen once for all the estimates below.

The ordinary bundle used in the proof is
\begin{align}
 \int_{\mathfrak m_*}|f|^{k(k+1)}
    &\ll X^{k(k+1)-k-1+\epsilon},\label{eq:ordinary-critical-bundle}\\
 \int_{\mathfrak M_*}|f|^p&\ll X^{p-k}\qquad(p>2k),\label{eq:ordinary-major-bundle}\\
 \int_{\mathfrak M_*\setminus\mathfrak N_c}|f|^p
    &\ll X^{p-k-c(p/k-2)/2}\qquad(p>2k).
       \label{eq:ordinary-pruning-bundle}
\end{align}
We give the field arguments, including the ingredients that are used elsewhere
for the auxiliary cubic sums.

\subsection{Packing and close pairs}

If $\mathfrak b$ is a nonzero fractional ideal, then in any translated
archimedean box of radii $R_v$,
\begin{equation}\label{eq:ordinary-packing}
 \#(\mathfrak b\cap\text{box})\ll_K
       1+\frac{\prod_vR_v^{d_v}}{N\mathfrak b}.
\end{equation}
Indeed, every nonzero difference $z\in\mathfrak b$ has
$|Nz|\ge N\mathfrak b$.  After dividing the $v$th coordinate by $R_v$,
the maximum-distance separation is at least
$(N\mathfrak b/\prod R_v^{d_v})^{1/n}$.
Packing disjoint balls in a fixed bounded box proves
\eqref{eq:ordinary-packing}, also for highly unequal radii and translated boxes.

The following close-pair observation is particularly useful.  Suppose
$|y_v|\le C P$, $a\in\mathfrak D^{-1}$, and
\[
 |\theta_vy_v+\eta_v-a_v|\le Z_v\qquad(v\mid\infty).
\]
If $X\prod Z_v^{d_v}\le c_K$ for a sufficiently small fixed $c_K$, then either
there is at most one pair $(y,a)$ or all the pairs lie on one affine $K$-line
$a=a_0+\gamma(y-y_0)$, and their number is at most
\begin{equation}\label{eq:ordinary-line}
 O_K\left(1+Q_\gamma^{-1}
       \prod_v\min\{P,Z_v/|\theta_v-\gamma_v|\}^{d_v}\right).
\end{equation}
A quotient by zero means infinity.
To prove this, take differences $(u_i,b_i)$ from one pair.
The determinants $b_1u_2-b_2u_1\in\mathfrak D^{-1}$ have conjugates
$O(PZ_v)$, hence norm less than $N\mathfrak D^{-1}$, so vanish.
Distinct pairs cannot have the same $y$ by the same norm argument.
They therefore have one slope $\gamma=b/u\in K$.
Their $y$-coordinates lie in a coset of $\mathfrak a_\gamma$ and in the
intersection of the original box with boxes of radii
$O(Z_v/|\theta_v-\gamma_v|)$; apply \eqref{eq:ordinary-packing}.
Without the small-volume hypothesis the elementary bound is
\begin{equation}\label{eq:ordinary-line-trivial}
 \#\{(y,a)\}\ll_K X\left(1+\prod_v Z_v^{d_v}\right),
\end{equation}
by counting $a$ separately for each $y$.

\subsection{The shifted kernel and the critical minor-arc moment}

Choose a smooth compactly supported $v_0$, put $H=P^{k-1}$, and define
$D_H(t)=\sum_{h\in\mathcal O}v_0(h/H)e_K(-th)$.
Poisson summation gives, for every fixed large $L$,
\[
 |D_H(t)|\ll H^n\sum_{a\in\mathfrak D^{-1}}
             \prod_v(1+H|t_v-a_v|)^{-L}.
\]
Let $Y$ be a fixed coordinate box of algebraic integers with side comparable
to $P$, so $\#Y\asymp_KX$.  Uniformly in $\eta\in V$ and
$\theta\in\mathfrak m_*$,
\begin{equation}\label{eq:ordinary-shifted-kernel}
 \frac1{\#Y}\sum_{y\in Y}|D_H(k\theta y+\eta)|\ll X^{k-2}.
\end{equation}
Here is the full counting argument.
Decompose the displayed Poisson majorant into boxes
$Z_v=z_v/H$, $z_v\ge1$ dyadic, and write $B=\prod z_v^{d_v}$.
If $B\le c_KX^{k-2}$, \eqref{eq:ordinary-line} applies to the close-pair
count $N_B$.  Apart from a single pair it gives
\[
 \frac{N_B}{B}\ll 1+
 \frac{X}{Q_\gamma
   \prod_v\max\{z_v,P^k|k\theta_v-\gamma_v|\}^{d_v}}
 \ll_K1.
\]
The last inequality follows from $\theta\notin\mathfrak M_*$:
the conductor of $\gamma/k$ is at most $k^nQ_\gamma$, and
$A_P(\theta-\gamma/k)\le A_P(k\theta-\gamma)$.
If $B>c_KX^{k-2}$, \eqref{eq:ordinary-line-trivial} gives
$N_B/B\ll X/B+X^{2-k}\ll1$, since $k\ge3$.
Summing the convergent dyadic weights and using $H^n/X=X^{k-2}$ proves
\eqref{eq:ordinary-shifted-kernel}.

Define the full Vinogradov sum over a fixed enlarged coordinate box by
\[
 F(\boldsymbol\alpha)=\sum_{x\in\mathcal B(CP)}
    e_K(\alpha_1x+\cdots+\alpha_kx^k),\qquad
 J_{u,k}(CP)=\int_{\mathbb T_K^k}|F(\boldsymbol\alpha)|^{2u}
                 d\boldsymbol\alpha.
\]
For any positive integer $u$, the shifted-kernel argument gives
\begin{equation}\label{eq:ordinary-localized-general}
 \int_{\mathfrak m_*}|f(\theta)|^{2u}d\theta
       \ll X^{k(k-1)/2-1}J_{u,k}(CP).
\end{equation}
For completeness, insert coefficients of $x,\ldots,x^{k-2}$ by periodic
Fourier kernels which equal one on every possible difference of sums of these
powers.  Their suprema cost
$X^{1+\cdots+(k-2)}=X^{(k-1)(k-2)/2}$.
After integrating those coefficients, expansion imposes
$\sigma_j=\sum_{i\le u}(x_i^j-y_i^j)=0$ for $1\le j\le k-2$.
A common translation by $z\in Y$ preserves these equations and gives
\[
 \sigma_k(\mathbf x-z,\mathbf y-z)
      =\sigma_k(\mathbf x,\mathbf y)-kz\sigma_{k-1}(\mathbf x,\mathbf y).
\]
Fourier-complete each shifted weight using
$w((x-z)/P)=\int\widehat w(\xi)e_K(\xi(x-z)/P)\,d\xi$.
The Fourier $L^1$ norm is fixed, and all $x$ now lie in one enlarged box.
Introduce $\alpha_{k-1}$ by $D_H$ with $v_0$ equal to one on all possible
$P^{1-k}\sigma_{k-1}$.
Averaging over $z$ gives the kernel
$D_H(\alpha_{k-1}+kz\theta)$ times factors of modulus one.
Apply \eqref{eq:ordinary-shifted-kernel}.
The remaining product contains $2u$ full Vinogradov sums, each with only a
linear coefficient shifted by a Fourier frequency.
Extend the coefficient domain to $\mathbb T_K^k$ and use H\"older with
exponent $2u$.  Translation of $\alpha_1$ preserves Haar measure, so every
factor is $J_{u,k}(CP)$.
The total cost is
$X^{(k-1)(k-2)/2+k-2}=X^{k(k-1)/2-1}$, proving
\eqref{eq:ordinary-localized-general}.

The ready-made deep input is Wooley's number-field theorem
\cite[Theorem~15.1]{Wooley2019}, applied to $t,t^2,\ldots,t^k$,
whose Wronskian is nonzero.  In its volume normalization it gives
\[
 J_{k(k+1)/2,k}(CP)\ll X^{k(k+1)/2+\epsilon}.
\]
Substitution in \eqref{eq:ordinary-localized-general} proves
\eqref{eq:ordinary-critical-bundle}.  Thus the moment theorem itself is
already stated over number fields; the localization to these product arcs is
the argument just proved, rather than a rational minor-arc theorem assumed to transfer.

\subsection{The classical Hua seeds}

The elementary field Hua estimates needed for the smooth-moment iteration are
\begin{equation}\label{eq:ordinary-Hua}
 \int_{\mathbb T_K}|f|^{2^r}\ll X^{2^r-r+\epsilon}
       \qquad(1\le r\le k).
\end{equation}
Here is a short proof, to keep the seeds independent of rational-only inputs.
By orthogonality a bounded weighted moment is at most a fixed multiple of
the unweighted moment on a containing box.
The case $r=1$ follows from $x^k=y^k$.
Assume the result for $r=J<k$.
Repeated Cauchy on the additive lattice gives
$|f|^{2^J}\ll X^{2^J-J-1}\sum_{x,\mathbf h}
e_K(\theta\Delta_{h_1}\cdots\Delta_{h_J}x^k)$, with the usual vertex-support
restrictions; the aggregate on the right is nonnegative.
Multiply by $|f|^{2^J}$ and integrate.
The equation so counted is
$\sum_{i\le2^{J-1}}(u_i^k-v_i^k)=\Delta_{h_1}\cdots\Delta_{h_J}x^k$.
If its common value is zero, the induction hypothesis supplies
$O(X^{2^J-J+\epsilon})$ choices of $\mathbf u,\mathbf v$, and there are
$O(X^J)$ choices of $x,\mathbf h$: a zero shift is immediate, while otherwise
the nonconstant polynomial in $x$ has boundedly many roots.
If its common value is nonzero, every $h_i$ divides it, so ideal divisors and
bounded-height units give $O(X^\epsilon)$ choices for all shifts, followed by
boundedly many $x$.  Thus the total count is $O(X^{2^J+\epsilon})$ in both
cases.  Multiplying by the prefactor proves the next Hua moment.

\subsection{Complete sums, Poisson summation, and pruning}

If $\psi$ is a primitive additive character modulo an integral ideal
$\mathfrak a$, $a$ is a unit modulo $\mathfrak a$, and $b$ is arbitrary, then
\begin{equation}\label{eq:ordinary-complete}
 \left|\sum_{x\bmod\mathfrak a}\psi(ax^k+bx)\right|
       \ll_{K,k,\epsilon}(N\mathfrak a)^{1-1/k+\epsilon}.
\end{equation}
We recall a uniform local proof because the linear twists are required by Poisson.
At a good prime of residue cardinality $q$, not dividing $2k(k-1)$,
write the modulus as $\pi^e$ and split into shells $v(x)=t$.
The deep shell $t\ge\lceil e/k\rceil$ has normalized size $O(q^{-e/k})$.
Outside it, put $r=v(b)$.  If $r\ne(k-1)t$, the residual polynomial on units
has unit derivative, so all shells with residual conductor at least two vanish.
Conductor one is bounded by the finite-field Weil theorem
\cite{Weil1948}.
There is at most one critical shell $r=(k-1)t$; its stationary points are
simple and its normalized size is
$O_k(q^{-t-(e-kt)/2})=O_k(q^{-e/k})$.
Thus the constant is uniform at every good prime.
At each of the finitely many bad primes, the local van der Corput theorem
\cite[Proposition~3.3]{Cluckers2011}, applied to the $k$th derivative $k!a$,
gives $O_{K_\mathfrak p,k}(q^{-e/k})$ uniformly in the integral lower coefficients.
Chinese remaindering and the ideal-divisor bound absorb the product of
good-prime constants into $(N\mathfrak a)^\epsilon$, proving
\eqref{eq:ordinary-complete}.

For $w=\prod_vw_v$ supported in archimedean annuli, let
\[
 I_v(t,u)=\int_{K_v}w_v(z)e_{K_v}(tz^k-uz)\,dz.
\]
Uniform stationary phase and integration by parts give, for every large fixed $L$,
\begin{equation}\label{eq:ordinary-stationary}
 |I_v(t,u)|\ll(1+|t|)^{-d_v/2}
                (1+|u|/(1+|t|))^{-L}.
\end{equation}
At a real place use the nonzero second derivative on the annulus.
At a complex place the real Hessian determinant of $\operatorname{Re}(tz^k)$
is a nonzero constant times $-|t|^2|z|^{2k-4}$.
After a finite partition, the stationary-phase theorem
\cite[\S7.7]{Hormander2003} applies uniformly in the linear coefficient;
outside the gradient range repeated integration by parts gives the second factor.

For $\mathfrak a=\mathfrak a_\gamma$, $Q=N\mathfrak a$, ideal Poisson gives
\[
 f(\gamma+\beta)=\frac XQ
  \sum_{\xi\in\mathfrak a^{-1}\mathfrak D^{-1}}
    C_{\mathfrak a}(\gamma,\xi)
    \prod_v I_v(P^k\beta_v,P\xi_v),
 \quad
 C_{\mathfrak a}(\gamma,\xi)=
    \sum_{r\bmod\mathfrak a}e_K(\gamma r^k+\xi r).
\]
Primitivity and \eqref{eq:ordinary-complete} apply uniformly in every $\xi$.
Equation~\eqref{eq:ordinary-packing}, with radii
$(1+P^k|\beta_v|)/P$, bounds the sum of the decaying Fourier factors by
$O_K(1+QA_P(\beta)/X)$.
Hence, on the broad arcs,
\begin{equation}\label{eq:ordinary-broad-point}
 |f(\gamma+\beta)|\ll
       XQ^{-1/k+\epsilon}A_P(\beta)^{-1/2}.
\end{equation}
On $QA_P(\beta)\le X^c$ the same argument, using the norm lower bound for
each nonzero dual frequency, gives for every fixed $B>0$
\begin{equation}\label{eq:ordinary-core-formula}
 f(\gamma+\beta)=XG(\gamma)V(P^k\beta)+O(X^{-B}),
 \quad
 G(\gamma)=Q^{-1}C_{\mathfrak a}(\gamma,0),\quad
 V(t)=\prod_v I_v(t_v,0).
\end{equation}
Indeed the first nonzero normalized maximum frequency is at least
$c_K(X/(QA_P(\beta)))^{1/n}$; enough integrations by parts make its entire
tail an arbitrary power of $QA_P(\beta)/X$.

There are at most $Q$ center classes of conductor $\mathfrak a$.
Changing variables $t=P^k\beta$ in \eqref{eq:ordinary-broad-point} therefore
reduces the broad $p$th moment to
\[
 X^{p-k}\sum_{\mathfrak a}Q^{1-p/k+p\epsilon}
         \int_V A_1(t)^{-p/2}\,dt.
\]
Both factors converge for $p>2k$ and sufficiently small $\epsilon$, proving
\eqref{eq:ordinary-major-bundle}.
For the part outside $\mathfrak N_c$, every broad-arc parametrization has
$QA_1(t)>X^c$.  Put $a=p/k-2>0$, $\epsilon=a/(4p)$, and $\tau=a/2$.
Use
$\mathbf1_{\{QA_1(t)>X^c\}}\le X^{-c\tau}Q^\tau A_1(t)^\tau$.
The ideal exponent becomes $-1-a/4<-1$, and
$p/2-\tau=1+p(k-1)/(2k)>1$ makes every real and complex integral converge.
This proves \eqref{eq:ordinary-pruning-bundle}.
No disjointness of the broad arcs is needed for these upper bounds.

\paragraph{Imported results and adaptations.}
The only deep global moment input in this section is Wooley's theorem already
formulated over number fields.  The other cited ready-made inputs are finite-field
Weil, local-field van der Corput, and archimedean stationary phase.
The conductor geometry, shifted-kernel localization, and product-arc pruning
are proved above.  In particular, the rational minor-arc estimate is not used
as an unproved transfer principle.


% Working mathematical proof for independent audit.  All powers are volume powers.
% Written by the k=5/6 analytic agent, 2026-11-25.
\section{The sixth-power smooth moment: a field transfer requiring only quadratic counts}
\label{sec:g6-transfer}

The rational templates are Vaughan--Wooley, \emph{Further improvements in
Waring's problem}, Acta Math.~174 (1995), Lemmas~3.2--3.7, and
\emph{Further improvements in Waring's problem, IV}, Acta Arith.~94
(2000), 203--285, Lemmas~2.1--2.2 and~3.1 and Theorems~3.4--3.5.
The proof below establishes the particular field analogues required here;
the numerical certificate is then unchanged.

Fix a number field $K$ of degree $n$.  In this section $X=P^n$ is the
volume of an ordinary root box, and $\mathbb T_K=(K\otimes\mathbb R)/\mathfrak D^{-1}$
has probability Haar measure.  We write $\mathrm e_K(t)=e(\operatorname{Tr}t)$.
Constants may depend on $K$, on the fixed number of differencing steps, and on
fixed archimedean boxes.  A box of volume parameter $Y\ge1$ means a fixed
multiple of $\{|x_v|\le Y^{1/n}:v\mid\infty\}$; its lattice cardinality is $O(Y)$.
The smooth set $\mathcal A(Y,R)$ consists of its nonzero algebraic integers
whose prime ideal divisors have norm at most $R$.  The zero term can be added
without changing any exponent below.

An exponent $\lambda_s$ is permissible if, for every $\epsilon>0$, some
$\eta>0$ gives
\[
 \int_{\mathbb T_K}\left|\sum_{x\in\mathcal A(Y,R)}\mathrm e_K(\alpha x^6)\right|^{2s}
 d\alpha\ll Y^{\lambda_s+\epsilon}
 \qquad(1\le R\le Y^\eta).
\]
We always take $s\le\lambda_s\le2s$.  Only finitely many uses of a permissible
estimate occur in any proof below, so $R=X^\eta$ can be made sufficiently small
for every quotient scale that occurs.

\subsection{Two divisor facts}

For fixed $C$ and every $\epsilon>0$, the following assertions hold uniformly
when all displayed nonzero algebraic integers have archimedean size at most $X^C$.

\begin{enumerate}
\item The number of divisors $d\in\mathcal O_K$ of a fixed nonzero $a$, with
all $|d_v|\le X^C$, is $O(X^\epsilon)$.  More generally the number of ideals
of norm at most $X^C$ supported on the prime ideals dividing $(a)$ is
$O(X^\epsilon)$.
\item If $abc\ne0$, the number of solutions of
\begin{equation}\label{eq:g6-binary}
 ax^2+by^2=c,
 \qquad x,y\in\mathcal O_K,\quad |x_v|,|y_v|\le X^C,
\end{equation}
is $O(X^\epsilon)$, uniformly in $a,b,c$.
\end{enumerate}

For the first assertion, count ideal divisors and then generators.  Ratios of
generators are units of logarithmic height $O(\log X)$; Dirichlet's unit lattice
therefore contributes only a power of $\log X$.  For the stronger supported-ideal
assertion, split the supporting primes at a sufficiently large fixed norm $T$.
The finitely many primes below $T$ contribute a power of $\log X$.
For $\sigma>0$ the large-prime Euler product is bounded by
\[
 \prod_{\substack{\mathfrak p\mid(a)\\N\mathfrak p>T}}
 (1-N\mathfrak p^{-\sigma})^{-1}
 \le \exp\!\left(O\!\left(\frac{\log X}{T^\sigma\log T}\right)\right).
\]
Rankin's inequality, first choosing $\sigma$ small and then $T$ large, gives
the assertion.

To prove the second assertion, put $L=K(\sqrt{-ab})$ and
$z=ax+y\sqrt{-ab}$, so that $N_{L/K}z=ac$.
When $[L:K]=2$, the ideal $(z)$ divides $(ac)\mathcal O_L$.
The number of such ideals is $O(X^\epsilon)$ uniformly in $L$:
$[L:\mathbb Q]\le2n$, and a rational prime has at most $2n$ prime ideals above it,
so the ordinary divisor bound applied to the rational norm suffices.
For each ideal, admissible generators differ by units of logarithmic height
$O(\log X)$.  This count is uniform in $L$: non-torsion algebraic units of
degree at most $2n$ have a positive uniform lower bound for their logarithmic
height, by Northcott's theorem and Kronecker's theorem.  Packing in the
logarithmic unit space gives $O((\log X)^{2n})$ possibilities; the number of
torsion units is bounded in terms of $n$.  Recovering $x,y$ from $z$ is injective.
If $-ab=d^2$ in $K$, then $d\in\mathcal O_K$ and
$(ax-dy)(ax+dy)=ac$, so the first assertion applies directly.
The hypothesis $c\ne0$ is essential; zero right sides are handled separately below.

\subsection{Localizing the ordinary variables}

One may prove a permissible full-box moment by proving it for smooth roots in
arbitrarily small fixed boxes away from zero at every embedding.  Here are the
details that allow this reduction uniformly in the scale.
Partition nonzero roots into dyadic ranges for each $|x_v|$.  There are
$O((\log X)^{r_1+r_2})$ ranges, since $|Nx|\ge1$ supplies a polynomial lower
bound for every $|x_v|$.  Roots with $|Nx|\le X^{1/2}$ number
$O(X^{1/2+\epsilon})$, by counting ideals and bounded-height units; their
$2s$th moment is trivially $O(X^{s+\epsilon})$.
For each remaining annular box let $Y$ be the product of its archimedean radii,
with local dimensions counted.  Thus $X^{1/2}\ll Y\ll X$.
Multiplication by one unit balances these radii to fixed multiples of $Y^{1/n}$.
It preserves ideal smoothness and the mean value, since multiplication of
$\alpha$ by a unit's sixth power is a torus automorphism.
Cover the resulting compact annulus by finitely many boxes of relative diameter
$\delta$, where $\delta>0$ is a sufficiently small constant depending only on
$K$ and the maximal differencing depth.  H\"older's inequality for the sum of
these pieces costs only a power of $\log X$.  Since $R=X^\eta\le Y^{2\eta}$,
a sufficiently small initial smoothness exponent works for all pieces.

For one such piece, replace two smooth roots by ordinary roots in that same
small box.  The other smooth roots can subsequently range over full boxes.
The two distinguished ordinary roots, and every pair produced from them by
efficient differencing, retain their actual small support.
In centered notation define
\[
 \Psi_j(z;\mathbf h;\mathbf m)=
 \Delta^*_j\bigl((z-\textstyle\sum_{i\le j}h_i m_i^6)^6;
                 2\mathbf h;\mathbf m\bigr),
 \qquad
 \Delta^*_1 f(z;h;m)=m^{-6}(f(z+hm^6)-f(z)).
\]
These are integral polynomials.  Their leading term is
$2^j(6)_j(h_1\cdots h_j)z^{6-j}$.
At each embedding, the actual support has $z$ within $O_j(\delta P)$ of a
fixed nonzero center of size comparable to $P$, and
\begin{equation}\label{eq:g6-small-increments}
 |h_{i,v}m_{i,v}^6|\le C_j\delta P\qquad(1\le i\le j).
\end{equation}
Fixed factors of two change the center and the $C_j$, harmlessly.
The central change of variables is $z=z_1+z_2$,
$m^6h=z_1-z_2$; divisibility by $2$ can be retained as a support restriction.
More explicitly, for exponent $k$ the centered polynomial identity multiplies
the differenced equation by $2^k$, doubles the earlier shift parameters, and
replaces every remaining quotient root $x$ by $2x$.
This statement applies at every repeated step, and also with $k=7$ below.
The retained parity restrictions only reduce a solution count; subsequently
enlarging the bounding boxes changes their scales by fixed constants, and the
primes above $2$ belong to the smoothness range for sufficiently large $X$.
Consequently $\delta$ can be chosen so that, whenever all $h_i\ne0$,
\begin{equation}\label{eq:g6-nonzero}
 \Psi_j(z;\mathbf h;\mathbf m)\ne0,
 \qquad \partial_z\Psi_j(z;\mathbf h;\mathbf m)\ne0
 \quad(0\le j\le5).
\end{equation}
Indeed, after dividing by $h_1\cdots h_j$, the leading monomial dominates
all remaining monomials at each embedding.  The same is true for the derivative.
An ordinary differencing shift is a difference of two points in the small box;
it therefore retains \eqref{eq:g6-small-increments}.

In all subsequent counts the coupled restrictions
\eqref{eq:g6-small-increments} are retained.  The separate bounding boxes for
$h_i,m_i$ are used only to count parameter choices, never to replace the support
when \eqref{eq:g6-nonzero} is needed.  This is the step that removes the
positivity assumption implicit in the rational auxiliary counts.

\subsection{The fundamental efficient-differencing inequality}

Let $F_\Psi(\alpha)=\sum_{z,\mathbf u}\mathrm e_K(\alpha\Psi(z,\mathbf u))$,
where $z$ ranges in a box of volume $X$, the auxiliary parameter tuples number
at most $U$, every polynomial has bounded positive degree and coefficients of
polynomial height, and $\Psi'(z,\mathbf u)\ne0$ on the retained support.
Let $g_Q$ be the full smooth sum of volume $Q$, and put
$S_s=\int|F_\Psi|^2|g_Q|^{2s}$.
Choose $1\le M\le Q$.

Let $h_K$ be the order of the ideal class group.  Any sequence of $h_K$ ideal
classes has a nonempty zero-sum subsequence.  Partition the prime ideal factors
of a smooth ideal into such principal blocks, each with at most $h_K$ factors,
leaving fewer than $h_K$ factors.  If the part of $(x)$ coprime to
$(\Psi'(z,\mathbf u))$ has norm greater than $MR^{h_K}$, successive principal
blocks supply a divisor $(m)$ of $(x)$, coprime to $\Psi'$, with
\[
 M<N(m)\le MR^{h_K}.
\]
Choose a generator balanced by units.  There are $O(MR^{h_K+\epsilon})$
possible generators, and $x/m$ is a smooth algebraic integer in a fixed box
of volume $Q/M$.  Powers of $R$ will always be absorbed in $X^\epsilon$.

If no such divisor exists, the prime-to-$\Psi'$ part of $(x)$ has norm at most
$MR^{h_K}$.  For fixed $z,\mathbf u$ the number of these bad $x$ is at most
$MX^\epsilon$: count the prime-to-$\Psi'$ ideal and use the supported-ideal
bound for its complementary ideal, then count bounded-height generators.
This bound also covers roots of norm at most $MR^{h_K}$.

The contribution of a bad variable is at most
$S_s^{1/2}(QMX^\epsilon S_{s-1})^{1/2}$.
For clarity, if
$G_x=\sum_{z,\mathbf u:,x\text{ bad}}\mathrm e_K(\alpha\Psi(z,\mathbf u))$,
then
\[
 \left|\sum_x\mathrm e_K(\alpha x^6)G_x\right|^2
 \le O(Q)\sum_x|G_x|^2.
\]
After multiplication by $|g_Q|^{2s-2}$ and integration, a fixed pair of auxiliary
tuples admits at most $MX^\epsilon$ common bad $x$, giving the assertion.

For the good variables, extract a principal divisor separately from each of the
$2s$ smooth roots.  The dependence of this extraction on the derivative is
handled as follows.  Let $\mathcal M$ be the set of balanced extracted divisors,
so that $L=\#\mathcal M\ll MX^\epsilon$.  For $W\ne0$ put
\[
 F_W(\alpha)=\sum_{\substack{z,\mathbf u\text{ on the retained support}\\
              ((W),(\Psi'(z,\mathbf u)))=1}}
       \mathrm e_K(\alpha\Psi(z,\mathbf u)),
 \quad g=g_{Q/M},\quad
 V_m=\int|F_m(\alpha)|^2|g(m^6\alpha)|^{2s}\,d\alpha.
\]
For tuples $\mathbf m,\mathbf n\in\mathcal M^s$ put
$W=\prod_i m_i$ and $W'=\prod_i n_i$.
After enlarging the quotient-root boxes, orthogonality bounds the good count by
the sum over these tuples of
\[
 \int |F_W(\alpha)F_{W'}(\alpha)|
       \prod_{i=1}^s|g(m_i^6\alpha)g(n_i^6\alpha)|\,d\alpha.
\]
H\"older bounds each integral by the product of $2s$ factors of the form
$(\int|F_W|^2|g(m_i^6\alpha)|^{2s})^{1/(2s)}$ and their $W',n_i$ analogues.
The coefficient set of $F_W$ is contained in that of $F_{m_i}$, because
coprimality to $W$ implies coprimality to $m_i$.  Orthogonality, with its
nonnegative solution counts, therefore bounds that integral by $V_{m_i}$.
Consequently the total good count is at most
\[
 \left(\sum_{m\in\mathcal M}V_m^{1/(2s)}\right)^{2s}
 \le L^{2s-1}\sum_{m\in\mathcal M}V_m.
\]
Both derivative coprimalities are present in every $V_m$.

Here is the remaining Cauchy step explicitly.  Set $\mathfrak b=(m)^6$.
For $\rho\in\mathcal O_K/\mathfrak b$ and fixed $\mathbf u$, let
$\mathcal R(\rho,\mathbf u)$ be the nonsingular residue classes $r$ satisfying
$\Psi(r,\mathbf u)\equiv\rho\pmod{\mathfrak b}$.
Their number is at most $D^{\omega((m))}\ll X^\epsilon$, by Hensel lifting
and the Chinese remainder theorem.  Write
$F_{\mathbf u,r}=\sum_{z\equiv r\, (\mathfrak b)}
\mathrm e_K(\alpha\Psi(z,\mathbf u))$, retaining the original support.
Orthogonality kills pairs with different $\Psi$ residues, and hence
\[
 V_m=\sum_\rho\int
 \left|\sum_{\mathbf u}\sum_{r\in\mathcal R(\rho,\mathbf u)}
                   F_{\mathbf u,r}(\alpha)\right|^2
 |g(m^6\alpha)|^{2s}\,d\alpha.
\]
Cauchy's inequality inside each square costs at most $UX^\epsilon$ and leaves
the sum of $|F_{\mathbf u,r}|^2$.  Its underlying equations have the same
auxiliary tuple on both sides and $z\equiv w\pmod{(m)^6}$.
Write $z-w=m^6h$.  The case $h=0$ contributes at most
$XUM(Q/M)^{\lambda_s+\epsilon}$; the other cases are bounded by
$T_{\rm next}=\int|F_{\rm next}||g_{Q/M}|^{2s}$, with all actual support
restrictions retained.  Thus the good contribution is at most
\[
 X^\epsilon U M^{2s-1}
 \{XUM(Q/M)^{\lambda_s}+T_{\rm next}\}.
\]

The bad contribution can be absorbed into this same explicit diagonal bound
without a lower bound for the density of smooth algebraic integers.
Indeed, $S_0\ll XU^2$, by fixing $z,\mathbf u,\mathbf v$ and solving the
nonconstant equation for $w$.  H\"older gives
$S_{s-1}\le S_0^{1/s}S_s^{1-1/s}$.
If the bad contribution dominates, then
\[
 S_s\ll X^\epsilon XU^2(QM)^s
 \le X^\epsilon XU^2M^{2s}(Q/M)^{\lambda_s},
\]
since $Q\ge M$ and $\lambda_s\ge s$.
In every case, therefore,
\begin{equation}\label{eq:g6-fundamental}
 S_s\ll X^\epsilon U M^{2s-1}
 \{XUM(Q/M)^{\lambda_s}+T_{\rm next}\}.
\end{equation}
This is the required field version of the fundamental efficient-differencing
inequality.  Its proof is the congruence argument of Wooley's 1990 thesis,
Chapter~6, with principal ideal blocks and the preceding explicit treatment
of bad variables and the diagonal term.

\subsection{The auxiliary second and fourth moments actually needed}

Put $M_i=X^{\phi_i}$, $H_i=XM_i^{-6}$,
$\Phi_j=\sum_{i\le j}\phi_i$,
$\widetilde M_j=\prod_{i\le j}M_i$,
$\widetilde H_j=\prod_{i\le j}H_i$, and $U_j=\widetilde H_j\widetilde M_j$.
Let $F_j$ be the sum of $\mathrm e_K(\alpha\Psi_j)$ over the retained small
support, $h_i\ne0$, with $h_i,m_i$ in boxes of volume $O(H_i)$ and
$O(M_iR^{h_K})$ respectively.  All $R$ losses below are absorbed into $X^\epsilon$.
Assume
\begin{equation}\label{eq:g6-parameters}
 0\le\phi_j\le\cdots\le\phi_1\le1/6,
 \qquad \Phi_I+6(\phi_{I-1}+\phi_I)\le2\quad(3\le I\le j).
\end{equation}
Then the required estimates are
\begin{align}
 \int|F_j|^2&\ll X^{1+\epsilon}U_j &&(1\le j\le4),\label{eq:g6-aux2}\\
 \int|F_j|^4&\ll X^{2+\epsilon}U_j^3 &&(1\le j\le3),\label{eq:g6-aux4}\\
 \int|F_4|^4&\ll X^{2+\epsilon}U_4^3\widetilde M_4.\label{eq:g6-aux4special}
\end{align}

Here is a proof retaining the exact low-degree algebra.  Since every phase value
is nonzero and divisible by $h_1\cdots h_j$, the divisor bound and Cauchy's
inequality reduce the second moment, at cost $X^\epsilon$, to solutions of
\[
 \Psi_j(z;\mathbf h;\mathbf m)=\Psi_j(w;\mathbf h;\mathbf n).
\]
For odd residual degree $6-j$, the phase is also divisible by $z$.
The same argument forces $z=w$.  Unless the two residual polynomials coincide,
$z$ has only $O(1)$ possibilities; this costs at most
$\widetilde H_j\widetilde M_j^2\ll XU_j$.
Coincidence implies the coefficient equations
\begin{equation}\label{eq:g6-K}
 \sum_{i\le j}h_i^{2r}(m_i^{12r}-n_i^{12r})=0
 \quad(1\le r\le\lfloor(6-j)/2\rfloor).
\end{equation}
For residual degree two the equation is
$3(z^2-w^2)+\sum h_i^2(m_i^{12}-n_i^{12})=0$.
If the second summand is nonzero, factor $z^2-w^2$ and use the divisor bound;
otherwise $z=w$, since the small support excludes $z+w=0$.
For residual degree four put
$\Xi_r=\sum h_i^{2r}m_i^{12r}$ and
$\Gamma_r=\sum h_i^{2r}n_i^{12r}$.
The equation becomes
\[
 15(u^2-v^2)=10(\Xi_1^2-\Gamma_1^2)+2(\Xi_2-\Gamma_2),
 \qquad u=z^2+\Xi_1,\quad v=w^2+\Gamma_1.
\]
The small-increment support excludes $u+v=0$ as well as $z+w=0$.
Factoring first $u^2-v^2$, then $z^2-w^2$ when necessary, gives
$O(X^\epsilon)$ choices unless $z=w$, $\Xi_1=\Gamma_1$, and
$\Xi_2=\Gamma_2$.  Consequently, if $K_j$ counts \eqref{eq:g6-K},
\begin{equation}\label{eq:g6-reduction}
 \int|F_j|^2\ll X^{1+\epsilon}K_j+X^{1+\epsilon}U_j.
\end{equation}

For $j=1$, $K_1\ll H_1M_1$ by the bounded number of twelfth roots.
For $j=2$, it suffices to retain the first equation.  Put
$A_i(\alpha)=\sum_{h_i}|\sum_{m_i}\mathrm e_K(\alpha h_i^2m_i^{12})|^2$.
H\"older and the fourth Hua moment for degree twelve give
$\int A_i^2\ll H_i^2M_i^{2+\epsilon}$; multiplication by nonzero $h_i^2$
preserves Haar measure.  Hence $K_2\ll X^\epsilon U_2$.

For $j=3,4$ retain only
$\sum h_i^2a_i=0$, where $a_i=m_i^{12}-n_i^{12}$.
Indices with $a_i=0$ contribute $O(H_iM_i)$ and can be removed.
For two remaining indices, fixing one triple $(h_i,m_i,n_i)$ determines the
other in $O(X^\epsilon)$ ways, by factoring
$h^2(m^6-n^6)(m^6+n^6)$.
The symmetric geometric mean bounds this contribution by
$X^\epsilon H_iH_jM_iM_j$.
For three remaining indices the same argument gives
$X^\epsilon(\prod H_i)^{2/3}(\prod M_i)^{4/3}$.
This is at most $X^\epsilon\prod H_iM_i$ under
$M_1M_2M_3\le H_2H_3$; the ordering in \eqref{eq:g6-parameters} makes this
condition sufficient for every three-index subset.
For four nonzero $a_i$, split according as
$c=h_1^2a_1+h_2^2a_2$ is zero.  If $c=0$, apply the two-index bound twice.
Otherwise fix all $m_i,n_i,h_1,h_2$ and apply \eqref{eq:g6-binary} to
$a_3h_3^2+a_4h_4^2=-c$.
The result is
$X^\epsilon H_1H_2\widetilde M_4^2\le X^\epsilon U_4$,
because $\widetilde M_4\le H_3H_4$.
These are precisely the two parameter inequalities in
\eqref{eq:g6-parameters}.  Thus \eqref{eq:g6-aux2} follows.

One ordinary differencing step in $z$, followed by Cauchy's inequality, gives
\begin{equation}\label{eq:g6-fourth-reduction}
 \int|F_j|^4\ll XU_j^2\int|F_j|^2
 +U_j^2\int|F_{j+1}^{(0)}|^2,
\end{equation}
where the new divisor is $m_{j+1}=1$ and the new shift has volume $X$.
Its zero shift is the first term; nonzero shifts retain the small support.
Explicitly, writing $F_j=\sum_{\mathbf u}f_{\mathbf u}$ and expanding
$\sum_{\mathbf u}|f_{\mathbf u}|^2$ gives
$|F_j|^2\ll XU_j^2+U_j|G|$, where $G$ is the sum over nonzero ordinary
shifts.  Thus, with $B=\int|F_j|^4$,
$B\ll XU_j^2\int|F_j|^2+U_jB^{1/2}(\int|G|^2)^{1/2}$.
Absorption and the centered change of variables give
\eqref{eq:g6-fourth-reduction}; orthogonality bounds
$\int|G|^2$ by $\int|F_{j+1}^{(0)}|^2$.
For $j\le3$, the preceding coefficient counts give
$\int|F_{j+1}^{(0)}|^2\ll X^{2+\epsilon}U_j$.
Although appending a zero $\phi$ removes a coefficient equation in some cases,
the proofs above already use only the first equation when that matters.
This proves \eqref{eq:g6-aux4}.
For $j=4$, the fifth difference is a fixed nonzero constant times
$zh_1\cdots h_5$, independent of the $m_i$.
The divisor bound therefore gives
$\int|F_5^{(0)}|^2\ll X^{2+\epsilon}\widetilde H_4\widetilde M_4^2$,
which proves \eqref{eq:g6-aux4special}.

\subsection{Exact recurrence and the twentieth moment}

Put $q_j=1-\Phi_j$, and let $T_j(s)$ denote an upper exponent for
$\int|F_j||g_{X^{q_j}}|^{2s}$.
For $0<t<2s$ and $p=2s-t$, Cauchy's inequality followed by
\eqref{eq:g6-fundamental} gives
\begin{align*}
 T_j(s)\le\tfrac12\bigl[&\lambda_tq_j+j-5\Phi_j+(2p-1)\phi_{j+1}\\
 &+\max\{1+j-6\Phi_j+\Phi_{j+1}+\lambda_pq_{j+1},T_{j+1}(p)\}\bigr].
\end{align*}
The initial joint moment of order $2S$ has upper exponent
\[
 \max\{1+2(S-1)\phi_1+\lambda_{S-1}q_1,
             (2S-3)\phi_1+T_1(S-1)\}.
\]
The proven terminal second and fourth moments have exponents
\[
 A_{d,\ell}=2^\ell-\ell+(2^\ell-1)d-5(2^\ell-1)\Phi_d
                +\mathbf1_{d+\ell=6}\Phi_d,
 \qquad\ell\in\{1,2\},
\]
for the depths used in the eight-row certificate.
H\"older supplies the terminal mixed exponent with the weights specified there.
Thus every arithmetic update in the certificate is a proved field inequality.
The classical Hua moments over $K$ supply the initial orders
$(r,\mu_r)=(0,0),(2,1),(4,2),(8,5),(16,12),(32,27)$.
The exact finite contraction certificate therefore proves
\begin{equation}\label{eq:g6-twentieth}
 \int_{\mathbb T_K}|g(\alpha)|^{20}\,d\alpha
 \ll X^{14.22+\epsilon}.
\end{equation}
For each fixed $\epsilon$ only finitely many iterations are required; choose the
initial smoothness exponent after allocating the total error budget among them.
Together with the product-root endgame and primitive local solubility, this is
the analytic input for $G_{\mathbb K}(6)\le24$.


\section{An eight-row exact certificate for the sixth-power moment}
\label{sec:g6-certificate}

Write $\mu_r=\lambda_{r/2}$ for an exponent of the full $r$th moment.
We fix $\mu_4=2$. All decimals and parameter numerators in this section
are exact rational numbers. The following eight rows suffice:
\begin{center}
\small
\begin{tabular}{@{}r l c c l r@{}}
\toprule
$S$ & Cauchy path & $(\ell_1,\ell_2)$ & $u$ & $10^5(\phi_1,\ldots,\phi_d)$ & $b_{2S}$\\
\midrule
3 & empty & $(1,1)$ & $8$ & $(4545)$ & 3.0915\\
4 & empty & $(2,2)$ & $8$ & $(8333)$ & 4.334\\
5 & $(5)$ & $(1,2)$ & $10$ & $(10673,5080)$ & 5.726\\
6 & $(6)$ & $(2,2)$ & $32/3$ & $(11855,8751)$ & 7.233\\
7 & $(7,6)$ & $(1,2)$ & $12$ & $(12981,10763,5551)$ & 8.852\\
8 & $(8,7)$ & $(2,2)$ & $40/3$ & $(13784,12076,8562)$ & 10.562\\
9 & $(9,8,8)$ & $(1,1)$ & $16$ & $(14593,13194,10858,1059)$ & 12.356\\
10 & $(10,9,9)$ & $(1,2)$ & $18$ & $(15303,14283,12101,6908)$ & 14.220\\
\bottomrule
\end{tabular}
\end{center}

Here is a complete specification of the arithmetic. A path $(t_1,\ldots,t_{d-1})$
defines $s_1=S-1$ and $s_{j+1}=2s_j-t_j$. Put
$\Phi_j=\sum_{i\le j}\phi_i$ and $q_j=1-\Phi_j$.
At the terminal depth put $c=2s_d/u$ and choose nonnegative $a,b$ with
\[
 a+b+c=1,\qquad 2^{\ell_1}a+2^{\ell_2}b=1.
\]
If $\ell_1=\ell_2$, take $a=2^{-\ell_1}$ and $b=0$.
Define
\[
 A_{d,\ell}=2^\ell-\ell+(2^\ell-1)d
       -5(2^\ell-1)\Phi_d+\mathbf1_{d+\ell=6}\Phi_d,
\]
and initialize
\[
 T_d=aA_{d,\ell_1}+bA_{d,\ell_2}+cq_d\mu_u.
\]
For $j=d-1,\ldots,1$, put $p=2s_j-t_j$ and use
\[
 T_j=\frac12\left\{q_j\mu_{2t_j}+j-5\Phi_j+(2p-1)\phi_{j+1}
 +\max\left(1+j-6\Phi_j+\Phi_{j+1}+q_{j+1}\mu_{2p},T_{j+1}\right)\right\}.
\]
Finally the row output is
\[
 \mathcal T_{2S}(\boldsymbol\mu)=
 \max\left\{1+2(S-1)\phi_1+q_1\mu_{2S-2},
                         (2S-3)\phi_1+T_1\right\}.
\]
Only two nonintegral moments occur, and their meaning is fixed as
\[
 \mu_{32/3}=(2\mu_{10}+\mu_{12})/3,
 \qquad \mu_{40/3}=(\mu_{12}+2\mu_{14})/3.
\]
These are ordinary H\"older interpolations. The terminal estimates and every
backwards step are exactly those proved in Section~\ref{sec:g6-transfer}.

Substitution of the table's vector $\boldsymbol b$ gives:
\begin{center}
\begin{tabular}{@{}r r r r@{}}
\toprule
$r$ & $b_r$ & $\mathcal T_r(\boldsymbol b)$ & $b_r-\mathcal T_r(\boldsymbol b)$\\
\midrule
6 &3.0915&3.091234850&0.000265150\\
8 &4.334&4.333865305&0.000134695\\
10&5.726&5.725616976&0.000383024\\
12&7.233&7.232682700&0.000317300\\
14&8.852&8.851804270&0.000195730\\
16&10.562&10.561620060&0.000379940\\
18&12.356&12.355567340&0.000432660\\
20&14.220&14.219701320&0.000298680\\
\bottomrule
\end{tabular}
\end{center}
The displayed outputs are exact terminating decimals. All parameter orderings
hold. The largest auxiliary side expression
$\Phi_I+6(\phi_{I-1}+\phi_I)$ is $1.99991<2$.

The map is monotone and has sup-norm contraction constant at most $23/25$.
To check this without a numerical eigensolver, the total coefficient of the
variable moment exponents at the terminal step is at most $q_d$; a backwards
step replaces this upper bound by at most $(q_j+q_{j+1})/2\le q_j$.
For $S\ge4$, both initial branches have coefficient sum at most
$q_1\le0.91667<23/25$. For $S=3$ the initial diagonal uses only the fixed
$\mu_4=2$, and the other branch has coefficient sum $q_1/2<1/2$.
Maxima preserve this common bound. If necessary, replace a row output by its
maximum with $r/2$, to maintain the convention $\mu_r\ge r/2$ used in the
fundamental lemma; the table and the contraction are unchanged.

Hua's anchors $(r,\mu_r)=(2,1),(4,2),(8,5),(16,12),(32,27)$ and
interpolation give the initial vector
\[
 (3.5,5,6.75,8.5,10.25,12,13.875,15.75)
\]
at orders $6,8,\ldots,20$. Its largest excess over $\boldsymbol b$ is $1.53$.
Consequently finite iteration gives
$\boldsymbol\mu^{(j)}\le\boldsymbol b+1.53(23/25)^j\mathbf1$.
For each prescribed $\epsilon>0$, choose a finite $j$ making this excess
smaller than $\epsilon/2$, and allocate the remaining $\epsilon/2$ among
the finitely many analytic operations. This proves that every $b_r$ is
permissible, in particular
\[
 \int|g_6|^{20}\ll X^{14.22+\epsilon}.
\]
Alternatively, already after 79 rounds the numerical exponent is at most
$14.222108287293<14+2/9$, leaving a strictly positive endgame margin.

Using the permissible excess $\Delta_{20}=0.22$, the twenty-four-variable
criterion in Section~\ref{sec:product-endgame} has
\[
 D=4-18\Delta_{20}=1/25,\qquad D/42=1/1050.
\]
Four ordinary roots and twenty product roots therefore give the required
broad-minor power saving. The primitive local bound $h_6\le12$ is more than
sufficient at twenty-four roots, completing the sixth-power assertion of
Theorem~\ref{thm:main-transfer}.


% Supplement to the general efficient-differencing and quadratic-count section.
\section{The auxiliary moments for seventh powers}
\label{sec:g7-auxiliary}

Use the same divisor, binary-quadratic, small-support, and fundamental
congruence lemmas as for sixth powers, replacing the exponent $6$ by $7$.
The proof of that fundamental inequality uses only a fixed polynomial degree,
Hensel lifting at nonsingular roots, and principal ideal blocks; its statement
and proof therefore apply verbatim with $m^7$ in the congruence.
Put
\[
 M_i=X^{\phi_i},\quad H_i=XM_i^{-7},\quad
 \Phi_j=\sum_{i\le j}\phi_i,\quad
 U_j=\prod_{i\le j}H_iM_i.
\]
For $j\le4$ require
\begin{equation}\label{eq:g7-aux-conditions}
 0\le\phi_j\le\cdots\le\phi_1\le1/7,
 \qquad
 \Phi_I+7(\phi_{I-1}+\phi_I)\le2\quad(3\le I\le j).
\end{equation}
Then the estimates needed below are
\begin{equation}\label{eq:g7-auxmoments}
 \int_{\mathbb T_K}|F_j(\alpha)|^2\,d\alpha
       \ll X^{1+\epsilon}U_j,\qquad
 \int_{\mathbb T_K}|F_j(\alpha)|^4\,d\alpha
       \ll X^{2+\epsilon}U_j^3
 \qquad(1\le j\le4).
\end{equation}

For $j=2,3,4$ the residual degrees are $5,4,3$.  Their second moments are
exactly the odd-degree or quadratic-factor arguments already proved in the
sixth-power case.  The resulting coefficient equations are
\[
 \sum_{i\le j}h_i^{2r}(m_i^{14r}-n_i^{14r})=0.
\]
Their counts use only the first equation for $j\ge2$.
The pair conditions \eqref{eq:g7-aux-conditions} are precisely
$\prod_{i\le I}M_i\le H_{I-1}H_I$ for $I=3,4$.
Thus the same binary-quadratic argument bounds the coefficient count by
$X^\epsilon U_j$.  Appending one ordinary difference, with new divisor $1$
and shift of volume $X$, gives the fourth moment by Cauchy.  At $j=4$ the
appended residual degree is $2$; its last coefficient cancels in the
coincidence equations because the new divisor is $1$.  The earlier four
coefficient variables are still controlled by the same pair conditions.
Consequently this appended second moment is $O(X^{2+\epsilon}U_4)$,
which gives the fourth-moment assertion without the extra $M$ loss that
occurs for sixth powers at depth four.

Only the second moment at $j=1$, of residual degree six, requires an additional
input.  We give the argument explicitly.

\subsection{A ready-made affine curve theorem over $K$}

Binyamini, Cluckers and Kato \cite[Corollary~12]{BCK2025}, \emph{Proceedings of the London Mathematical
Society} 130 (2025), article e70016, Corollary~12, prove that for an
irreducible affine curve $C/K$ of degree $d>1$ in fixed affine space,
\begin{equation}\label{eq:g7-BCK}
 \#\{x\in C(\mathcal O_K): |x_v|\le Y^{1/n}\text{ for every }v\}
 \ll_{K,m} d^2Y^{1/d}(\log Y)^\kappa.
\end{equation}
This is their volume-height normalization, the definition preceding~(28); the constant
is independent of the coefficients of the curve.  All number fields are
allowed, including fields with no real embedding.  Only the fixed case
$d=3$, $m=2$ is used here.

\subsection{Residual degree six}

The nonzero phase-value divisor argument first forces the same $h$ on both
sides.  The centered polynomial has the form
\[
 \Psi_1(z;h;m)=h\sum_{r=0}^{3}c_r(hm^7)z^{2r},
\]
where $c_3$ is a fixed nonzero integer, and $c_r$ are integral polynomials.
Write $R_0$ for the solutions with $c_0(hm^7)=c_0(hn^7)$ and $R_1$ for
the remainder.  Here the explicit identity is
\[
 \Psi_1(z;h;m)=2h\{7z^6+35h^2m^{14}z^4
                       +21h^4m^{28}z^2+h^6m^{42}\}.
\]
Thus the constant-term equality defining $R_0$ gives $m^{42}=n^{42}$.
For each fixed $h,m$ there are at most $42$ choices for $n$, and for each
$z$ the equation for $w$ has at most six roots.  Consequently
\[
 R_0\ll XH_1M_1.
\]

Fix $h,m,n$ contributing to $R_1$ and set
\[
 F(x,y)=\sum_{r=0}^{3}\{c_r(hm^7)x^r-c_r(hn^7)y^r\}.
\]
Its constant term is nonzero and its highest homogeneous part is
$c_3(x^3-y^3)$.  We must count $F(z^2,w^2)=0$ with all conjugates of
$z,w$ of size $O(X^{1/n})$.
An irreducible cubic factor is handled by \eqref{eq:g7-BCK} in a box of
volume $O(X^2)$, yielding $O(X^{2/3+\epsilon})$ pairs $(z^2,w^2)$.
Each such pair has at most four preimages $(z,w)$.

For clarity, factors of degree at most two cause no extra loss over general
$K$.  A linear factor has nonzero coefficients of both $x$ and $y$, since
its leading term divides $x^3-y^3$, and nonzero constant term.
After substitution it is a nonzero-right-side binary quadratic equation
$az^2+bw^2=c$, so the uniform binary-quadratic divisor lemma gives
$O(X^\epsilon)$ solutions.
A quadratic factor has nondegenerate highest homogeneous part: its two
linear factors over an algebraic closure have distinct slopes among the
three cube roots of unity.  Complete the square by an affine change of
variables over $K$.  If the resulting constant is zero, the conic splits
over an algebraic closure; if it does not split over $K$, its $K$-points lie
in the intersection of conjugate lines and there are $O(1)$ of them.
If it splits over $K$, treat its linear factors as above.
If the completed constant is nonzero, diagonalize the nondegenerate binary
quadratic form and apply the same binary-quadratic divisor lemma.
The standard fixed-degree coefficient bound for polynomial factors gives
factor coefficients of polynomial height, uniformly in $h,m,n$.
All changes of variables can be cleared of denominators of polynomial
height, so its uniformity in the coefficients applies.
Thus quadratic factors also contribute $O(X^\epsilon)$.
Factors defined only over an extension of $K$ can alternatively be disposed
of by intersecting their distinct Galois conjugates, giving $O(1)$ points.

There are $O(H_1M_1^2X^\epsilon)$ choices for $h,m,n$, so
\[
 R_1\ll X^{2/3+\epsilon}H_1M_1^2
       \ll X^{1+\epsilon}H_1M_1,
\]
because $M_1\le X^{1/7}<X^{1/3}$.
This proves the remaining second moment in \eqref{eq:g7-auxmoments}.
One ordinary difference now has residual degree five and supplies its fourth
moment using the already proved odd-degree argument.

\paragraph{Status of the input.}
The affine curve estimate \eqref{eq:g7-BCK} is imported as a theorem already
stated over number fields.  Its application to the efficient-differencing
polynomials, including the low-degree factors and the narrow complex supports,
is the adaptation proved here.  No numerical smooth-moment exponent from a
rational-only theorem is imported.


% Auxiliary proof for the septic transfer. Every size variable denotes volume.
\section{A cubic auxiliary dissection over an arbitrary number field}
\label{sec:g7-cubic}

All parameters in this section are volumes.  Fix $K$, write $n=[K:\mathbb Q]$,
and use probability Haar measure on $\mathbb T_K=V_K/\mathfrak D^{-1}$.
For $T\ge1$ put
\[
 A_T(\beta)=\prod_{v\mid\infty}\max(1,T^{1/n}|\beta_v|)^{d_v},\qquad
 \mathcal H_T(\alpha)=\min_{\gamma\in K}Q_\gamma A_T(\alpha-\gamma).
\]
The conductor norm $Q_\gamma$ and the affine field-line lemma are as in the
quartic argument.  Multiplication by a nonzero algebraic integer is a
surjective endomorphism of $\mathbb T_K$ and preserves its probability measure.

\subsection{A linear-kernel consequence of the bilinear large sieve}

The elementary bilinear sampling bound is
\begin{equation}\label{eq:g7-bilinear}
 \sum_{z\in\mathcal B(J)}\left|\sum_{a\in\mathcal B(N)}c(a)e_K(\theta za)\right|^2
 \ll_K\left(J+N+\frac{JN}{\mathcal H_{JN}(\theta)}\right)\sum_a|c(a)|^2.
\end{equation}
For $J\le N$ this follows by the Gaussian Gram-matrix proof of the elementary
large sieve and the affine field-line lemma: the maximum number of sample
frequencies in a box of side $N^{-1/n}$ is
$O_K(1+J/\mathcal H_{JN}(\theta))$.  Partition the $z$ box into a fixed number
of smaller boxes before applying the line lemma.  For $J>N$, transpose the
finite matrix $(e_K(\theta za))_{z,a}$ and use duality.  This gives the same
bound with $J,N$ interchanged.  The argument works for fixed translates,
subsets, and fixed multiples of the boxes.  For each fixed nonzero
$c\in\mathcal O_K$, the heights $\mathcal H_T(c\alpha)$ and
$\mathcal H_T(\alpha)$ are comparable up to constants depending on $K,c$:
multiply or divide an approximating center by $c$ and use the conductor
index bound and the fixed archimedean multiplication factors.

Here is the corresponding positive kernel estimate.  Let
\[
 D_N(t)=\sum_{z\in\mathcal O_K}
      \exp(-c|z|^2/N^{2/n})e_K(tz).
\]
Choose the Euclidean norm compatibly with the trace pairing.  Poisson
summation writes this as a positive sum of Gaussians centered on
$\mathfrak D^{-1}$, each of height $O_K(N)$ and width $N^{-1/n}$.
Then
\begin{equation}\label{eq:g7-kernel}
 \sum_{h\in\mathcal B(J)}D_N(\theta h)
 \ll_K J+N+\frac{JN}{\mathcal H_{JN}(\theta)}.
\end{equation}
Indeed, enlarge the $h$ box by a fixed factor.  Every difference in the
original box has $\gg_K J$ representations as $r-s$ in the enlarged box.
Because $D_N\ge0$, its sum over the original box is at most a constant times
$J^{-1}\sum_{r,s}D_N(\theta(r-s))$.  The latter double sum is a quadratic
form in the Gram matrix of $(e_K(\theta rz))_{r,z}$.  Its norm is bounded by
\eqref{eq:g7-bilinear}.  Gaussian tails in $z$ are handled by dyadic boxes;
monotonicity of $\mathcal H_T$ in $T$ and Gaussian decay make their total
convergent.  This proves \eqref{eq:g7-kernel}.
For Fourier transforms of smooth compactly supported weights, replace the
Gaussian kernel by a rapidly decreasing positive majorant; a dyadic
superposition of Gaussian kernels proves the same bound.  When the kernel
parameter decreases from $N$ to $N/2^r$, use that
$\mathcal H_t(\theta)/t$ is nonincreasing in $t$, directly from its
definition; consequently $(JN/2^r)/\mathcal H_{JN/2^r}(\theta)$
is at most $JN/\mathcal H_{JN}(\theta)$.  The rapid tail weights
absorb the corresponding powers of $2^r$.

\subsection{The cubic family}

Let $j=k-3$, and put
\[
 M_i=X^{\phi_i},\quad H_i=XM_i^{-k},\quad
 M=\prod_{i\le j}M_i,\quad H=\prod_{i\le j}H_i,\quad
 Q=X/M,\quad T=Q^k=X^3H.
\]
We require $T\gg_KX^2$, which will always hold in the septic application.
The relevant auxiliary polynomial has the form
\[
 \Psi_j(z;\mathbf h;\mathbf m)
   =a(\mathbf h)\{c_k z^3+B(\mathbf h,\mathbf m)z\},\qquad
 a(\mathbf h)=h_1\cdots h_j,
\]
where $c_k\ne0$ is fixed and integral, and $B$ is integral after multiplication
by a fixed denominator depending only on $k$.  Such fixed denominators can be
cleared once, changing only finitely many bad prime ideals and fixed lattices.
We sum over nonzero $h_i$ in boxes of volume $O(H_i)$ and over $m_i$ in boxes
of volume $O(M_iX^\epsilon)$.  Retain the actual small-support restrictions
from efficient differencing.  In particular $z$ lies in a fixed small box
around a nonzero archimedean center of size $X^{1/n}$, and
$|B_v|\ll_KX^{2/n}$ uniformly.  The Hessian of the cubic phase in $z$ is
nonzero on that box at every archimedean place.

Let $F_j(\alpha)$ be this sum, allowing a bounded smooth weight in $z$ with
uniformly bounded scaled derivatives and any added linear archimedean phase.
To arrange this smoothness without a boundary hypothesis, start with a
nonnegative $C^\infty$ weight that is at least one on the distinguished
ordinary root box and supported in a slightly larger tiny box.  The initial
solution count is thereby enlarged.  Every Cauchy or differencing step
produces products of translates of this weight.  The affine changes in $z$
have coefficients bounded in terms of the fixed differencing depth, so
all derivatives remain uniformly bounded at scale $X^{1/n}$.  Its support
retains the small-increment restrictions and the nonzero center.  Extra
indicators may be enlarged to such products before the last absolute-value
estimate, since the underlying orthogonality count has nonnegative weights.
Alternatively, ordinary Fourier completion of sharp box intersections costs
only $X^\epsilon$; the estimates below are uniform in its linear frequencies.

Define
\[
 \mathfrak M=\{\alpha:\mathcal H_T(\alpha)\le C_KX\},\qquad
 \mathfrak m=\mathbb T_K\setminus\mathfrak M.
\]
Take a fixed sufficiently small $C_K$ when disjointness is wanted; changing
it by a fixed factor changes no bound below.

\begin{proposition}[Cubic auxiliary major and minor estimates]
\label{prop:g7-cubic}
Under the preceding assumptions,
\begin{align}
 \sup_{\alpha\in\mathfrak m}|F_j(\alpha)|
       &\ll_{K,\epsilon}X^{3/4+\epsilon}HM,
       \label{eq:g7-cubic-minor}\\
 \int_{\mathfrak M}|F_j(\alpha)|^4\,d\alpha
       &\ll_{K,\epsilon}(XHM)^4T^{-1}X^\epsilon.
       \label{eq:g7-cubic-major}
\end{align}
\end{proposition}

\begin{proof}[Minor-arc estimate]
Apply two ordinary differences in $z$ to a cubic summand and use H\"older
for the $h,m$ sums.  Zero differences contribute $O(X^3H)$ before the final
H\"older factor.  Every remaining linear coefficient is a fixed nonzero
multiple of $h_1\cdots h_j u_1u_2$.  For any prescribed nonzero value,
the number of tuples $(\mathbf h,u_1,u_2)$ producing it is $O(X^\epsilon)$,
by ideal divisors followed by the bounded-height unit count.  Such values
lie in a box of volume $O(HX^2)=O(T/X)$.
The linear sums are bounded by positive kernels of width $X^{-1/n}$.
Equation~\eqref{eq:g7-kernel}, with $J=T/X$ and $N=X$, bounds their total by
\[
 X^\epsilon\left(T/X+X+T/\mathcal H_T(\alpha)\right)
       \ll_K X^\epsilon T/X
 \qquad(\alpha\in\mathfrak m).
\]
Consequently the sum over $h$ of the fourth powers of the cubic sums is
$O(X^{3+\epsilon}H)$, uniformly in $m$.  H\"older over the $HM$ parameter
choices now gives $|F_j|^4\ll X^{3+\epsilon}(HM)^4$, as asserted.
\end{proof}

\subsection{Cubic complete sums and the major arcs}

For a good prime ideal $\mathfrak p$, $q=N\mathfrak p$, define
\[
 w_3(\mathfrak p^{3u+1})=q^{-u-1/2},\qquad
 w_3(\mathfrak p^{3u+2})=w_3(\mathfrak p^{3u+3})=q^{-u-1}
 \quad(u\ge0).
\]
At the finitely many bad primes, including those above $6c_k$, take
$w_3(\mathfrak p^e)=q^{-e/3}$.  Extend multiplicatively and put $w_3(1)=1$.
For any cubic polynomial with fixed unit leading coefficient at good primes,
and any primitive additive character modulo an ideal $\mathfrak q$,
its normalized complete sum is
\begin{equation}\label{eq:g7-complete}
 \ll_{K,\epsilon}(N\mathfrak q)^\epsilon w_3(\mathfrak q),
\end{equation}
uniformly in all integral lower coefficients, including a linear twist.
For completeness, at good primes translate to $az^3+bz$.
At modulus $\mathfrak p$ the finite-field Weil bound gives $2q^{-1/2}$.
For higher powers, stationary residue classes satisfy $3az^2+b=0$.
Simple classes give quadratic Gauss bounds $O(q^{-e/2})$.
A multiple class is $z\equiv0$, with $\mathfrak p\mid b$.
If $\mathfrak p\parallel b$, its contribution is zero for $e\ge3$;
if $\mathfrak p^2\mid b$, substitution $z=\pi y$ gives precisely the
factor $q^{-1}$ times the normalized sum at level $e-3$.
The base cases $e=2,3$ give $q^{-1}$.  These recurrences prove the displayed
weight, with an absolute factor per prime absorbed by $(N\mathfrak q)^\epsilon$.
At bad primes one may apply the local-field van der Corput estimate to the
third derivative.  Proposition~3.3 of Cluckers \cite{Cluckers2011}, \emph{Expositiones Mathematicae}
29 (2011), 371--386, gives $O_K(q^{-e/3})$ uniformly here: after division by
$6c_k$ the third derivative has fixed absolute value, and all other coefficients
have bounded local Gauss norm.  This is a ready-made local-field result.

Two consequences, both elementary Euler-product calculations, are
\begin{align}
 \sum_{N\mathfrak q\le Y}(N\mathfrak q)w_3(\mathfrak q)^4
       &\ll_{K,\epsilon}Y^\epsilon,\label{eq:g7-wsum}\\
 \sum_{\mathfrak d\mid\mathfrak q}
       \frac{w_3(\mathfrak q/\mathfrak d)}{N\mathfrak d}
       &\ll_{K,\epsilon}(N\mathfrak q)^\epsilon w_3(\mathfrak q).
       \label{eq:g7-wconv}
\end{align}
For the first, the good-prime Euler factor is $1+O(q^{-1})$ and the remaining
terms are summable; for the second the local ratios form a geometrically
convergent sum up to a bounded factor per prime.

\begin{proof}[Major-arc estimate]
Write $\alpha=\gamma+\beta$ on an arc with conductor $\mathfrak q$ and
$(N\mathfrak q)A_T(\beta)\le C_KX$.
For fixed $h,m$, put $a=h_1\cdots h_j$ and
\[
 \mathfrak q_a=\mathfrak q/(\mathfrak q+(a)).
\]
This is exactly the conductor of $a\gamma$.
Poisson summation modulo $\mathfrak q_a$, followed by
\eqref{eq:g7-complete}, bounds the cubic sum by
\begin{equation}\label{eq:g7-poisson}
 \ll X^{1+\epsilon}w_3(\mathfrak q_a)
       A_{X^3}(a\beta)^{-1/2}.
\end{equation}
Here are the archimedean details.  At a real place, the second derivative is
bounded away from zero after division by its scale, and the second derivative
test gives $(1+X^{3/n}|a_v\beta_v|)^{-1/2}$.
At a complex place the real Hessian of the real part of the holomorphic cubic
has determinant a nonzero constant times $|a_v\beta_v z_v|^2$;
two-dimensional stationary phase gives the square of this factor.
These estimates are uniform in the linear Poisson and completion frequencies.
Outside the natural gradient box repeated integration by parts gives arbitrary
decay.  Ideal packing bounds the total number of relevant dual frequencies by
\[
 O_K\left(1+\frac{N\mathfrak q_a}{X}A_{X^3}(a\beta)\right)=O_K(1),
\]
because $|a_v|\ll H^{1/n}$ and the product-arc inequality holds.
Dyadic enlargement of the gradient box, with enough integrations by parts,
proves the same bound for the tails.  This proves \eqref{eq:g7-poisson} for
the complete sum, including all Poisson modes.

It remains to average over $h$.  On an archimedean dyadic box with radii
$r_{i,v}$ for $h_i$, let $U_i=\prod_v r_{i,v}^{d_v}$.
There are no nonzero $h_i\in\mathfrak b$ in such a box unless
$U_i\gg_KN\mathfrak b$; when nonempty, ideal packing gives
$O_K(U_i/N\mathfrak b)$ choices.
Split the condition $\mathfrak d\mid(h_1\cdots h_j)$ among divisor ideals
of the individual $h_i$.  The number of splittings is $O(X^\epsilon)$.
Thus summing \eqref{eq:g7-poisson} on this dyadic box, and using
\eqref{eq:g7-wconv}, gives its trivial cardinality times
$X^\epsilon w_3(\mathfrak q)$ times the archimedean decay.
The sum over dyadic radii converges toward their largest values: at each
place the counting factor has degree $d_v$ whereas the decay removes only
$d_v/2$.  Therefore
\[
 |F_j(\gamma+\beta)|
 \ll X^{1+\epsilon}HM\,w_3(\mathfrak q) A_T(\beta)^{-1/2}.
\]
Finally $\int_{V_K}A_T(\beta)^{-2}\,d\beta\ll_KT^{-1}$,
and there are at most $N\mathfrak q$ center classes of conductor
$\mathfrak q$.  Sum the fourth power of the preceding bound using
\eqref{eq:g7-wsum}.  Overlaps, if any, only enlarge this upper bound.
This proves \eqref{eq:g7-cubic-major}.
\end{proof}

\subsection{The pure smooth terminal process}

Let $\lambda_t$ be any permissible smooth $2t$-moment exponent, extended to
real $t\ge0$ by H\"older interpolation.  For every $u\ge0$, Proposition
\ref{prop:g7-cubic} implies
\begin{equation}\label{eq:g7-M1-terminal}
 \int|F_j(\alpha)g_Q(\alpha)^{2u}|\,d\alpha
 \ll X^{1+\epsilon}HM
 \left\{X^{-1/4}Q^{\lambda_u}
       +Q^{-k/4+(3/4)\lambda_{4u/3}}\right\}.
\end{equation}
The minor term is the supremum estimate times the full smooth moment.
The major term is H\"older with exponents $4$ and $4/3$, using
$T=Q^k$.  This is the pure smooth $M_1$ process; it needs no extra ordinary
variables inside the terminal integral and no smooth major-arc asymptotic.


\section{A twelve-row exact certificate for the septic smooth moment}
\label{sec:g7-certificate}

The preceding analytic lemmas imply a particularly small certificate:
\begin{equation}\label{eq:g7-smooth28}
 \int_{\mathbb T_K}|g_X(\alpha)|^{28}\,d\alpha
       \ll_{K,\epsilon}X^{21.11765+\epsilon}
       \ll_{K,\epsilon}X^{21+59/500+\epsilon}.
\end{equation}
Here $g_X$ is the ideal-smooth seventh-power sum in a fixed number field.
Every decimal in this section is an exact terminating rational number.
We give the complete data and recurrence, so the numerical assertion is
reproducible without an external table of rational permissible exponents.

\subsection{The recurrence being certified}

Write $\lambda_t$ for a permissible $2t$-moment exponent, with
$\lambda_0=0$, $\lambda_1=1$, $\lambda_2=2$.
For nonintegral $t$, use linear interpolation between the adjacent integer
entries, justified by H\"older.  A row specifies $s,j,\phi_1,\ldots,\phi_j$;
put $\Phi_i=\sum_{r\le i}\phi_r$, $q_i=1-\Phi_i$, and $u=s-j$.
The terminal exponent $C_j$ is as follows.
For process $A_j^{(\ell)}$, put $p=2^\ell$ and
\[
 \kappa_\ell=p-\ell+(p-1)(j-6\Phi_j),\qquad
 C_j=p^{-1}\kappa_\ell+(1-p^{-1})q_j
                \lambda_{u/(1-p^{-1})}.
\]
For $B_{j,t}^{(\ell)}$, put
\[
 c=u/t,\quad a=2(1-c)-2/p,\quad b=2/p-(1-c),
 \qquad C_j=a\kappa_{\ell-1}+b\kappa_\ell+cq_j\lambda_t.
\]
The displayed rows have $a,b,c\ge0$, $a+b+c=1$ and
$a(p/2)+bp=1$, exactly the required H\"older identities.
For the pure cubic process $M_4$, equation~\eqref{eq:g7-M1-terminal} gives
\[
 C_4=5-6\Phi_4+
 \max\left\{-\frac14+q_4\lambda_{s-4},\quad
      q_4\left(-\frac74+\frac34\lambda_{4(s-4)/3}\right)\right\}.
\]

Now proceed backwards for $i=j-1,\ldots,1$.  Put $u_i=s-i-1$ and
$t_i=s-i+1$, and replace the current $C_{i+1}$ by
\begin{align*}
 D_i&=1+i-7\Phi_i+\Phi_{i+1}+q_{i+1}\lambda_{u_i},\\
 E_i&=\tfrac12\{q_i\lambda_{t_i}+i-6\Phi_i
                       +(2u_i-1)\phi_{i+1}\},\\
 C_i&=E_i+\tfrac12\max(D_i,C_{i+1}).
\end{align*}
The fundamental efficient-differencing inequality and the intermediate
Cauchy inequalities then produce the permissible exponent
\begin{equation}\label{eq:g7-map}
 \mathcal F_s(\boldsymbol\lambda)
  =(2s-3)\phi_1+
       \max\{1+\phi_1+(1-\phi_1)\lambda_{s-1},C_1\}.
\end{equation}
All quantities in this recurrence are explicit rational combinations of
entries of $\boldsymbol\lambda$ and a finite number of maxima.
In iteration use $\max\{s,\mathcal F_s(\boldsymbol\lambda)\}$, and,
if desired, its minimum with the trivial exponent $2s$.  This preserves
the hypothesis $\lambda_s\ge s$ needed for the fundamental absorption.
For the certificate below $s<\mathcal F_s(\boldsymbol b)<b_s<2s$,
so these harmless clamps do not change any tabulated check.

\subsection{All certificate data}

In the table, $m_i=10^5\phi_i$.  Dashes mean that the corresponding parameter
does not occur.  Put $b_0=0$, $b_1=1$, $b_2=2$, and take $b_s$ from the table.
The last column is an integer lower bound for
$10^8\{b_s-\mathcal F_s(\boldsymbol b)\}$, obtained by the exact recurrence
above.  It is strictly positive in every row.

\begin{center}
\small
\begin{tabular}{r c r r r r r r}
\hline
$s$ & Process & $b_s$ & $m_1$ & $m_2$ & $m_3$ & $m_4$ & Slack bound\\
\hline
3 & $A_{1}^{(1)}$ & 3.06493 & 3196 & -- & -- & -- & 52177 \\
4 & $A_{1}^{(2)}$ & 4.26513 & 6818 & -- & -- & -- & 8692 \\
5 & $B_{2,5}^{(2)}$ & 5.59014 & 8699 & 3541 & -- & -- & 10185 \\
6 & $A_{2}^{(2)}$ & 7.01542 & 9641 & 6937 & -- & -- & 12539 \\
7 & $B_{3,6}^{(2)}$ & 8.54214 & 10564 & 8803 & 4059 & -- & 14896 \\
8 & $A_{3}^{(2)}$ & 10.15371 & 11203 & 9889 & 6902 & -- & 12594 \\
9 & $B_{4,8}^{(2)}$ & 11.84806 & 11874 & 10797 & 8946 & 4151 & 16152 \\
10 & $A_{4}^{(2)}$ & 13.60673 & 12330 & 11453 & 9898 & 6609 & 13579 \\
11 & $A_{4}^{(2)}$ & 15.42553 & 12805 & 12028 & 10870 & 8586 & 14177 \\
12 & $A_{4}^{(2)}$ & 17.29695 & 13252 & 12659 & 10684 & 10188 & 17123 \\
13 & $M_4$ & 19.19477 & 13392 & 12953 & 12116 & 10714 & 14754 \\
14 & $M_4$ & 21.11765 & 13559 & 13070 & 12326 & 10846 & 15886 \\
\hline
\end{tabular}
\end{center}

Every $A$ and $B$ row satisfies the ordering and pair inequalities
\eqref{eq:g7-aux-conditions}; these can be checked directly from the integers
$m_i$.  Every $M_4$ row satisfies $7q_4>2$, the hypothesis of the cubic
dissection.  All required interpolation arguments lie in $[0,14]$.

\subsection{Why a strict supersolution proves permissible exponents}

The map $\mathcal F=(\mathcal F_s)_{3\le s\le14}$, including the
constant clamps just specified, is monotone in the entries
of its argument and is a contraction in the sup norm, with constant at most
$97/100$.  Here is a direct check of the latter claim, avoiding numerical
spectral-radius arguments.  Each terminal expression has total coefficient
of the $\lambda$ variables at most $q_j$.  At a backwards step, the
coefficients contributed by $E_i$ sum to $q_i/2$, while those contributed by
either $D_i/2$ or $C_{i+1}/2$ sum to at most $q_{i+1}/2$.
Since $q_{i+1}\le q_i$, the resulting sum is at most $q_i$.
At the initial maximum both branches therefore have coefficient sum at most
$q_1=1-\phi_1$.  Linear interpolation has nonnegative coefficients summing
to one, and a maximum preserves a common Lipschitz bound.
Every row has $\phi_1\ge0.03196>0.03$, so $q_1<97/100$.

Begin with the permissible vector
\[
 \lambda_0^{(0)}=0,\qquad \lambda_1^{(0)}=1,\qquad
 \lambda_2^{(0)}=2,\qquad \lambda_s^{(0)}=2s\quad(3\le s\le14).
\]
The first two nontrivial entries follow from the second and fourth moments;
the remaining entries use only the trivial bound.  Higher Hua estimates can
accelerate the initial numerical search but are not needed for this proof.
There is a finite $D>0$ with
$\boldsymbol\lambda^{(0)}\le\boldsymbol b+D\mathbf1$.
Applying the twelve rows simultaneously, their strict supersolution property
and monotonicity give
\[
 \boldsymbol\lambda^{(r)}
 \le \boldsymbol b+(97/100)^rD\mathbf1.
\]
Each finite iterate is permissible by the proved analytic lemmas.
For a prescribed $\epsilon>0$, choose finite $r$ with
$(97/100)^rD<\epsilon/2$ and choose the smoothness parameter small enough
for those finitely many operations.  The remaining permissible-exponent
$\epsilon/2$ loss proves the bound with exponent $b_s+\epsilon$.
This proves \eqref{eq:g7-smooth28}.

\subsection{The thirty-one-variable analytic interface}

Take the slightly weaker convenient excess $\Delta=59/500=0.118$ in
\eqref{eq:g7-smooth28}.  The localized ordinary critical moment supplied by
the number-field VMVT framework has order $56$ and broad-minor exponent
$56-7-1=48$.  Thus H\"older, with
$6/56+25/28=1$, gives
\[
 \int_{\mathfrak m_*}|f_X(\alpha)|^6|g_X(\alpha)|^{25}\,d\alpha
 \ll X^{(6/56)48+(25/28)(21+\Delta)+\epsilon}
 =X^{24-1/560+\epsilon}.
\]
The smooth bootstrap from moment $28$ to moment $32$ has positive margin
\[
 (32-28)-(56-32)\Delta=4-24\Delta=146/125>0.
\]
Consequently the established bootstrap supplies the required $32$nd moment
for the broad-to-narrow annulus.  There use H\"older with ordinary exponent
$p=192/7$, since $6/p+25/32=1$ and $p>2\cdot7$.
The ordinary major-arc pruning estimate therefore saves a positive power.
These statements provide exactly the analytic input for the product-root major-arc endgame in Section~\ref{sec:product-endgame} with six ordinary and twenty-five smooth roots.
No thirtieth smooth moment is needed.


% Working proof, 2026-11-24.  Root agent owns this file.
\section{A product-root endgame over an arbitrary fixed number field}
\label{sec:product-endgame}

Throughout this section $K$ is fixed, $n=[K:\mathbb Q]$, $V=K\otimes\mathbb R$,
$X=P^n$, and $\mathbb T_K=V/\mathfrak D^{-1}$ has probability Haar measure.
All exponents below are exponents of the root volume $X$.
The constants and the sufficiently-large threshold may depend on $K$.
For $\gamma\in K/\mathfrak D^{-1}$ write
\[
 \mathfrak a_\gamma=\{a\in\mathcal O_K:a\gamma\in\mathfrak D^{-1}\},
 \quad Q_\gamma=N\mathfrak a_\gamma,
 \quad A_P(\beta)=\prod_{v\mid\infty}\max(1,P^k|\beta_v|)^{d_v}.
\]
Let $\mathfrak M_*$ be the union of the product arcs
$Q_\gamma A_P(\alpha-\gamma)\le C_KX$, and let $\mathfrak m_*$ be its complement.
We use the ordinary weighted sums $f$ from the preceding analytic framework.
Their weights are nonnegative, smooth, supported in compact annuli away from zero,
and at least one on the support of the smooth sum under consideration.
The two ordinary estimates from Section~\ref{sec:ordinary-bundle} used here are
\begin{align}
 \int_{\mathfrak m_*}|f|^{w}&\ll X^{w-k-1+\varepsilon},\qquad w=k(k+1),
 \label{eq:ordinary-critical-here}\\
 \int_{\mathfrak M_*\setminus\mathfrak N_c}|f|^p
 &\ll X^{p-k-c(p/k-2)/2}\qquad(p>2k),
 \label{eq:ordinary-pruning-here}
\end{align}
where $\mathfrak N_c$ is the union of $Q_\gamma A_P(\alpha-\gamma)\le X^c$.
Also $\int_{\mathfrak M_*}|f|^p\ll X^{p-k}$ for $p>2k$.
These are number-field estimates, not rational estimates assumed to transfer.

Let $g$ be the sum over algebraic integers in a fixed bounded annulus of scale $P$
whose prime ideal factors have norm at most $X^\eta$; extra fixed archimedean
restrictions may be imposed.  Write $U_r=\int_{\mathbb T_K}|g|^r$.
The smoothness exponent is sufficiently small for the finitely many moment estimates
being used, and is then held fixed.

\begin{proposition}[Critical-moment absorption]
\label{prop:bootstrap-here}
Suppose $u<S<w$ are even integers, $S>2k$, and
\[
 U_u\ll X^{u-k+\Delta+\varepsilon},
 \qquad D=(S-u)-(w-S)\Delta>0.
\]
Then $U_S\ll X^{S-k}$, with no $X^\varepsilon$ loss.
\end{proposition}
\begin{proof}
Orthogonality and nonnegative coefficients give
$U_S\le\int |f|^{S-u}|g|^u$.
On $\mathfrak m_*$ apply H\"older with weights
$(S-u)/w$, $(w-S)/w$, and $u/w$ to the moments of $f$, $g$, and $g$
of orders $w$, $u$, and $S$, respectively.
On $\mathfrak M_*$ use weights $(S-u)/S$ and $u/S$ and the ordinary
major-arc $S$th moment.
Writing $T=X^{-(S-k)}U_S$ yields
\[
 T\le C\{X^{-D/w+\varepsilon}T^{u/w}+T^{u/S}\}.
\]
Choose $\varepsilon<D/(2w)$.  Both powers of $T$ are strictly smaller than one,
so this inequality bounds $T$ independently of $P$.
The argument uses no asymptotic formula for smooth algebraic integers.
\end{proof}

\subsection{Replacing smooth-number distribution by finite products}
Fix an integer $\ell>1/\eta$ and nonnegative smooth weights $w_1,\ldots,w_\ell$
in small annuli of $V$, positive at all real places.  Define
\[
 h(\alpha)=\sum_{y_1,\ldots,y_\ell\in\mathcal O_K}
 \prod_{j=1}^{\ell}w_j(y_j/P^{1/\ell})
 e\!\left(\operatorname {Tr}\bigl(\alpha(y_1\cdots y_\ell)^k\bigr)\right).
\]
Each tuple supplies one $k$th-power root $y_1\cdots y_\ell$.
Every prime ideal dividing this root has norm at most
$C_KX^{1/\ell}<X^\eta$ for sufficiently large $X$.
The support can be placed inside the annulus used to define $g$.

For any prescribed $\epsilon>0$, the coefficient $b(x)$ of a nonzero root $x$
in $h$ satisfies $0\le b(x)\ll X^\epsilon$.
Indeed, the ideals $(y_1),\ldots,(y_\ell)$ factor $(x)$, so the number of ideal
tuples is bounded by its $\ell$-fold ideal divisor function, which is $O(X^\epsilon)$.
For each ideal tuple the ratios of two admissible generators are units with
all archimedean absolute values bounded above and below by fixed constants;
there are only $O_{K,\ell}(1)$ choices.  Consequently, for every fixed even $r$,
\begin{equation}\label{eq:product-moment-domination}
 \int|h|^r\ll X^\epsilon U_r.
\end{equation}
The same argument applies to joint even moments with ordinary factors.
The symbol $\epsilon$ can be made arbitrarily small by applying the divisor bound
with a correspondingly smaller exponent.  Real moments follow by interpolation.
The order of choices matters.  First choose the error exponent in the given
smooth-moment theorem small relative to the strict minor-arc saving, then fix a
permissible $\eta$ and prove the loss-free bootstrap moment.  Next choose
$\ell>1/\eta$ and $c<1/\ell$.  Finally choose the exponent in the independent
ideal-divisor bound small relative to $c$.  Thus the annular argument does not
require the potentially circular conditions $\eta\ll\epsilon\ll\eta$.

There is a direct major-arc formula for $h$.  Put
\begin{align*}
 H_\ell(\gamma)&=Q_\gamma^{-\ell}
 \sum_{\mathbf y\bmod\mathfrak a_\gamma}
 e\!\left(\operatorname {Tr}(\gamma(y_1\cdots y_\ell)^k)\right),\\
 V_\ell(t)&=\int_{V^\ell}\prod_j w_j(z_j)
 e\!\left(\operatorname {Tr}(t(z_1\cdots z_\ell)^k)\right)d\mathbf z.
\end{align*}
For any $0<c<1/\ell$ and any fixed $B>0$, Poisson summation in the $\ell$ factors gives
\begin{equation}\label{eq:product-core}
 h(\gamma+\beta)=XH_\ell(\gamma)V_\ell(P^k\beta)+O(X^{-B})
 \qquad(Q_\gamma A_P(\beta)\le X^c).
\end{equation}
Here $B$ is arbitrary because all factor weights are smooth.  To see the uniformity,
split each factor into residue classes modulo $\mathfrak a_\gamma$.
In the resulting Poisson formula, each nonzero dual coordinate belongs to
$\mathfrak a_\gamma^{-1}\mathfrak D^{-1}$ and has norm at least $c_K/Q_\gamma$.
The normalized dual radii are $(1+P^k|\beta_v|)/P^{1/\ell}$.
Their product has volume $A_P(\beta)/X^{1/\ell}$, so every nonzero mode has
normalized maximum radius at least
$c_K(X^{1/\ell}/(Q_\gamma A_P(\beta)))^{1/n}$.
Integration by parts once a dual coordinate exceeds the corresponding phase
gradient, followed by ideal-lattice packing, gives an arbitrary power of
$Q_\gamma A_P(\beta)/X^{1/\ell}$.
The trivial bound for the complete mixed sum is sufficient for these nonzero modes.

The pure complete factors satisfy
\begin{equation}\label{eq:product-complete}
 |H_\ell(\gamma)|\ll_{K,k,\ell,\epsilon}Q_\gamma^{-1/k+\epsilon}.
\end{equation}
For a prime-power conductor $\mathfrak p^E$, condition on
$T=\sum_{j<\ell}v_\mathfrak p(y_j)$.
The probability of a given $T$ is at most
$C_\ell(T+1)^{\ell-2}q^{-T}$, and the last-factor normalized $k$th-power sum
is at most $C_{K,k}q^{-\max(E-kT,0)/k}$.
Summing gives $C_{K,k,\ell}(E+1)^{\ell-1}q^{-E/k}$.
The finitely many primes dividing $k$ or the different are absorbed in the constant;
Chinese remaindering and the ideal divisor bound prove \eqref{eq:product-complete}.
The archimedean factors are rapidly decreasing: on the supports, differentiation
in any one factor has nonvanishing gradient of size comparable to $|t_v|$.

\paragraph{The representation main term and its normalization.}
For $a+b=N$, put $L=a+b\ell$ and let
\[
 \mathcal F(\mathbf x,\mathbf y)
   =\sum_{i=1}^a x_i^k+
      \sum_{j=1}^b(y_{j,1}\cdots y_{j,\ell})^k.
\]
Orthogonality makes
$R(\nu)=\int_{\mathbb T_K}f(\alpha)^ah(\alpha)^b e_K(-\alpha\nu)\,d\alpha$
the nonnegative weighted count of these representations.
The core arcs are disjoint for sufficiently large $X$.
Indeed, distinct centers of conductor norms $Q,Q'$ satisfy
$|N(\gamma-\gamma')|\gg_K(QQ')^{-1}$, whereas an overlap forces
$|N(\gamma-\gamma')|\ll_K X^{2c-k}(QQ')^{-1}$.
Here $2c<k$, since $c<1/\ell<1$ and the exponents under consideration are at least five.

Whenever the minor and annular estimates in the next subsection provide power
savings, the two core Poisson formulae give
\[
 R(\nu)=X^{N-k}\mathfrak S(\nu)\mathfrak J(\nu/P^k)
            +o(X^{N-k}),
\]
uniformly for normalized targets in the compact set used below, where
\[
 \mathfrak S(\nu)=\sum_{\gamma\in K/\mathfrak D^{-1}}
      G(\gamma)^aH_\ell(\gamma)^b e_K(-\gamma\nu),
 \qquad
 \mathfrak J(t)=\int_V V_f(\beta)^aV_\ell(\beta)^b e_K(-\beta t)\,d\beta.
\]
Here $V_f(t)=\int_V w(z)e_K(tz^k)\,dz$ is the ordinary archimedean
factor, and all measures have the normalization fixed in
Section~\ref{sec:ordinary-bundle}.  The conductor majorant is summable
because $N>2k$; the archimedean factors decay rapidly.  These same majorants justify extending the truncated core
series and integral, with a uniform error tending to zero.
Chinese remaindering and finite-ring character orthogonality give
\[
 \mathfrak S(\nu)=\prod_{\mathfrak p}\sigma_{\mathfrak p}(\nu),\qquad
 \sigma_{\mathfrak p}(\nu)=\lim_{E\to\infty}
 q^{-E(L-1)}\#\{(\mathbf x,\mathbf y)\bmod\mathfrak p^E:
                       \mathcal F(\mathbf x,\mathbf y)\equiv\nu\},
 \quad q=N\mathfrak p.
\]
For completeness, a primitive local representation gives a target-uniform
positive lower bound.  Replace each product root $x$ by $(x,1,\ldots,1)$.
At a unit coordinate the derivative has valuation $h=v_{\mathfrak p}(k)$.
Fix the other $L-1$ variables modulo $\mathfrak p^{2h+1}$ at that exact solution.
Every further choice of those variables lifts in the distinguished coordinate
by Hensel's lemma.  Thus for every $E\ge2h+1$ the normalized count is at least
$q^{-(2h+1)(L-1)}$, independently of the target.
At all sufficiently large primes the complete-sum estimates give
$\sigma_{\mathfrak p}(\nu)=1+O(q^{1-N/k+\epsilon})$.
Choose $\epsilon<N/k-2$.  This summable good-prime tail, combined with the
positive lower bounds at the finitely many remaining primes, makes
$\mathfrak S(\nu)$ positive uniformly in $\nu$.

Finally take $P=|N\nu|^{1/(kn)}$ and balance $\nu$ using $k$th powers of
totally positive units, for example squares of units.  Their logarithmic
lattice has full rank, so the normalized targets lie in a compact set with
every embedding bounded away from zero.  Small root boxes near equal $k$th
roots, positive at the real places, give a uniformly positive singular
integral on a finite cover.  Undoing the totally positive unit preserves
the signs of all real roots and preserves their absolute norms.
Every resulting root has norm $O_K(P^n)=O_K(|N\nu|^{1/k})$, exactly as required
by the stated convention for $G_{\mathbb K}(k)$.

\subsection{The resulting numerical gates}
If the hypotheses of Proposition~\ref{prop:bootstrap-here} hold with $N=S$,
the representation integral $f^{S-u}h^u$ has minor-arc bound
\[
 \int_{\mathfrak m_*}|f|^{S-u}|h|^u
 \ll X^{S-k-D/w+\epsilon}.
\]
For the broad major arcs outside $\mathfrak N_c$, use H\"older on the ordinary
$S$th moment and the product $S$th moment.  By \eqref{eq:ordinary-pruning-here},
\eqref{eq:product-moment-domination}, and Proposition~\ref{prop:bootstrap-here},
this has a fixed power saving, once $\epsilon$ is chosen sufficiently small.
Equation~\eqref{eq:product-core} then supplies a positive core main term of
order $X^{S-k}$.  Hence the stated local hypothesis implies $G_{\mathbb K}(k)\le S$.

In particular, the following are sufficient smooth-moment gates:
\[
\begin{array}{c|c|c|c|c}
 k&S&u&w&\hbox{required excess}\\\hline
6&24&20&42&\Delta_{20}<2/9\\
8&38&34&72&\Delta_{34}<2/17\\
9&46&40&90&\Delta_{40}<3/22
\end{array}
\]
The eighth-power field moment in Theorem~\ref{thm:eighth-moment} satisfies
the middle gate. The ninth-power gate remains to be established over $K$;
a favorable rational numerical certificate alone does not establish it.

For $k=7$, a single estimate $U_{28}\ll X^{21+\Delta+\epsilon}$ with
$\Delta<3/25$ suffices.  First bootstrap $U_{32}\ll X^{25}$,
using $(S,u,w)=(32,28,56)$ and $\Delta<1/6$.
Then use six ordinary roots and twenty-five product roots.  The minor-arc
H\"older weights $6/56$ and $25/28$ add to one and give the saving
\[
 \delta=\frac{3-25\Delta}{28}>0.
\]
On the broad major-arc annulus take ordinary order $p=192/7>14$ and product
order $32$, with weights $7/32$ and $25/32$.
The same core argument proves $G_{\mathbb K}(7)\le31$ under the primitive local hypothesis.

For $k=5$, the joint estimate proved in Section~\ref{sec:k5-certificate}, namely
$\int|f|^2|g|^{14}\ll X^{11+\Delta+\epsilon}$ with $\Delta<1/13$
gives the seventeen-root minor estimate for $f^4h^{13}$ by weights
$1/14$ and $13/14$ on the ordinary thirtieth moment and that joint moment.
Coefficientwise domination gives $U_{16}\le\int|f|^2|g|^{14}$.
Bootstrap $U_{18}$ from this $U_{16}$ estimate first; the same product-weight and power-core
argument handles the remaining arcs, choosing ordinary order $72/5>10$
and product order $18$.  This formulation keeps the joint-moment hypothesis
distinct from a pure sixteenth-moment hypothesis.


% Auxiliary result for the higher-power transfer investigation.
% All parameters H,J,N,V below measure volumes, not side lengths.
\section{A number-field cubic coefficient-sampling lemma}
\label{sec:coefficient-sampling-common}

Let $K$ be fixed, $n=[K:\mathbb Q]$, $V_K=K\otimes\mathbb R$,
$\mathbb T_K=V_K/\mathfrak D^{-1}$, and
$E(z)=\exp(2\pi i\operatorname{Tr}_{K/\mathbb Q}z)$.
For each archimedean place write $d_v=[K_v:\mathbb R]$.
For $T\ge1$, let $\mathcal B(T)$ be the algebraic integers satisfying
$|x_v|\le C T^{1/n}$ at every archimedean place, where the fixed
constant $C$ may be enlarged finitely many times.
Thus $\#\mathcal B(T)\ll_K T$.

Put
\begin{align*}
 Q_\gamma&=N\{a\in\mathcal O_K:a\gamma\in\mathfrak D^{-1}\},
 &A_T(\beta)&=\prod_{v\mid\infty}
              \max(1,T^{1/n}|\beta_v|)^{d_v},\\
 \mathcal H_T(\alpha)&=\min_{\gamma\in K}Q_\gamma A_T(\alpha-\gamma).
\end{align*}
The minimum includes all lifts of centers and is invariant modulo
$\mathfrak D^{-1}$. It exists: any fixed upper bound on the expression
bounds the conductor norms, leaves finitely many center classes, and
bounds every archimedean error.

\begin{lemma}[Product-height Dirichlet approximation and separation]
\label{lem:g789-height}
For $T\ge1$ one has $\mathcal H_T(\alpha)\ll_K T^{1/2}$.
Moreover there is $c_K>0$ such that if
\[
 \alpha=\gamma_i+\beta_i,\qquad
 (Q_{\gamma_1}A_T(\beta_1))
 (Q_{\gamma_2}A_T(\beta_2))<c_KT,
\]
then the two lifted centers are equal: $\gamma_1=\gamma_2$.
\end{lemma}
\begin{proof}
Apply Minkowski's theorem in $V_K^2$ to the lattice
$\{(b,a-\alpha b):b\in\mathcal O_K,\ a\in\mathfrak D^{-1}\}$,
whose covolume is fixed independently of $\alpha$, using the bounds
\[
 |b_v|\le C_KT^{1/(2n)},\qquad
 |a_v-\alpha_vb_v|\le C_KT^{-1/(2n)}.
\]
For sufficiently large $T$ the resulting nonzero pair has $b\ne0$,
because a nonzero element of $\mathfrak D^{-1}$ has absolute norm
at least $(N\mathfrak D)^{-1}$. Set $\gamma=a/b$.
Then $Q_\gamma\le|Nb|$, and
\[
 A_T(\alpha-\gamma)
 \le \prod_v(C_KT^{1/(2n)}/|b_v|)^{d_v}
 \ll_K T^{1/2}/|Nb|.
\]
Bounded $T$ is harmless. For separation, a nonzero difference of the
two centers has norm at least
$(Q_{\gamma_1}Q_{\gamma_2}N\mathfrak D)^{-1}$, whereas
\[
 |N(\beta_1-\beta_2)|\le2^nT^{-1}A_T(\beta_1)A_T(\beta_2).
\]
Choosing $c_K$ sufficiently small proves the assertion.
\end{proof}

\begin{lemma}[A bilinear sampling bound]
\label{lem:g789-bilinear}
If $J,N\ge1$, $c(a)$ is supported on $\mathcal B(N)$, and
$\theta\in V_K$, then
\[
 \sum_{z\in\mathcal B(J)}
       \left|\sum_a c(a)E(\theta za)\right|^2
 \ll_K\left(J+N+\frac{JN}{\mathcal H_{JN}(\theta)}\right)
                         \sum_a|c(a)|^2.
\]
The same assertion holds for arbitrary subsets of the displayed
boxes and arbitrary fixed translates of those boxes.
\end{lemma}
\begin{proof}
First suppose $J\le N$. The multidimensional elementary large-sieve
inequality bounds the left side by $O_K(NM)\sum|c(a)|^2$, where $M$
is the maximum number of sample frequencies $\theta z$ in a translate
of an archimedean box of side $O_K(N^{-1/n})$ modulo
$\mathfrak D^{-1}$. One direct proof chooses a positive Gaussian
majorant for the coefficient box, applies Poisson summation to its
Gram matrix, and bounds its operator norm by the maximum row sum.
Covering successive frequency annuli by boxes of side $N^{-1/n}$
and using Gaussian decay bounds that row sum by $O_K(NM)$.

Partition the $z$ box into a fixed number of smaller boxes if
necessary. The affine field-line lemma in Section~\ref{sec:ordinary-bundle}
applies to the close pairs $(z,a)$, because their volume product is
$O_K(J/N)$, made sufficiently small by this fixed subdivision.
It gives either at most one point or a slope $\gamma\in K$ and
\[
 M\ll_K1+Q_\gamma^{-1}
       \prod_v\min\{J^{1/n},N^{-1/n}/|\theta_v-\gamma_v|\}^{d_v}
 =O_K\left(1+\frac{J}{Q_\gamma A_{JN}(\theta-\gamma)}\right).
\]
This is at most $O_K(1+J/\mathcal H_{JN}(\theta))$.
The large-sieve bound proves the result when $J\le N$.
When $J>N$, duality of the finite matrix
$(E(\theta za))_{z,a}$ interchanges the two boxes and proves the
bound with $J$ in place of $N$. Fixed translates merely insert
unit-modulus row and column factors; subsets can be handled by
zero extension or deletion of rows.
\end{proof}

\begin{lemma}[Counting small norms in an ideal]
\label{lem:g789-smallnorm}
Let $\mathfrak b$ be an integral ideal, $b=N\mathfrak b$, and
$H\ge1$. For $U\ge b$ one has
\[
 \#\{h\in\mathfrak b\cap\mathcal B(H):0<|Nh|\le U\}
       \ll_{K,\varepsilon}H^\varepsilon U/b.
\]
There are no such nonzero points when $U<b$.
\end{lemma}
\begin{proof}
For a counted element, $(h)\mathfrak b^{-1}$ is an integral ideal
of norm at most $U/b$. There are $O_K(U/b)$ such ideals.
For any fixed principal ideal $(h)$, its generators in the box
number $O_K((1+\log H)^{r_1+r_2-1})$: upper bounds for all conjugates,
together with $|Nh|\ge1$, confine their logarithmic vectors to a box
of side $O_K(1+\log H)$, and the unit logarithms form a fixed lattice.
This proves the asserted estimate, including uniformity in
$\mathfrak b$.
\end{proof}

\begin{proposition}[Refined cubic coefficient sampling over $K$]
\label{prop:g789-sampling}
Let $H,J,N\ge1$, put $V=H^3JN$, and choose
\[
 1\le Y\le c_K\min\{J,N,(JN/H^3)^{1/2}\}.
\]
Choose a lifted approximation $\alpha=\gamma+\beta$ with
$L=Q_\gamma A_V(\beta)\ll_K V^{1/2}$, as in
Lemma~\ref{lem:g789-height}. If $c(a)$ is supported on $\mathcal B(N)$,
then
\[
 \sum_{\substack{h\in\mathcal B(H)\\h\ne0}}
 \sum_{z\in\mathcal B(J)}
       \left|\sum_a c(a)E(\alpha h^3za)\right|^2
 \ll_{K,\varepsilon}
 V^\varepsilon JHN\left(Y^{-1}+\frac{H^2}{L}\right)
                          \sum_a|c(a)|^2.
\]
In particular $L$ can be replaced by the smaller intrinsic height
$\mathcal H_V(\alpha)$. The exclusion of $h=0$ is essential; in
applications the variables are products of nonzero shifts.
\end{proposition}
\begin{proof}
Write $W=JN$ and apply Lemma~\ref{lem:g789-bilinear} with
$\theta=\alpha h^3$. Since $Y\le\min(J,N)$, the contribution
of $J+N$, and that of all $h$ with $\mathcal H_W(\alpha h^3)\ge Y$,
is $O_K(JHN/Y)\sum|c(a)|^2$.

For each remaining $h$, choose
$\alpha h^3=\gamma_h+\beta_h$ with
$Q_{\gamma_h}A_W(\beta_h)<Y$. The center $\gamma_h/h^3$ is an
approximation to $\alpha$ at scale $V$ with height at most $C_KH^3Y$.
Indeed its conductor norm is at most $|Nh|^3Q_{\gamma_h}$, and
\[
 |Nh|^3 A_V(\beta_h/h^3)
 =\prod_v\max\{|h_v|^3,V^{1/n}|(\beta_h)_v|\}^{d_v}
 \ll_K H^3 A_W(\beta_h).
\]
The upper bound on $Y$ makes the product of this height and $L$
smaller than the separation constant times $V$.
Lemma~\ref{lem:g789-height} therefore gives
$\gamma_h=h^3\gamma$ and $\beta_h=h^3\beta$ exactly.

Let $\mathfrak a$ be the conductor of $\gamma$, $Q=N\mathfrak a$,
and set $A=A_V(\beta)$. The conductor of $h^3\gamma$ is
$\mathfrak a/(\mathfrak a+(h)^3)$. Group the nonzero $h$ according
to the ideal $\mathfrak d=\mathfrak a+(h)^3$.
There are $O_{K,\varepsilon}(V^\varepsilon)$ possible divisor ideals.
For each fixed class put
\[
 \mathfrak b=\prod_{\mathfrak p}
       \mathfrak p^{\lceil v_\mathfrak p(\mathfrak d)/3\rceil},
 \qquad b=N\mathfrak b.
\]
Then $h\in\mathfrak b$ and $N\mathfrak d\le b^3$.
Since $|h_v|\ll_KH^{1/n}$ at every place, one has
\[
 A_W(h^3\beta)\gg_K
       \max\{1,A|Nh|^3/H^3\}.
\]
Consequently the reciprocal-height sum in this class is at most
\[
 \frac{b^3}{Q}
 \sum_{0\ne h\in\mathfrak b\cap\mathcal B(H)}
       \min\{1,C_KH^3/(A|Nh|^3)\}.
\]
Put $T=C_KH/A^{1/3}$, enlarging the fixed constant if necessary.
Lemma~\ref{lem:g789-smallnorm} and dyadic summation bound the inner
sum by
\[
 O_{K,\varepsilon}\left(H^\varepsilon
 \begin{cases}T/b,&T\ge b,\\(T/b)^3,&T<b.\end{cases}\right).
\]
In the first case $b^2T\le T^3$, and in the second case the factor
$b^3$ cancels. Thus either case contributes at most
$O_{K,\varepsilon}(V^\varepsilon H^3/(QA))$.
Summing the divisor classes and recalling $QA=L$ proves
\[
 \sum_{h:\,\mathcal H_W(\alpha h^3)<Y}
       \mathcal H_W(\alpha h^3)^{-1}
 \ll_{K,\varepsilon}V^\varepsilon H^3/L.
\]
Multiplication by $JN\sum|c(a)|^2$ completes the proof.
\end{proof}

\paragraph{Scope.}
This establishes the number-field counterpart of the refined sampling
estimate used in the project's eighth- and ninth-power certificates.
It does not by itself transfer their permissible smooth moments:
the initial efficient-differencing inequalities, the auxiliary seed
moments, and the restricted-height smooth mean theorem must still be
established for the chosen number-field smooth sets.


\section{An auxiliary degree-ten moment over every number field}
\label{sec:tenth-moment}

The fifth-power transfer uses a degree-ten coefficient moment.  The following
bound is obtained over the field itself and does not import the rational
coefficient table.
\begin{proposition}\label{prop:tenth44}
Fix a number field $K$.  Write $X=P^{[K:\mathbb Q]}$ and let
$g_{10}(\alpha;X,R)$ be the ideal-smooth tenth-power sum in a fixed bounded
archimedean box of scale $P$.  For every $\epsilon>0$ there is $\eta>0$ such that
\[
 \int_{\mathbb T_K}|g_{10}(\alpha;X,R)|^{44}\,d\alpha
             \ll_{K,\epsilon}X^{34.227+\epsilon}
       \qquad(1\le R\le X^\eta).
\]
In particular the excess $1/3$ is permissible at order forty-four.
\end{proposition}

Only the fundamental principal-block differencing lemma, the ordinary Hua
second and fourth moments, the auxiliary estimates proved next, and the
already established cubic dissection are required.  The fixed number of
variables in this moment assertion is independent of the field; implicit
constants and the allowable smoothness parameter may depend on it.

\subsection{The auxiliary second and fourth moments needed here}

Set $k=10$, $M_i=X^{\phi_i}$, $H_i=XM_i^{-10}$,
$\Phi_j=\sum_{i\le j}\phi_i$, and $U_j=\prod_{i\le j}H_iM_i$.
Use the centered polynomials $\Psi_j$ and the small ordinary supports in the
preceding efficient-differencing argument.  The required auxiliary estimates are
\begin{equation}\label{eq:tenth-aux}
 \int |F_j|^2\ll X^{1+\epsilon}U_j,
 \qquad
 \int |F_j|^4\ll X^{2+\epsilon}U_j^3\qquad(1\le j\le6).
\end{equation}
Assume the parameters are ordered,
$0<\phi_j\le\cdots\le\phi_1<1/10$, and
\begin{equation}\label{eq:tenth-pairs}
 \Phi_I+10(\phi_{I-1}+\phi_I)<2\qquad(3\le I\le j).
\end{equation}
For each second moment whose residual degree is even and at least six,
add the condition $\Phi_j<1/3$.  For a fourth moment impose that same
condition if the once-ordinarily-differenced residual degree is even and
at least six.  Every parameter row below meets these requirements.

First consider coefficient coincidences.  Retaining just the first necessary
coefficient equation reduces their count to
\[
              \sum_{i\le j}h_i^2(m_i^{20}-n_i^{20})=0.
\]
Write $a_i=m_i^{20}-n_i^{20}$.  Remove indices with $a_i=0$; each contributes
$O(H_iM_i)$, since a nonzero element has at most twenty twentieth roots.
For two remaining indices, put
$A_i(\alpha)=\sum_{h_i}|\sum_{m_i}e_K(\alpha h_i^2m_i^{20})|^2$.
The fourth Hua moment of degree twenty gives
$\int A_i^2\ll H_i^2M_i^{2+\epsilon}$, so Cauchy bounds the two-index count
by $X^\epsilon H_1H_2M_1M_2$.
For three nonzero indices, fix two triples $(h_i,m_i,n_i)$.  Their sum is
nonzero whenever the third term is nonzero.  The factorization
\[
       h^2(m^{10}-n^{10})(m^{10}+n^{10})=c\ne0
\]
and the algebraic-integer divisor bound determine the remaining triple in
$O(X^\epsilon)$ ways.  Taking the geometric mean over the three choices of
free index gives
$X^\epsilon(\prod H_i)^{2/3}(\prod M_i)^{4/3}$, at most
$X^\epsilon\prod H_iM_i$ under \eqref{eq:tenth-pairs}.

For $r\ge4$ remaining nonzero indices, fix all $m_i,n_i$, and split according
as $c=\sum_{i<r-1}a_ih_i^2$ vanishes.  When $c=0$, induction and the
two-index estimate apply to the two groups separately.  When $c\ne0$, fix
the first $r-2$ shifts.  The equation
$a_{r-1}h_{r-1}^2+a_rh_r^2=-c$ has $O(X^\epsilon)$ solutions by the uniform
binary-quadratic norm lemma.  Its total cost is
\[
 X^\epsilon\Bigl(\prod_{i<r-1}H_i\Bigr)
                   \Bigl(\prod_{i\le r}M_i\Bigr)^2
 \le X^\epsilon\prod_{i\le r}H_iM_i,
\]
because $\prod_{i\le r}M_i<H_{r-1}H_r$.
The same inequality for any subset of $r$ indices follows from the one for
the first $r$: the $M_i$ decrease and the $H_i$ increase.  This completes
the induction and proves $K_j\ll X^\epsilon U_j$ for all $j\le6$.

For clarity, the even-residual curve argument needed in
\eqref{eq:tenth-aux} is recalled in a form that also handles unequal constant
terms.  The phase-value divisor argument first makes the two shift tuples
equal.  With residual degree $2d\ge6$, write
\[
 \Psi_j(z;\mathbf h;\mathbf m)
    =(\prod h_i)\sum_{r=0}^{d}c_r(\mathbf h,\mathbf m)z^{2r}.
\]
If the constant coefficients agree on the two sides, subtract them.
The remaining value is divisible by $z^2$; it is nonzero on the chosen
small supports, since all increments are small relative to the nonzero
ordinary center and the leading term dominates.  Divisor counting and
Cauchy therefore force the same $z$ on the two sides.  Unless all
coefficients agree, the resulting polynomial determines $z$ in $O(1)$
ways.  Its cost is $O(X^\epsilon HM^2)$, where $H=\prod H_i$ and
$M=\prod M_i$, and the coincident-polynomial contribution is
$O(X^{1+\epsilon}K_j)$.

If the constant coefficients differ, set $x=z^2$, $y=w^2$.  The resulting
curve $F(x,y)=0$ has nonzero constant term and highest homogeneous part
$c(x^d-y^d)$ with $c\ne0$.  Irreducible factors of degree at least three
have $O(X^{2/3+\epsilon})$ integral points in the box of volume $X^2$,
by the number-field affine curve theorem of Binyamini--Cluckers--Kato \cite[Corollary~12]{BCK2025}.  Linear factors have nonzero coefficients of both variables
and nonzero constant, so substitution gives a nonzero-right-side binary
quadratic equation in $z,w$.  Quadratic factors have nondegenerate highest
homogeneous part, since the slopes in $x^d-y^d$ are distinct.  Completing
the square and diagonalizing gives the same uniform norm-equation count,
with polynomial-height coefficients after clearing denominators.  A zero
completed constant yields geometric lines; lines defined over $K$ are
handled as above, and the remaining $K$-points lie in intersections of
conjugate lines.  Components defined only over an extension similarly
contribute $O(1)$ points.  These low-degree components cost $O(X^\epsilon)$.
There are $O(X^\epsilon HM^2)$ parameter choices, so the noncoincident part
is at most $X^{2/3+\epsilon}HM^2$, which is $O(X^{1+\epsilon}U_j)$ when
$M<X^{1/3}$.  This proves the needed even-residual second moments.

For odd residual degrees the phase is divisible by $z$ and the elementary
divisor argument gives the second moment directly; for residual degrees
two and four the elementary factorizations from the sixth-power argument
apply without a curve-product restriction.  Finally one ordinary difference
and Cauchy give the fourth moment in \eqref{eq:tenth-aux}.  The new divisor
is fixed at one, its shift has volume $X$, and its coefficient term cancels
in the coefficient-coincidence equations.  Apply the same second-moment
argument to this enlarged tuple.  All residual degrees occurring in these six auxiliary bounds and
after one ordinary difference are at least three.  Thus no residual-linear
multiplicity loss is used.

The cubic dissection itself was proved with general $j=k-3$.  At $k=10$,
$j=7$ its two terminal bounds are
\[
 \int |F_7|\,|g_Q|^{2u}
 \ll X^{1+\epsilon}\Bigl(\prod H_iM_i\Bigr)
 \left\{X^{-1/4}Q^{\lambda_u}
   +Q^{-10/4+(3/4)\lambda_{4u/3}}\right\},
\]
provided $10(1-\Phi_7)>2$.  Here $\lambda_t$ denotes the exponent of the
$2t$-th smooth moment.  No auxiliary pair condition is needed in this cubic
terminal.  All its displayed parameter rows satisfy the stated strict
condition.  This supplies every analytic input for the exact recurrence below.


\subsection{A twenty-row exact certificate}

For $0\le s\le22$ write $\lambda_s$ for a permissible exponent at full
moment order $2s$, and fix $\lambda_0=0$, $\lambda_1=1$, $\lambda_2=2$.
Use ordinary linear interpolation at nonintegral indices.  For a row with
terminal depth $j$, put $u=s-j$, $\Phi_i=\sum_{r\le i}\phi_r$ and
$q_i=1-\Phi_i$.  For $A_j^{(\ell)}$ set $p=2^\ell$,
\[
 \kappa_\ell=p-\ell+(p-1)(j-9\Phi_j),\qquad
 C_j=p^{-1}\kappa_\ell+(1-p^{-1})q_j\lambda_{u/(1-p^{-1})}.
\]
For $B_{j,t}^{(\ell)}$, take
$c=u/t$, $a=2(1-c)-2/p$, $b=2/p-(1-c)$, and
\[
 C_j=a\kappa_{\ell-1}+b\kappa_\ell+cq_j\lambda_t.
\]
Every displayed row has $a,b,c\ge0$, $a+b+c=1$ and $a(p/2)+bp=1$.
For the cubic process $M_7$ set
\[
 C_7=8-9\Phi_7+
 \max\left\{-\frac14+q_7\lambda_{s-7},\quad
 q_7\left(-\frac{10}{4}+\frac34\lambda_{4(s-7)/3}\right)\right\}.
\]
Proceed backwards for $i=j-1,\ldots,1$, taking $u_i=s-i-1$, $t_i=s-i+1$:
\begin{align*}
 D_i&=1+i-10\Phi_i+\Phi_{i+1}+q_{i+1}\lambda_{u_i},\\
 E_i&=\frac12\{q_i\lambda_{t_i}+i-9\Phi_i+(2u_i-1)\phi_{i+1}\},\\
 C_i&=E_i+\frac12\max(D_i,C_{i+1}).
\end{align*}
The new permissible exponent is
\[
 \mathcal F_s(\boldsymbol\lambda)
   =(2s-3)\phi_1+
    \max\{1+\phi_1+q_1\lambda_{s-1},C_1\}.
\]
As before, clamp each output below by $s$ and above by $2s$ if necessary.
These clamps leave the following exact supersolution checks unchanged.

All decimal exponents in the table are exact rationals.  In its last column
$m_i=10^6\phi_i$.  Put $b_0=0$, $b_1=1$, $b_2=2$.
\begin{center}\small
\begin{tabular}{r c r l}
\hline
$s$&Process&$b_s$&$(m_1,\ldots,m_j)$\\\hline
3&$A_{1}^{(1)}$&3.059&$(14055)$\\
4&$A_{1}^{(2)}$&4.187&$(43306)$\\
5&$A_{2}^{(1)}$&5.396&$(54419,4144)$\\
6&$B_{2,6}^{(2)}$&6.677&$(60561,37484)$\\
7&$A_{2}^{(2)}$&8.027&$(65402,52459)$\\
8&$B_{3,7}^{(2)}$&9.445&$(69569,61371,36047)$\\
9&$A_{3}^{(2)}$&10.925&$(72921,66690,51503)$\\
10&$A_{4}^{(2)}$&12.468&$(76424,71431,62513,36792)$\\
11&$A_{4}^{(2)}$&14.068&$(79297,75107,68187,52487)$\\
12&$A_{5}^{(2)}$&15.722&$(82168,78422,73311,63489,37614)$\\
13&$A_{5}^{(2)}$&17.425&$(84622,81317,76802,69403,52369)$\\
14&$B_{6,11}^{(2)}$&19.174&$(86978,83979,80265,74636,63970,36608)$\\
15&$A_{6}^{(2)}$&20.962&$(88956,86310,82900,78089,69494,50134)$\\
16&$A_{6}^{(2)}$&22.786&$(90880,88512,85044,80465,75297,62716)$\\
17&$M_7$&24.642&$(92576,90641,88425,85972,83083,79668,75000)$\\
18&$M_7$&26.522&$(93766,92184,90188,87784,84903,81007,74999)$\\
19&$M_7$&28.424&$(94865,93377,91689,89425,86442,82205,75000)$\\
20&$M_7$&30.344&$(95784,94467,92853,90807,87806,83166,74999)$\\
21&$M_7$&32.279&$(96546,95377,93905,91877,88947,84028,74999)$\\
22&$M_7$&34.227&$(97173,96131,94778,92829,89806,84712,74999)$\\
\hline
\end{tabular}\end{center}

Exact substitution gives
$\mathcal F_s(\boldsymbol b)<b_s$ in all twenty rows; the smallest slack is
\[
           \frac{318527}{500000000}=0.000637054>0.
\]
The smallest prefix-pair slack in the $A/B$ rows is $1/250000$.
Whenever the even-residual curve argument is needed, the largest product
exponent is $\Phi_j=137539/500000=0.275078<1/3$.
The $M_7$ rows all have $10q_7>2$.  Every interpolation argument lies in
$[0,22]$.  The complete standalone program below verifies these assertions
by exact rational arithmetic, without an optimizer or external data.

For completeness, the supersolution is sufficient without borrowing any
published smooth-moment table.  Each terminal has total coefficient of the
moment variables at most $q_j$.  A backwards step has coefficient mass at
most $(q_i+q_{i+1})/2\le q_i$; both initial branches therefore have mass at
most $q_1$.  Interpolation has nonnegative coefficient mass one, and maxima
preserve a common Lipschitz bound.  All rows have $\phi_1>1/100$, so the
clamped map is monotone and has sup-norm contraction constant at most $99/100$.
Start with the permissible vector $\lambda_0=0$, $\lambda_1=1$,
$\lambda_2=2$, and $\lambda_s=2s$ for $3\le s\le22$.
For some fixed $D>0$ it is bounded by $\boldsymbol b+D\mathbf1$.
After any finite number $r$ of simultaneous updates it is bounded by
$\boldsymbol b+D(99/100)^r\mathbf1$.  Given $\epsilon>0$, choose finite
$r$ making this last error less than $\epsilon/2$, and then choose the
smoothness parameter small enough for those finitely many operations.
The remaining admissibility loss proves
$\lambda_{22}\le b_{22}=34.227$ in the usual permissible-exponent sense.
This proves Proposition~\ref{prop:tenth44}.


% Fifth-power transfer: the nontrivial coefficient lemma.
\section{A second-moment estimate for the decic coefficient sums}
\label{sec:k5-coeff2}

All size parameters in this section denote volumes.  The field is fixed,
$n=[K:\mathbb Q]$, and $\mathcal H_T$ denotes the conductor--product height
used in the preceding report.  The notation $\mathcal B(L)$ means a fixed
archimedean box with side lengths comparable to $L^{1/n}$.  Fixed fractional
lattices in place of $\mathcal O_K$ do not change any assertion below.

The following is the precise number-field form of the coefficient estimate
needed from Vaughan--Wooley's fifth-power argument \cite[Lemma~4.3]{VW1995}.  It is useful to state
it for the coefficient degree $2d$, although only $d=5$ is used.

\begin{proposition}\label{prop:k5-coeff2}
Let $d\ge3$, $J,H,M\ge1$, and $T=JH^3M^{2d}$.  Suppose
\[
 1\le U\le c_K\min\{M,JH^{-3},T^{1/(4d)},
                  (JH^3M^{1-2d})^{1/6}\}.
\]
Let $\mathcal C$ be any subset of a sufficiently small fixed box about a
nonzero point at every embedding, of scale $M^{1/n}$.  For each $h$ one may
include any fixed lower-degree additive twist in the $m$ sum, provided it
is independent of $j$.  Then
\[
 E_1(\alpha):=\sum_{\substack{h\in\mathcal B(H)\\h\ne0}}
 \sum_{j\in\mathcal B(J)}
 \left|\sum_{m\in\mathcal C}e_K(\alpha jh^3m^{2d}+b_h(m))\right|^2
 \ll_{K,d,\epsilon}T^\epsilon JHM^2
       \{\mathcal H_T(\alpha)^{-1/d}+U^{-1}\}.
\]
The same bound holds for smooth weights and for the fixed lattices arising
from class representatives.  No positivity condition at complex places is
being imposed.
\end{proposition}

\paragraph{Conventions for the proof.}
A nonzero algebraic integer has absolute norm at least one.  After clearing
a fixed denominator this also applies, with a fixed positive lower bound,
to every lattice occurring below.  An arbitrary conductor ideal can be
made principal by multiplying it by one of finitely many fixed class
representatives.  Its chosen denominator therefore has norm comparable
to the conductor norm.  In all congruence counts it is the conductor ideal,
not that auxiliary denominator, which is used.

We use anisotropic Dirichlet approximation in the following elementary
form.  Given positive radii $R_v$ with $\prod_vR_v^{d_v}\gg_K1$, Minkowski's
theorem supplies $a\in\mathfrak D^{-1}$ and nonzero $q\in\mathcal O_K$ with
\[
 |q_v|\ll_KR_v,\qquad |(q\beta-a)_v|\ll_K R_v^{-1}.
\]
The constants can be made harmless by reducing the fixed constant $c_K$
in the proposition.  All later identities follow by estimating the norm
of a cross product and showing that it is smaller than its nonzero lower
bound.  This replaces the rational proof's use of a nonzero integer of
absolute value less than one.

A smooth positive majorant of a linear sum over $\mathcal B(J)$ is a
periodic rapidly decreasing kernel $D_J$.  Dyadic boxes for its errors,
and dyadic boxes for all nonzero variables, cost only $T^\epsilon$.
Thus the estimates below may be proved with specified radii at every
embedding.  Products of these radii, with local dimensions $d_v$, give
the volume expressions displayed below.  For example the contribution
where the error kernel has normalized product size at least $U$ is
$O(T^\epsilon JHM^2/U)$.

\begin{proof}
Expanding the square and summing over $j$, using a smooth majorant before
taking absolute values, reduces the assertion to the untwisted positive
kernel sum.  Terms with $m_1=m_2$ contribute $O(JHM)$ and are acceptable.
Because the $m$ box is sufficiently small, $m_1^{2d}=m_2^{2d}$ within this
box implies $m_1=m_2$.  Put $D_0=m_2^{2d}-m_1^{2d}$ for the other pairs.

\emph{First identification: the cubic shift.}
For a fixed pair and an error box of normalized product size at most $U$,
approximate $\alpha D_0$ with denominator radii
$(J_vH_v^3/U_v)^{1/2}$, where $\prod U_v^{d_v}=U$.
If $\alpha D_0h^3$ is in that error box around $a_0\in\mathfrak D^{-1}$,
then the norm of the cross product between this approximation and $a_0/h^3$
is
\[
 \ll_K (UH^3/J)^{1/2}.
\]
It is therefore zero.  Write the resulting common center as $\gamma_{D_0}$.
Its conductor $\mathfrak q_{D_0}$ divides $(h)^3$.
The sum over such $h$, using ideal-lattice packing on dyadic boxes, is at most
\[
 T^\epsilon JH V_{D_0}^{-1/3},\qquad
 V_{D_0}=N\mathfrak q_{D_0}\,
 A_{JH^3}(\alpha D_0-\gamma_{D_0}).
\]
Here is a lattice proof of the summation, including thin boxes.  If
$\mathfrak b=\prod\mathfrak p^{\lceil v_\mathfrak p(\mathfrak q)/3\rceil}$,
then $h\in\mathfrak b$ and $N\mathfrak b\ge(N\mathfrak q)^{1/3}$.
For $A=A_{JH^3}(\beta)$ one has
$A_J(h^3\beta)\gg\max(1,A|Nh|^3/H^3)$.
The elementary ideal count
$\#\{h\in\mathfrak b\cap\mathcal B(H):0<|Nh|\le B\}
\ll H^\epsilon B/N\mathfrak b$, with no points for $B<N\mathfrak b$,
and dyadic summation give
$\sum_h A_J(h^3\beta)^{-1}\ll H^\epsilon H/(A^{1/3}N\mathfrak b)$.
This proves the displayed bound without a boundary or real-positivity assumption.
We may discard the pairs for which $V_{D_0}\ge U^3$ at the required cost.

\emph{Second identification: remove the divided difference.}
Group pairs according to their greatest common divisor ideal
$\mathfrak g=(m_1,m_2)$, of norm $G$.  Make $\mathfrak g$ principal using a
fixed class representative and choose a balanced generator $g$.
Write $m_1=gn$, $m_2=g(n+l)$.  The variables $n,n+l$ belong to fixed fractional
lattices, are primitive up to fixed primes, and have scale $N=M/G$.
The polynomial
\[
 D(n,l)=\frac{(n+l)^{2d}-n^{2d}}{l}
\]
has absolute value comparable to $N_v^{2d-1}$ at every embedding.
Indeed, both arguments are in the same tiny box away from zero, so the
leading derivative $2d n^{2d-1}$ dominates the other terms.
This is the place where localization replaces real positivity.

For fixed $g,l$, approximate $\alpha g^{2d}l$ with denominator radii
$(J_vH_v^3N_v^{2d-1})^{1/2}$, obtaining a center $\gamma_{g,l}$.
A comparison with $\gamma_{D_0}/D(n,l)$ has cross-product norm at most
\[
 \ll_K V_{D_0}
       \left(\frac{N^{2d-1}}{JH^3}\right)^{1/2}
 \le C_K U^3\left(\frac{M^{2d-1}}{JH^3}\right)^{1/2}<1.
\]
Consequently $\gamma_{D_0}=D(n,l)\gamma_{g,l}$.
If $\mathfrak q$ is the conductor of $\gamma_{g,l}$, then
$\mathfrak q_{D_0}=\mathfrak q/\mathfrak q_1$, with
$\mathfrak q_1\mid(D(n,l))$.
Moreover $n(n+l)$ is prime to $\mathfrak q_1$ outside a fixed set of primes.
The number of primitive solutions of $D(n,l)\equiv0\pmod{\mathfrak q_1}$
in residue classes is $O_{K,d,\epsilon}((N\mathfrak q_1)^{\epsilon})$.
To see the uniformity even at primes dividing $2d$, let
$t=v_{\mathfrak p}(l)$.  For a primitive solution modulo $\mathfrak p^E$,
both $n$ and $n+l$ are units.  The map
$n\mapsto r=1+l/n\pmod{\mathfrak p^{E+t}}$ is injective, and
$D(n,l)\equiv0\pmod{\mathfrak p^E}$ implies
$r^{2d}\equiv1\pmod{\mathfrak p^{E+t}}$.
The number of unit $2d$th roots modulo any power of a fixed prime is
bounded in terms of $K,d$: beyond the fixed derivative depth
$v_{\mathfrak p}(2d)$, Hensel lifting gives only a bounded number of
lifts around each root.  At good primes the bound is $2d$.
Multiplying the local bounds gives the stated $T^\epsilon$ factor.
Fixed denominators from class representatives are cleared first; their
fixed valuations only change the constant.

Put
\[
 V_{g,l}=N\mathfrak q\,
 A_{JH^3N^{2d-1}}(\alpha g^{2d}l-\gamma_{g,l}).
\]
The preceding comparison of $D(n,l)$ with its archimedean scale gives
$V_{D_0}\asymp_K V_{g,l}/N\mathfrak q_1$.
Ideal packing gives at most
$T^\epsilon(1+N/N\mathfrak q_1)$ choices for $n$ in each possible class.
The terms $N\mathfrak q_1>N$ contribute $O(T^\epsilon JHM)$ in total:
there are $O(T^\epsilon)$ possible $n$ for each $g,l$ and divisor class,
and their weight is $V_{D_0}^{-1/3}\le1$;
in each dyadic $G$ range the number of $g,l$ pairs is $O(T^\epsilon M)$.
For the rest, $(N\mathfrak q_1)^{1/3}/N\mathfrak q_1\le1$.
Thus the remaining sum is bounded by
\[
 T^\epsilon JHM\sum_gG^{-1}\sum_l V_{g,l}^{-1/3}.
\]
Here and below the summation over $g$ means one balanced generator per
ideal and finitely many fixed class choices; its count in a dyadic range
of $G$ is $O_K(GT^\epsilon)$.
We may further restrict to $V_{g,l}<(U/G)^3$, at a cost
$O(T^\epsilon JHM^2/U)$.

\emph{Third identification: remove the difference.}
For fixed $g$, approximate $\alpha g^{2d}$ with denominator radii
$(J_vH_v^3N_v^{2d})^{1/2}$, obtaining $\gamma_g$.
Comparison with $\gamma_{g,l}/l$ has cross-product norm
\[
 \ll_K\frac{V_{g,l}}{(JH^3)^{1/2}N^{d-1}}<1.
\]
The previous bound on $V_{g,l}$ and the hypotheses on $U$ ensure the last
inequality.  Hence $\gamma_{g,l}=l\gamma_g$ exactly.
Splitting according to the divisor ideal of the conductor that divides $l$,
and summing $l$ on dyadic archimedean boxes, gives
\[
 \sum_l V_{g,l}^{-1/3}\ll T^\epsilon N V_g^{-1/3},\qquad
 V_g=Q_{\gamma_g}A_{JH^3N^{2d}}(\alpha g^{2d}-\gamma_g).
\]
For an explicit verification, put $A=A_{JH^3N^{2d}}(\beta)$.
Then $A_{JH^3N^{2d-1}}(l\beta)\gg A|Nl|/N$.
If $l$ belongs to a divisor ideal $\mathfrak b$, the preceding ideal norm
count gives
$\sum_{0\ne l\in\mathfrak b\cap\mathcal B(N)}|Nl|^{-1/3}
\ll T^\epsilon N^{2/3}/N\mathfrak b$.
Combining this with the conductor factor yields divisor weights
$(N\mathfrak b)^{-2/3}$, whose sum is $O(T^\epsilon)$.
This proves the displayed linear-variable estimate.
We have therefore reduced to
\[
 T^\epsilon JHM^2\sum_gG^{-2}V_g^{-1/3}.
\]
Terms with $V_g\ge(U/G)^3$ are again acceptable, so assume the reverse
inequality.  In particular $G<U$.

\emph{Fourth identification: remove the common divisor.}
Choose one Dirichlet center $\gamma$ for $\alpha$ at scale $T$.
Comparison with $\gamma_g/g^{2d}$ has cross-product norm at most
\[
 \ll_K V_gG^{2d}T^{-1/2}
 \le C_KU^3G^{2d-3}T^{-1/2}
 \le C_KU^{2d}T^{-1/2}<1.
\]
Thus $\gamma_g=g^{2d}\gamma$.  Write $Q=Q_\gamma$ and $A=A_T(\alpha-\gamma)$.
The conductor of $g^{2d}\gamma$ is obtained by dividing the conductor of
$\gamma$ by its greatest common divisor with $(g)^{2d}$.
The chosen generators are balanced, so the archimedean factor in $V_g$
is comparable to $A$.
Finally
\[
 \sum_{\mathfrak g}(N\mathfrak g)^{-2}
 \left(\frac{Q}{N(\mathfrak q+(\mathfrak g)^{2d})}\right)^{-1/3}
 \ll_{K,d,\epsilon}T^\epsilon Q^{-1/d}.
\]
This is an elementary Euler-product calculation: at a prime with
$v_\mathfrak p(\mathfrak q)=e$, the term of valuation $a$ has size
$q^{-2a-(e-\min(e,2da))/3}$, and its largest possible size is
$O(q^{-e/d})$ because $d\ge3$.  Summing the geometric tails and absorbing
the bounded factors per prime into $T^\epsilon$ proves the claim.
Since $A^{-1/3}\le A^{-1/d}$,
\[
 \sum_gG^{-2}V_g^{-1/3}\ll T^\epsilon(QA)^{-1/d}
 \le T^\epsilon\mathcal H_T(\alpha)^{-1/d}.
\]
Combining the four reductions proves the proposition.
\end{proof}


\section{The sharper cubic terminal for fifth powers}
\label{sec:k5-terminal}

Write $M_i=X^{\phi_i}$, $H_i=XM_i^{-5}$, $H=H_1H_2$,
$M=M_1M_2$, $Q=X/M$, and $T=Q^5=X^3H$.
The centered second difference for fifth powers is a cubic in $z$.
We adapt the sharper dissection of \cite[Lemma~6.1]{VW1995}.
After one ordinary difference in $z$, its quadratic part and its two
coefficient parts separate as
\[
 K(\alpha;h,h_1,h_2)
 L_1(\alpha;h,h_1,h_2)L_2(\alpha;h,h_1,h_2),
\]
where, up to fixed nonzero constants in $K$,
\[
 K=\sum_z e_K(\alpha hh_1h_2 z^2+\xi z),\qquad
 L_i=\sum_{m_i}e_K(c\alpha hh_1h_2h_i^2m_i^{10}+b_i(m_i)).
\]
The ordinary variable $h$ has volume $X$.

The support issue deserves to be explicit.  Partition each extracted
principal divisor into a dyadic norm annulus and then into finitely many
tiny archimedean boxes.  Its generator is balanced, so on one piece all
$|m_{i,v}|$ are comparable to $M_i^{1/n}$, up to the harmless original
smoothness loss.  The actual constraints
$|h_{i,v}m_{i,v}^5|\ll\delta X^{1/n}$ are then enlarged to independent
small boxes for $h_i$.  Every enlarged vertex remains in a slightly larger
tiny box about the same nonzero center.  This enlargement is made in the
nonnegative orthogonality count before an absolute value is introduced.
Consequently it preserves the nonvanishing of the differenced polynomial
and its derivative and yields the Cartesian auxiliary sums used here.
Smooth completion can alternatively be applied to the original vertex
weights.  Fourier inversion gives a linear twist in $z$ and twists of
degree five in $m_i$; the coefficient of the latter depends on $h_i$ and the
completion frequencies but is independent of the linear sampling variable
$hh_{3-i}$.  The estimates below are uniform in precisely these twists.

\subsection{The higher coefficient moment}
The coefficient-sampling proof in Section~\ref{sec:coefficient-sampling-common} has the
following strengthening.  If $V=H_0^3JN$ and
\[
 Y\le c_K\min\{J,N,(JN/H_0^3)^{1/2}\},
\]
then
\begin{equation}\label{eq:k5-sampling-strong}
 \sum_{0\ne h\in\mathcal B(H_0)}\sum_{j\in\mathcal B(J)}
 \left|\sum_a c_h(a)e_K(\alpha h^3ja)\right|^2
 \ll V^\epsilon JH_0N\{Y^{-1}+\mathcal H_V(\alpha)^{-1/3}\}
                      \sup_h\sum_a|c_h(a)|^2.
\end{equation}
The coefficients may depend on $h$ because the elementary bilinear
large-sieve bound is applied separately for each $h$.
Here are the additional details proving the exponent $-1/3$.
The small-height centers are identified as $h^3\gamma$, where
$L=Q_\gamma A_V(\alpha-\gamma)\ll V^{1/2}$, by the conductor separation
lemma.  Put $\mathfrak d=\mathfrak q+(h)^3$ and
$\mathfrak b=\prod\mathfrak p^{\lceil v_\mathfrak p(\mathfrak d)/3\rceil}$.
For $A=A_V(\alpha-\gamma)$ the norm-count argument bounds the contribution
of this divisor class by
\[
 V^\epsilon\frac{N\mathfrak d}{Q_\gamma}
       \frac{H_0}{A^{1/3}N\mathfrak b}
 \le V^\epsilon H_0L^{-1/3},
\]
since $N\mathfrak b\ge(N\mathfrak d)^{1/3}$ and
$N\mathfrak d\le Q_\gamma$.
There are only $V^\epsilon$ divisor classes.  The large-height classes and
the $J+N$ terms of the bilinear bound contribute $JH_0N/Y$.
This proves \eqref{eq:k5-sampling-strong}.

Apply this with $c_h(a)$ the coefficients of a $t$th power of a decic smooth
sum of volume $M_1$.  Its squared coefficient norm is bounded by the decic
$2t$th moment, even in the presence of the lower-degree twists: by
orthogonality and the triangle inequality, its absolute coefficient count
is dominated by the untwisted decic count.  Thus if $\mu_t$ is a permissible
decic $2t$-moment exponent, then
\[
 E_t(\alpha;J,H_0,M_1)
 \ll X^\epsilon JH_0M_1^{\mu_t+10}
             \{Y^{-1}+\mathcal H_{JH_0^3M_1^{10}}(\alpha)^{-1/3}\}.
\]

\subsection{The terminal estimate}
Let $\theta=\phi_1$, $\phi=\phi_2$, and write $U=X^u$.  Assume
\begin{gather}
 1/15\le\theta\le1/5,\quad 0<\phi\le\theta,\quad
 \phi<5\theta-1/3,\label{eq:k5-parameter1}\\
 0<u<\phi,\quad u<-1-5\theta+15\phi,\quad
 4u+\theta+\phi<1,\quad 6u+5\theta+24\phi<5.
 \label{eq:k5-parameter2}
\end{gather}
The strict inequalities absorb the small powers arising from principal
blocks and from dyadic completion.

Let $t\ge3$, and let $\mu_t$ be the decic $2t$-moment exponent.  Put
\[
 Z=XU^{1-1/t}\left(X^{1/3}M_1^{2t-10-\mu_t}\right)^{1/t},
 \qquad \rho=\frac{\log_X Z}{4}.
\]
On the cubic minor arcs $\mathcal H_T(\alpha)>C_KX$, ordinary differencing
and H\"older give
\[
 |F_2|^2\ll X(HM)^2+HM\,D^{1/2}E_1^{1/(2t)}E_2^{1/2-1/(2t)}.
\]
Here $D$ is the sum of $|K|^2$, $E_1$ the sum of $|L_1|^{2t}$, and
$E_2$ the sum of $|L_2|^{2t/(t-1)}$, over $h,h_1,h_2$.
The quadratic linear-kernel calculation gives
\[
 D\ll X^{2+\epsilon}H.
\]
In detail, after expanding $|K|^2$, the zero difference contributes $X^2H$;
for the other terms factor $z_1^2-z_2^2=(z_1-z_2)(z_1+z_2)$, group the
nonzero coefficient $hh_1h_2(z_1-z_2)$ by ideal divisors, and apply the
positive linear-kernel estimate at scale $T$.  Its minor-arc bound is
$T/X=X^2H$.

Group the product $hh_2$ in $E_1$, at a cost $X^\epsilon$.
Equation~\eqref{eq:k5-sampling-strong}, with
$J=XH_2$, $H_0=H_1$, $N=M_1^{10}$ and $Y=X^{1/3}$, is applicable by
\eqref{eq:k5-parameter1}.  It gives
\[
 E_1\ll X^{2/3+\epsilon}H M_1^{\mu_t+10}.
\]
For $E_2$, first remove the power exceeding two using $|L_2|\ll M_2$.
Proposition~\ref{prop:k5-coeff2}, with $J=XH_1$, $H_0=H_2$, $M=M_2$ and
$d=5$, applies by \eqref{eq:k5-parameter2}.  Since $U<X^{1/5}$, it yields
\[
 E_2\ll X^{1+\epsilon}H M_2^{2t/(t-1)}U^{-1}.
\]
Consequently, provided $\rho\le1/2$,
\begin{equation}\label{eq:k5-sharp-minor}
 \sup_{\mathfrak m}|F_2|\ll X^{1+\epsilon}HM Z^{-1/4}.
\end{equation}

The already-proved cubic complete-sum and Poisson estimates give
\[
 \int_{\mathfrak M}|F_2|^4\ll X^\epsilon(XHM)^4Q^{-5}.
\]
They apply uniformly to the completed support and linear twists just
specified, including all complex places and all nonzero Poisson modes.
For any permissible quintic smooth moments $\lambda_r$ and
$\lambda_{4r/3}$, H\"older on the major arcs and
\eqref{eq:k5-sharp-minor} on the minor arcs therefore prove
\begin{equation}\label{eq:k5-enhanced-terminal}
 \int|F_2g_Q^{2r}|
 \ll X^{1+\epsilon}HM
 \left\{Z^{-1/4}Q^{\lambda_r}
       +Q^{-5/4+(3/4)\lambda_{4r/3}}\right\}.
\end{equation}
This is the needed field version of the sharper fifth-power terminal.
Its proof uses the decic moment, the coefficient lemma proved above, and
the existing cubic dissection; it does not import a rational moment
estimate as though it were already a theorem over $K$.


\section{A six-row certificate for the quintic joint moment}
\label{sec:k5-certificate}

The fundamental principal-block inequality in Section~\ref{sec:g6-transfer}
applies with exponent five: its proof uses only the polynomial degree,
nonvanishing of the derivative on the actual support, nonsingular Hensel
lifting, and the ideal class group.  Write $\lambda_s$ for the exponent of
the pure smooth $2s$th moment, and also retain every new bound as a joint
bound with two distinguished ordinary variables.
For the initial values we may use
\[
 \lambda_0=0,\quad\lambda_1=1,\quad\lambda_2=2,\quad
 \lambda_s=2s\quad(3\le s\le8).
\]
Only the ordinary fourth Hua moment is needed at initialization.

The low auxiliary moments required here are
\[
 \int |F_j|^2\ll X^{1+\epsilon}U_j,\qquad
 \int |F_j|^4\ll X^{2+\epsilon}U_j^3\qquad(j=1,2),
 \quad U_j=\prod_{i\le j}H_iM_i.
\]
Their proof is the low-residual-degree algebra already used for six and seven.
For $j=1$, the first difference is
$2h(5z^4+10h^2m^{10}z^2+h^4m^{20})$.
Completing the square in $z^2$ reduces its coincidence equation to a
product of two algebraic integers; the divisor bound then gives the
second moment.  For $j=2$, the residual cubic is divisible by $z$.
The nonzero phase-value divisor argument forces the same $h_1,h_2,z$ on
both sides, apart from $X^\epsilon$ multiplicity; the remaining coincidence
condition is
\[
 h_1^2(m_1^{10}-n_1^{10})+h_2^2(m_2^{10}-n_2^{10})=0.
\]
Cauchy and the fourth Hua moment for degree ten bound its count by
$X^\epsilon H_1H_2M_1M_2$.
One ordinary difference gives the fourth moments from these second moments.
For $F_2$ the appended residual degree is two; the new divisor is one,
so its coefficient cancels from the coincidence equation and gives no
extra $M$ loss.  The small-support and divisor arguments, including the
handling of zero phase values, are exactly as in the preceding field proof.

\subsection{The map being certified}
Interpolate $\lambda_s$ linearly between adjacent integer half-orders.
For an $A$ row with depth $j$, moment index $\ell\in\{1,2\}$, and
parameters $\phi_1,\ldots,\phi_j$, put
\[
 p=2^\ell,\quad b=1-p^{-1},\quad
 \Phi_i=\sum_{r\le i}\phi_r,\quad q_i=1-\Phi_i,
\]
and set
\[
 C_j=\frac{p-\ell+(p-1)j-4(p-1)\Phi_j}{p}
                +bq_j\lambda_{(S-j)/b}.
\]
For $i=j-1,\ldots,1$ define, with $r=S-i-1$ and $t=S-i+1$,
\[
 C_i=\frac{q_i\lambda_t+i-4\Phi_i+(2r-1)\phi_{i+1}}2
       +\frac12\max\{1+i-5\Phi_i+\Phi_{i+1}
                         +q_{i+1}\lambda_r,C_{i+1}\}.
\]
The output bound for the joint $2S$th moment is
\[
 T_S(\lambda)=
 \max\{1+\phi_1+q_1\lambda_{S-1},C_1\}+(2S-3)\phi_1.
\]
This is just the fundamental inequality, one backward H\"older step when
$j=2$, and the stated auxiliary second or fourth moment.

For the two enhanced rows take $r=6$, $t=2(S-1)-r$, and
\[
 \theta=\phi_1,\quad\phi=\phi_2,\quad q=1-\theta-\phi,
 \quad
 \rho=\frac14\left\{1+\frac{21u}{22}
                  +\frac{1/3+(34-34.227)\theta}{22}\right\}.
\]
Here the decic $44$th moment exponent $34.227$ is supplied by its separate
exact certificate and field auxiliary proof.
Put
\begin{align*}
 B&=\tfrac12\{(1-\theta)\lambda_t+1-4\theta+(2r-1)\phi\},\\
 D&=2-4\theta+\phi+q\lambda_r,\\
 C&=3-4(\theta+\phi)
        +\max\{-\rho+q\lambda_r,
                 q(-5/4+(3/4)\lambda_{4r/3})\}.
\end{align*}
The enhanced output is
\[
 T_S(\lambda)=\max\{1+\lambda_{S-1}
             +(2S-2-\lambda_{S-1})\theta,
                         (2S-3)\theta+B+\tfrac12\max(D,C)\}.
\]
Equation~\eqref{eq:k5-enhanced-terminal} proves this formula, including its
joint-moment interpretation.

\subsection{The exact supersolution}
All entries in the next table are rational numbers; the displayed integers
in the parameter columns are divided by $10^6$.
\[
\begin{array}{c|c|c|c|r|r|r|c}
 S&\text{row}&j&\ell&10^6\theta&10^6\phi&10^6u&b_S\\\hline
3&A&1&1&68130&&&3.1364\\
4&A&1&2&105600&&&4.4389\\
5&A&2&2&136590&72270&&5.9254\\
6&A&2&2&151340&113110&&7.5422\\
7&\text{enhanced}&2&-&163790&139365&139360&9.2726\\
8&\text{enhanced}&2&-&169960&138337&138330&11.0764
\end{array}
\]
The six strict differences $b_S-T_S(b)$ are, exactly,
\[
 \frac{62257}{2000000000},\quad
 \frac{2147}{25000000},\quad
 \frac{89351}{1000000000},\quad
 \frac{37509}{250000000},\quad
 \frac{128469}{500000000},\quad
 \frac{41387}{125000000}.
\]
The program included below verifies all parameter restrictions, interpolated
indices, and these rational inequalities without floating-point arithmetic.

The map is monotone and has sup-norm contraction constant at most $47/50$.
Indeed every affine branch has nonnegative variable coefficients with total
at most $q_1=1-\theta$, and each listed $q_1$ is smaller than $47/50$.
An enhanced terminal has coefficient mass at most $q$, and the backward
H\"older step replaces this by at most $(q_1+q)/2\le q_1$.
The major-arc branch has still smaller coefficient mass.
Clamping every output below by $S$ preserves these properties and the
supersolution inequalities.  Starting from the elementary initial bounds,
finite iteration therefore gives, for every $\epsilon>0$,
\[
 \int_{\mathbb T_K}|f_5(\alpha)|^2|g_5(\alpha)|^{14}\,d\alpha
          \ll X^{11.0764+\epsilon}.
\]
There is no appeal to an attained numerical fixed point: the number of
iterations is chosen finite after $\epsilon$ is prescribed, and the
smoothness parameter is then chosen sufficiently small for those finitely
many applications.

\section{The additional uniform transfer}
\label{sec:k5-global}
The preceding joint moment has excess
\[
 \Delta=\frac{191}{2500}<\frac1{13}.
\]
The ordinary minor-arc thirtieth moment and H\"older with weights
$1/14,13/14$ give
\[
 \int_{\mathfrak m_*}|f_5|^4|h_5|^{13}
 \ll X^{12-17/35000+\epsilon}.
\]
Here $h_5$ is the finite-product root sum in Section~\ref{sec:product-endgame}, so its
joint moments are dominated by the smooth moments up to an arbitrarily
small independent divisor loss.  Coefficientwise domination gives the
pure sixteenth moment first, and the loss-free absorption lemma gives
$U_{18}\ll X^{13}$ because
\[
 2-12\Delta=\frac{677}{625}>0.
\]
The broad major-arc annulus is controlled with ordinary moment order
$72/5>10$ and product moment order $18$, with H\"older weights $5/18,13/18$.
The finite-product core has a positive main term by the previously proved
Poisson and local-density argument.  Primitive local width is at most ten
for fifth powers, and hence is below seventeen at every finite completion.
Totally positive unit balancing supplies the same positivity and norm
convention as in Section~\ref{sec:conclusions}.  We conclude
\[
 \boxed{G_{\mathbb K}(5)\le17\quad\hbox{for every number field }K.}
\]
Only the number of summands is uniform in $K$; constants and sufficiently
large target thresholds may depend on the field.


\section{Primitive local solubility for the sixth and seventh powers}
\label{sec:lh6-local}

Let $A$ be the integer ring of a finite extension of $\mathbb Q_p$,
with uniformizer $\pi$, residue field $\mathbb F_q$, and absolute
ramification index $e=v_\pi(p)$. Write $J_k(A)$ for the additive group
generated by the $k$th powers. A representation is \emph{primitive}
if at least one of its roots is a unit. In particular, primitiveness
is required even for the target zero.

\begin{theorem}\label{thm:lh6-local}
Every element of $J_6(A)$ is a sum of at most twelve sixth powers,
primitively. Every element of $J_7(A)$ is a sum of at most fourteen
seventh powers, primitively. These bounds hold for every finite
extension of every $\mathbb Q_p$.
\end{theorem}

\subsection{Finite-field compression and exact carries}

We use the following finite-field fact: if $H=(\mathbb F_q^\times)^m$
and $E=\mathbb F_p(H)$, every element of $E$ is a sum of at most $m$
members of $H\cup\{0\}$. This is the finite-field Waring estimate
recorded in Joly, Proposition~7.16, and used in the proof of
Ramanujam, Lemma~1 \cite{Joly1970,Ramanujam1963}. The additive span of
$H$ is $E$: it is closed under multiplication, contains one, and is
a finite integral domain. One can also see the width bound directly
from Kneser's addition theorem. For $B=H\cup\{0\}$, the additive
stabilizer of any proper sumset $rB$ is $H$-invariant, hence an
$E$-subspace of $E$ and therefore zero. Consequently
$|rB|\ge r|H|+1$ until $rB=E$. Taking
$r=(|E|-1)/|H|\le m$ proves the assertion.

When $p\nmid m$, Hensel lifting now represents every unit of $J_m(A)$
by at most $m$ $m$th powers. Every target can be represented
primitively by $m+1$ terms, by first subtracting one. If $m\ge3$ is
odd, $m$ terms suffice: a zero residue has the primitive representation
$1+(-1)^m=0$, and a nonzero residue already has a unit coordinate.
We shall use primitive compression by three squares at $p=3$, and
by three cubes at $p=2$.

Put $k=pm$, where $(p,m)=(3,2)$ or $(2,3)$, and define the integral
carry polynomial
\[
 C_p(t_1,\ldots,t_N)
   =\frac{(\sum_i t_i)^p-\sum_i t_i^p}{p},\qquad
 \frac{\partial C_p}{\partial t_i}
   =(\sum_jt_j)^{p-1}-t_i^{p-1}.
\]
For $p=3$, start with any signed sixth-power expression for a target
$\nu_0$ and put $t_i=\epsilon_i x_i^2$. For $p=2$ use an unsigned
expression and put $t_i=x_i^3$. Such an unsigned expression always
exists: there is $r\in\mathbb Z_2^\times$ with $r^6=-7$, so
$-1=6\cdot1^6+r^6$. Multiplying this identity by a sixth power
replaces every negative summand by positive summands.
For example, Hensel lifting gives $z\in\mathbb Z_2$ with
$z^2+z+2=0$, whence $(1+2z)^2=-7$; an odd cube root of $1+2z$
then supplies $r$.

Compress the exact sum $S=\sum_i t_i$ as
$S=\sum_{j=1}^r y_j^m$, with $r\le3$ and a unit $y_j$, and put
$s_j=y_j^m$. Then
\begin{equation}\label{eq:lh6-carry}
 \nu_0-\sum_jy_j^6=pD,
 \qquad D=C_p(s)-C_p(t)\in A.
\end{equation}
Here $\bar D\in E$, with $E=\mathbb F_p((\mathbb F_q^\times)^m)$.
Choose $h_1,\ldots,h_l\in H$, $l\le m$, whose sum is $\bar D$,
and put
\begin{equation}\label{eq:lh6-constant}
 c=p\sum_{i=1}^l[h_i].
\end{equation}
Brackets here mean the ordinary Teichm\"uller lift in $A$.
Each $[h_i]$ is a sixth power: Frobenius permutes $H$, and
Teichm\"uller lifting is multiplicative. Thus $c$ uses at most six
sixth powers, and $D-c/p\in\pi A$.

We need one extra cancellation. If $H\subseteq\mathbb F_p^\times$
and $e\ge2$, then
\begin{equation}\label{eq:lh6-second-carry}
 D\equiv n\pmod{\pi^2A}
 \quad\text{for some }n\in\{0,1,\ldots,p-1\}.
\end{equation}
Indeed $A/\pi^2A\simeq\mathbb F_q[\varepsilon]/(\varepsilon^2)$
has characteristic $p$. Denote the coefficient of $\varepsilon$
by a prime. Every nonunit $t_i,s_j$ is divisible by $\pi^2$, and
every unit has $(p-1)$st power one in the residue field. The derivative
formula and the exact equality of the sums therefore give
\[
 D'=\sum_i t_i'-\sum_j s_j'=0.
\]
Its constant coefficient belongs to $\mathbb F_p$. This proves
\eqref{eq:lh6-second-carry}. The correction is then simply $c=pn$.

Finally, a sixth-power congruence with a unit coordinate lifts exactly
from modulus $9A$ at $p=3$ and from modulus $8A$ at $p=2$.
For a unit root $x$, use respectively $x(1+3z)$ and $x(1+4z)$:
\[
 \frac{(1+3z)^6-1}{9}\equiv2z\pmod\pi,
 \qquad
 \frac{(1+4z)^6-1}{8}\equiv3z\pmod\pi.
\]
Both normalized polynomials have a unit derivative. Hensel's lemma
absorbs any error in the stated modulus, preserving the unit root.

\subsection{Residue characteristic three}

Suppose first that $H\nsubseteq\mathbb F_3^\times$. Choose
Teichm\"uller squares $a,b$ with $a/b$ outside $\mathbb F_3$.
Reserve sixth-power roots whose underlying square values are $a,b$,
and apply the compression to the target with those two terms removed.
Put $A_0=a^3$, $B_0=b^3$, and $\lambda=-a/b$. The polynomial
\[
 F(Z)=\frac{A_0(1+Z)^3+B_0(1+\lambda Z)^3-A_0-B_0}{3}
\]
is integral: its cubic terms cancel because $A_0+B_0\lambda^3=0$.
Moreover $F'(0)=B_0(\lambda-\lambda^3)$ is a unit.
After the correction \eqref{eq:lh6-constant}, solve
$F(z)=D-c/3$ with $z\in\pi A$. Replace the underlying square
values by $a(1+z)$ and $b(1+\lambda z)$; they remain squares because
principal-unit square roots exist at $p=3$. This gives an exact
primitive representation with at most $3+2+6=11$ terms.

The remaining case has $q=3$ and $H=\{1\}$: indeed a square subgroup
of a finite field of characteristic three has order $(q-1)/2$, so it
cannot lie in $\mathbb F_3^\times$ when $q>3$.
If $e=1$, compression and \eqref{eq:lh6-constant} leave an error
in $3\pi A=9A$, and the preceding lifting argument uses at most nine
terms. Suppose henceforth that $e\ge2$.

Reserve three roots equal to one and compress the remaining target
by at most three squares. Let $b=y_1^2$ be a unit underlying square
among the compressed terms, so $b\equiv1\pmod\pi$. Equation
\eqref{eq:lh6-second-carry} and the correction $c=3n$, using at most
six ones, leave an error $3\pi^2R$.
Set $b_1=b$, $b_2=b_3=b_4=1$, $z_4=-z_1-z_2-z_3$, and
\[
 \mathcal F(z_1,z_2,z_3)
  =\frac{\sum_{i=1}^4(b_i+\pi z_i)^3-\sum_{i=1}^4b_i^3}
          {3\pi^2}.
\]
This is integral. Its linear part is $((b^2-1)/\pi)z_1$,
its quadratic part is $bz_1^2+z_2^2+z_3^2+z_4^2$, and its
cubic part is $\pi$ times an integral polynomial, since
$\sum_i z_i^3$ is divisible by three when $\sum_i z_i=0$.
Modulo $\pi$ its quadratic matrix is $I_3+\mathbf1\mathbf1^t$,
of determinant four, and hence is nondegenerate.

A nondegenerate ternary quadratic polynomial over a finite field of
odd characteristic represents every value nonsingularly. To see this,
complete the square and diagonalize. The two-square-set argument
represents every nonzero value by the first two coordinates; a zero
value has a nonzero representation by fixing the third coordinate to
one and representing its negative contribution by the first two.
Hensel lifting therefore solves $\mathcal F(z)=R$. Every
$b_i+\pi z_i$ remains a square. Modifying these four existing roots
absorbs the entire error and uses at most
\[
 3\text{ compressed terms}+3\text{ reserved terms}
       +6\text{ constant terms}=12.
\]

\subsection{Residue characteristic two}

Here $m=3$ and $H=(\mathbb F_q^\times)^3$.
If $H\ne\{1\}$, choose a Teichm\"uller cube $a$ whose residue is
different from one, and choose $\alpha$ with $\alpha^3=a$.
Reserve the roots $(\alpha,1)$ before compression. Replacing them by
\[
 \alpha(1+z)^{1/3},\qquad(1-az)^{1/3},\qquad z\in\pi A,
\]
changes their sixth-power sum by $2F_a(z)$, where
\begin{equation}\label{eq:lh6-dyadic-pair}
 F_a(Z)=(a^2-a)Z+a^2Z^2.
\end{equation}
Its derivative at zero is a unit. The correction
\eqref{eq:lh6-constant} leaves $D-c/2\in\pi A$, which this pair
absorbs by Hensel lifting. The number of terms is at most
$3+2+6=11$.

Suppose $H=\{1\}$ and $e\ge2$. Reserve four roots, in pairs
$(1,1)$ and $(1+\pi,1)$. For the second pair put $a=(1+\pi)^3$.
Then $a\equiv1+\pi\pmod{\pi^2}$ and $a^2\equiv1\pmod{\pi^2}$,
so
\[
 G(Z)=\frac{F_a(\pi Z)}{\pi^2}
     =\frac{a^2-a}{\pi}Z+a^2Z^2\equiv Z+Z^2\pmod\pi.
\]
By \eqref{eq:lh6-second-carry}, $c=2n$ uses at most two ones
and leaves an error $2\pi^2R$. Choose $t_0$ lifting the square root
of $\bar R$ and solve $G(t)=R-t_0^2$ with $t\in\pi A$.
The first pair, changed using $z=\pi t_0$ in
\eqref{eq:lh6-dyadic-pair} with $a=1$, contributes
$2\pi^2t_0^2$. The second pair, changed using $z=\pi t$,
contributes $2\pi^2G(t)$. Their sum is exactly the required error.
All four roots stay units, and at most $3+4+2=9$ terms are used.

If $H=\{1\}$ and $e=1$, the unit cubes are exactly $1+2A$:
the Teichm\"uller factor is killed, and cubing is bijective on
principal units. Thus the unit sixth-power classes modulo $8A$ are
\[
 1+4V,\qquad V=\{z+z^2:z\in\mathbb F_q\},
\]
while nonunit sixth powers vanish. Their generated image is
$(\mathbb Z/8\mathbb Z)\cdot1+4V$. Every class $a+4v$ is a
sum of between one and eight unit power classes: choose its count
$r\equiv a\pmod8$, use one class $1+4v$, and use $r-1$ ones.
Lifting from $8A$ gives the primitive bound eight.

At every other residue characteristic, finite-field compression of
$\nu-1$ and appending $1^6$, followed by ordinary Hensel lifting,
gives at most seven terms. This completes the sixth-power assertion.

\subsection{The seventh-power assertion and its primary source}

Joly's Theorem~7.34 states that an odd prime exponent $d$ has ordinary
width at most $2d-1$ in every henselian local ring with finite residue
field \cite[Theorem~7.34, pp.~105--106]{Joly1970}. Apply this with
$d=7$ to $\nu-1\in J_7(A)$ and append $1^7$. The resulting fourteen
terms are primitive, including when $\nu=0$. Since the exponent is
odd, replacing a negative root by its negative converts signed
notation to ordinary sums without increasing the count.

For precise attribution, the source is J.-R.~Joly,
\emph{Sommes de puissances $d$-i\`emes dans un anneau commutatif},
Acta Arithmetica \textbf{17} (1970), 37--114,
doi:10.4064/aa-17-1-37-114. Theorem~7.32 on printed page105 gives
the $p$-adic-ring case; Lemma~7.33 and Theorem~7.34 on pages105--106
give the stated henselian formulation. Proposition~7.16, printed
page95, is the finite-field input used above. Ramanujam's relevant
Lemma~1 is in \emph{Sums of $m$-th powers in $p$-adic rings},
Mathematika \textbf{10} (1963), 137--146.


\section{Primitive local bounds at exponents five through nine}
\label{sec:higher-local-bounds}

For a local integer ring $A$, write $J_k(A)$ for the additive group
generated by its $k$th powers, and $h_k(A)$ for the least number of
summands with coefficient $+1$ representing every member of $J_k(A)$ with at least
one unit root. The zero target is included. The following local bounds
are sufficient for all five proposed global counts:
\[
 \begin{array}{c|ccccc}
 k&5&6&7&8&9\\\hline
 \text{proposed global variable count}&17&24&31&38&46\\
 \text{primitive local bound}&10&12&14&35&30.
 \end{array}
\]
The eighth-power entry is the new local result below. Its proof uses
two small exact finite algebra certificates, explicitly identified.
The other four entries do not require new computations.

For $k=5,7$, Joly's Theorem~7.34 in \emph{Sommes de puissances
$d$-i\`emes dans un anneau commutatif}, Acta Arithmetica~17 (1970),
37--114, pp.~105--106, gives ordinary width $2k-1$ for an odd prime
exponent over a henselian ring with finite residue field. Apply that
theorem to $\nu-1$ and append $1^k$ to obtain the displayed primitive
bound. Signed and positive sums agree for these odd exponents.
The direct carry proof in Section~\ref{sec:lh6-local} gives
$h_6(A)\le12$ at both bad residue characteristics; good primes require
at most seven summands. We now give the coefficient-module and lifting
arguments needed for the eighth- and ninth-power entries.

\subsection{Witt coefficients with a positive summand count}
\label{sec:higher-witt-coefficients}

Let $p$ be a prime, let $q$ be a power of $p$, and put
$R=\mathbb F_q[[T]]$ and $B=W(\mathbb F_q)$, the unramified complete
discrete valuation ring with residue field $\mathbb F_q$.
We use $p$-typical Witt vectors $W_{a+1}(R)$, with coordinates
$(z_0,\ldots,z_a)$, Teichm\"uller vectors $[x]=(x,0,\ldots,0)$,
and ghost polynomials
\[
 w_j(z)=\sum_{i=0}^j p^i z_i^{p^{j-i}}.
\]
Their addition and multiplication are the universal integral
polynomials making the ghost maps homomorphisms. In particular
$[x][y]=[xy]$. All constant Witt vectors form the coefficient ring
$W_{a+1}(\mathbb F_q)=B/p^{a+1}B$.

Choose positive integers $N_0,\ldots,N_a$ satisfying
$p^{j-i}N_i\ge N_j$ for $i<j$, and put
\[
 I_N=\{z:z_j\in T^{N_j}R\},\qquad
 M_N=W_{a+1}(R)/I_N.
\]
Weighted homogeneity of the Witt operations shows that $I_N$ is an
ideal: a monomial of Witt weight $p^j$ in its coordinates has
$T$-order at least $N_j$. Multiplication is separately homogeneous
in each input, giving the same assertion when only one input is in
$I_N$. Triangularity in the coordinates gives
$|M_N|=q^{\sum_jN_j}$: choose the coordinates successively modulo
$T^{N_j}$, using the coefficient of $z_j$ equal to one in the
$j$th addition polynomial.

\begin{lemma}\label{lem:higher-witt-coefficients}
Let $G$ be the actual additive span of the vectors $[x]$, $x\in R$.
Every class of $G$ modulo $I_N$ is a positive sum of at most
\begin{equation}\label{eq:higher-witt-cost}
 C_p(N)=L_0+(p-1)\sum_{s=1}^a p^{s-1}L_s,
 \qquad
 L_s=\sum_{j=s}^a\left\lceil\frac{N_j}{p^s}\right\rceil
\end{equation}
actual Teichm\"uller vectors. The bound is independent of $q$.
\end{lemma}
\begin{proof}
Write $\ell(M)=\log_q|M|$ for finite $B$-modules. In characteristic
$p$, multiplication by $p^s$ has the coordinate description
\[
 p^sz=(\underbrace{0,\ldots,0}_{s},
             z_0^{p^s},\ldots,z_{a-s}^{p^s}).
\]
Thus its image modulo $I_N$ has cardinality $q^{L_s}$.
One can count the kernel after any sufficiently large common
$T$-truncation: its coordinate conditions are exactly
$z_i\in T^{\lceil N_{i+s}/p^s\rceil}R$ for $i\le a-s$.
Consequently $\ell(p^sM_N)=L_s$.

The image $M$ of $G$ in $M_N$ is a $B$-submodule, since multiplying
an actual generator by a constant Teichm\"uller vector preserves its
form. Adjoin actual generators $x_i$ successively to obtain a
filtration $0=M_0\subsetneq\cdots\subsetneq M_r=M$ with cyclic
quotients $M_i/M_{i-1}\simeq B/p^{\lambda_i}B$.
A coefficient modulo $p^\lambda$ has a Teichm\"uller digit expansion;
regarding its digit $p^j[c]$ as $p^j$ copies of $[c]$ uses at most
$Q_\lambda=1+p+\cdots+p^{\lambda-1}$ positive scalars. Working
backwards through the filtration therefore represents a target by
at most $\sum_iQ_{\lambda_i}$ actual generators.

For every $s$, the natural surjections
$p^sM_i/p^sM_{i-1}\twoheadrightarrow p^s(M_i/M_{i-1})$ give
\[
 \ell(p^sM)\ge\sum_i\max(\lambda_i-s,0).
\]
The identity
\[
 Q_\lambda=\lambda+
   (p-1)\sum_{s\ge1}p^{s-1}\max(\lambda-s,0)
\]
then bounds the positive cost by
$\ell(M)+(p-1)\sum_{s\ge1}p^{s-1}\ell(p^sM)$.
Each length is at most the corresponding length for $M_N$, since
$p^sM\subseteq p^sM_N$. This proves
\eqref{eq:higher-witt-cost}. In particular, the argument never
replaces actual power generators by arbitrary invariant-factor
coordinates.
\end{proof}

Here is the link to local integral powers. For a local integer ring
$A$ of residue field $\mathbb F_q$, choose its unramified coefficient
ring $B\subseteq A$ and the surjection $D=B[[T]]\to A$ sending
$T$ to $\pi$. On $W_{a+1}(R)$, choosing coordinate lifts to $D$
defines an additive map
\begin{equation}\label{eq:higher-last-ghost}
 \Phi:W_{a+1}(R)\longrightarrow D/p^{a+1}D,
 \qquad \Phi(z)=\sum_{i=0}^a p^i\widetilde z_i^{p^{a-i}}.
\end{equation}
It is independent of the lifts: changing a lift by a multiple of
$p$ changes its $p^{a-i}$th power by a multiple of $p^{a-i+1}$.
Additivity follows by lifting the universal Witt addition and using
the ghost identity. In particular $\Phi([x])=\widetilde x^{p^a}$.
Lift an eligible target by any finite signed power expression to
$D$; its associated signed Teichm\"uller sum reduces to an element
of $G$.

Suppose a centered chart covers $I_N$ in the form
\[
 w=\sum_{i=1}^r[x_i]-\sum_{i=1}^r[b_i],
\]
retaining a unit among the $x_i$. Subtract the fixed baseline
$\sum[b_i]$ from the target and use Lemma~\ref{lem:higher-witt-coefficients}
modulo $I_N$. The remaining ideal element is supplied by the chart,
and the baseline cancels. The result is a primitive representation
by at most $C_p(N)+r$ powers modulo $p^{a+1}D$, and hence modulo
$p^{a+1}A$. Fixed nonzero constant residues in the chart are
preserved by choosing lifts of the roots with those constants.

\subsection{The ninth-power bound}
\label{sec:higher-ninth-local}

Only residue characteristic three requires attention. Use $p=3$,
$a=2$, and coordinate orders $N=(3,5,5)$. In $W_3(R)$ take centers
$(b_0,b_1,b_2)=(0,1,T)$ and perturbations
\[
 (\delta_0,\delta_1,\delta_2)
       =(T^3u_0,T^5u_1,T^3u_2).
\]
For $H=\sum[b_i+\delta_i]-\sum[b_i]$, the universal ghost
recurrence gives $H_j\in T^{N_j}R$. After dividing these coordinates
by $T^3,T^5,T^5$, the reduction modulo $T$ is
\[
 (u_0+u_2,\ u_1+u_2,\ u_1).
\]
Indeed the linear term of $H_1$ is $\sum b_i^2\delta_i$, and that
of $H_2$ is $\sum b_i^8\delta_i$; the other terms have strictly
higher $T$-orders. The displayed linear map has determinant $-1$,
so formal Hensel inversion covers every element of $I_N$.
The perturbed second root is $1+T^5u_1$, hence a unit.

The image lengths are $L_0,L_1,L_2=13,4,1$, so the positive
coefficient cost is
\[
 C_3(3,5,5)=13+2(4+3\cdot1)=27.
\]
The three baseline terms give thirty ninth powers modulo $27A$.
A unit root absorbs the remaining error: the polynomial
\[
 \frac{(1+3z)^9-1}{27}\equiv z\pmod\pi
\]
is integral and has a unit derivative. Hensel's lemma gives the
exact primitive representation with the same number of summands.
At all other residue characteristics, finite-field compression of
$\nu-1$, appending one, and ordinary Hensel lifting use at most ten
terms. Thus $h_9(A)\le30$ without any restriction on residue degree
or ramification.

\subsection{A primitive eighth-power bound of thirty-five}

\begin{theorem}\label{thm:local-eighth-thirtyfive}
For the integer ring $A$ of every finite extension of $\mathbb Q_p$,
\[
 h_8(A)\le35.
\]
More precisely, at dyadic places the following bounds hold:
\[
 \begin{array}{c|c}
 \text{residue degree and absolute ramification}&\text{bound}\\\hline
 f\ge3&19\\
 f=2&30\\
 f=1,\ e=1&32\\
 f=1,\ e=2&32\\
 f=1,\ e=3,4&35\\
 f=1,\ e=5&22\\
 f=1,\ e\ge6&33.
 \end{array}
\]
The entries $e=5$ and $e\ge6$ use the exact finite certificates
described below. Their sizes are respectively $8192$ residue classes
and $64$ polynomial identities plus eight carry relations.
\end{theorem}

Write $\pi$ for a uniformizer. A primitive eighth-power congruence
modulo $16\pi A$ lifts without increasing its number of summands.
Indeed, changing a unit root $x$ to $x(1+2\pi z)$ changes its power by
$16\pi x^8F(z)$, where $F\in A[z]$ and $F(z)\equiv z\pmod\pi$.
All higher binomial terms are divisible by an additional $\pi$ after
division by $16\pi$. Simple-root Hensel lifting applies.
At an odd prime the usual derivative argument gives $h_8(A)\le9$.

\subsection{Residue degrees at least two}

Work first in $W_4(\mathbb F_q[[T]])$. For a centered sum of four
Teichm\"uller vectors with centers $b_i$, its Jacobian has entries
$b_i^{2^j-1}$, $0\le j,i\le3$, with the first row equal to one.
This follows directly by differentiating the universal ghost identities.

If $f\ge3$, take centers $0,c_1,c_2,c_3$, where the three nonzero
constants are linearly independent over $\mathbb F_2$. The zero column
is $(1,0,0,0)^t$; the other minor is an invertible Moore matrix, after
multiplication of its columns by the $c_i$. The centered sum therefore
covers the ideal whose four coordinates belong to $TR$: substitute
perturbations $Tz$, divide the four coordinates by $T$, and apply
formal Hensel inversion to the resulting map with invertible
linear reduction.
By Lemma~\ref{lem:higher-witt-coefficients}, the ambient coefficient cost is
$4+3+2\cdot2+4\cdot1=15$, so the four-variable baseline gives
nineteen powers modulo sixteen.

If $f=2$, choose $\omega^2+\omega+1=0$ and centers $(0,1,\omega,T)$.
The determinant is
\[
 (1+\omega)(T^7+T)=T\cdot\text{a unit}.
\]
Multiplication by the adjugate and the substitution
$\eta=T^2\operatorname{adj}(J)z$ normalize the centered map to
$uz+TP(z)$ after division by $T^3$, with $u$ a unit. Thus the
isotropic ideal of coordinate order three is covered. Its $2$-power
image lengths are $12,6,2,1$, and its coefficient cost is
\[
 12+6+2\cdot2+4\cdot1=26.
\]
Adding the four-variable baseline gives thirty powers modulo sixteen,
by \eqref{eq:higher-last-ghost} and the baseline cancellation above.

To justify the last integral lift, lift the represented roots from
$R=\mathbb F_q[[T]]$ to $D=B[[T]]$. The residual target is in $16D$
by \eqref{eq:higher-last-ghost}. Both baselines retain two units
with distinct eighth-power residues. The final digit modulo
thirty-two can therefore be corrected using these same roots, with
no extra summands. Explicitly,
\[
 (1+2z)^8\equiv1+16(z^2+z)\pmod{32}.
\]
If their lifted roots are $y_1,y_2$, the attainable correction divided
by sixteen is
\[
 \bar y_1^8(z_1^2+z_1)+\bar y_2^8(z_2^2+z_2)\quad\text{in }R.
\]
The Artin--Schreier map $z\mapsto z^2+z$ is bijective on $TR$ and
has image the absolute-trace-zero hyperplane on $\mathbb F_q$.
Multiplication by either unit $\bar y_i^8$ preserves $TR$.
Its two constant images are distinct trace hyperplanes because the
nonzero constants $\bar y_1(0)^8$ and $\bar y_2(0)^8$ are distinct;
the nondegenerate trace pairing identifies their annihilator lines
with their inverse scalars. The sum of these images is consequently
all of $R$. Choose $z_1,z_2$, lift them to $D$, and replace the two
roots by $y_i(1+2z_i)$. This matches the target modulo $32D$ while
preserving both unit roots. After specialization, $32A\subseteq16\pi A$,
so the primitive lifting argument already proved completes the exact
representations with nineteen and thirty terms.

\subsection{A seven-variable weighted chart}

Let $D=\mathbb Z_2[[T]]$ and set
\[
 I=\{(z_0,z_1,z_2,z_3)\in W_4(D):
 z_0\in T^2D,\ z_1\in T^4D,\ z_2,z_3\in T^6D\}.
\]
The inequalities $2^{j-i}N_i\ge N_j$ for
$N=(2,4,6,6)$ show, by weighted homogeneity of the universal Witt
operations, that $I$ is an ideal. The same definition applies after
reduction modulo two.

\begin{lemma}\label{lem:eighth-seven-chart}
Every $w\in I$ has a representation
\[
 w=\sum_{i=0}^6[x_i]-\sum_{i=0}^6[b_i],
 \qquad (b_0,\ldots,b_6)=(0,1,T,T^2,0,0,0),
\]
where $x_1\equiv1\pmod{T}$.
\end{lemma}
\begin{proof}
Use perturbations
\[
 \delta_1=T^6u_1,\qquad\delta_2=T^3u_2,\qquad
 \delta_i=T^2u_i\quad(i=0,3,4,5,6).
\]
Write $H=\sum[b_i+\delta_i]-\sum[b_i]$ and put
$v=(u_0,u_3,u_4,u_5,u_6)$. The universal ghost recurrence shows that
$H_j$ is divisible by $T^{N_j}$. Modulo $(2,T)$, the normalized map is
\[
 \left(\sigma_1(v),\ u_2+u_3+\sigma_2(v),\ u_1+u_2,\ u_1\right).
\]
For the second component, one may use the exact formula
$H_1=\sum b_i\delta_i-\sum_{i<j}\delta_i\delta_j$.
In the last two components, all nonlinear ghost terms have orders
strictly greater than six, leaving the displayed linear terms.

The last two equations fix $u_1,u_2$. For the first two, it is enough
to solve $\sigma_1(v)=s$ and $\sigma_2(v)+u_3=c$ nonsingularly
over $\mathbb F_2$. For $(s,c)=(0,0),(0,1),(1,0),(1,1)$, choose
respectively zero, two, one, and three ones among the four coordinates
other than $u_3$, keeping $u_3=0$. The two Jacobian rows are independent:
their difference from $1$ is the vector $v+e_3$, which is neither zero
nor the all-one vector for these choices. The full Jacobian has rank
four. First lift the chosen solution through the $2$-adic coefficient
ring, then through the $T$-adic parameter, freezing the unused variables.
This solves every normalized target in $D^4$ and proves the lemma.
The root $x_1=1+T^6u_1$ remains a unit.
\end{proof}

\subsection{The finite coefficient certificate}

Put $\bar D=\mathbb F_2[[T]]$ and let $G$ be the actual additive
Teichm\"uller span in $W_4(\bar D)/\bar I$.
The following eight actual generators suffice:
\[
 [1],\ [1+T^5],\ [1+T^4],\ [1+T^3],\ [1+T^2],\
 [1+T],\ [T],\ [T+T^2].
\]
In this order, their successive cyclic quotient orders are
\[
 16,2,2,2,2,4,4,2.
\]
Consequently every element of $G$ is a positive sum of at most
\[
 15+1+1+1+1+3+3+1=26
\]
of the displayed generators.

Here is the precise finite verification used for that assertion.
The combined ghost map sends a vector to $2^{3-j}w_j$, truncated
at degree $N_j$, modulo sixteen for each $j=0,1,2,3$.
It is additive and has kernel exactly $\bar I$, by induction on the
coordinates. Every root can be reduced modulo $T^6$, leaving exactly
sixty-four binary polynomials. The standalone program in Appendix~\ref{app:eighth-ghost} checks their images
against the $8192$ normal forms with the coefficient ranges just
displayed. It also checks all eight carry relations, which make this
normal-form set an additive subgroup. If the generators above are
$g_0,\ldots,g_7$, these relations are
\[
 \begin{split}
 16g_0&=0,\quad 2g_1=2g_0,\quad2g_2=2g_0,\quad2g_3=2g_0,\\
 2g_4&=g_0+g_2,\quad4g_5=3g_0+g_2,\\
 4g_6&=12g_0+g_2+g_4+2g_5,\quad2g_7=2g_6.
 \end{split}
\]
The program also emits a coefficient witness for every binary root. All checks use exact integer arithmetic modulo sixteen.

\subsection{High ramification}

Suppose $f=1$ and $e\ge6$. Lift an eligible target through a finite
signed-power expression to $D$ and then to a Witt vector $g\in W_4(D)$
whose last ghost is that expression. Subtract the seven-term baseline,
reduce modulo two, and apply the preceding certificate modulo $\bar I$.
This uses at most twenty-six positive Teichm\"uller vectors. After
lifting them to $D$, the remaining Witt vector $\eta$ reduces into
$\bar I$. Choose $w\in I$ lifting this reduction. The last ghost of
$\eta-w$ is divisible by sixteen.

After specialization $T=\pi$, this remaining error belongs to
$16A\subseteq8\pi^6A$, since $e\ge6$. It can therefore be matched
exactly by the last ghost of a vector $V^3(T^6z)\in I$, after lifting
its coefficient from $A$ to $D$. Add that vector to $w$ and apply
Lemma~\ref{lem:eighth-seven-chart}. The fixed baseline cancels, leaving
at most $26+7=33$ positive eighth powers, one with a unit root.
No enlarged power-generated order is used in this reduction.

\subsection{Small ramification}

The rational ring has $h_8(\mathbb Z_2)=32$: unit eighth powers are
one modulo thirty-two, nonunit eighth powers vanish there, and a
primitive zero needs thirty-two unit terms. Any residue class is
represented by between one and thirty-two ones, followed by lifting.

For $e=2,3,4$, work modulo $16\pi A=\pi^{4e+1}A$.
The eighth power of every unit is $1+d$ with
\[
 v_\pi(d)\ge7\quad(e=2,3),\qquad
 v_\pi(d)\ge9\quad(e=4).
\]
For $e=2,4$, the two lowest terms in the binomial expansion cancel
modulo $\pi$, since the residue field is $\mathbb F_2$; direct use
of $(1+\pi y)^8$ gives these bounds. The product of any two such
deviations vanishes at the working precision. Thus the deviations
form an additive subgroup $D_0$: multiplication of unit eighth powers
corresponds to addition of deviations.

When $e=2$, every nonunit eighth power is zero or $\pi^8$, and
$\pi^8\equiv16$ at this precision. The generated image is consequently
$\mathbb Z\cdot1+D_0$. It is represented by at most thirty-two unit
eighth powers, using one unit power for the prescribed deviation.

When $e=3$, choose $\pi^3=2$, by extracting the cube root of the
unit in any initial uniformizer relation. The nonunit images are
zero and $\pi^8$, and $4\pi^8=0$ at precision $\pi^{13}$.
Thus at most three copies of $\pi^8$, together with at most
thirty-two unit eighth powers, represent every eligible class.

When $e=4$, the nonunit images are zero, $\pi^8$, and $\pi^{16}$.
Writing $\pi^4=2u$ shows
\[
 \pi^{16}\equiv16,\qquad4\pi^8\equiv16\pmod{\pi^{17}}.
\]
The same bound of thirty-five follows. All these representations
have a unit root and lift.

Finally, if $e=5$, extracting a fifth root of a principal unit gives
$A=\mathbb Z_2[\pi]$ with $\pi^5=2$. Eighth powers modulo $16\pi$
depend only on roots modulo $2\pi$. The latter quotient has sixty-four
classes, represented by
\[
 a_0+a_1\pi+\cdots+a_4\pi^4,
 \qquad0\le a_0<4,\quad a_i\in\{0,1\}\ (i>0).
\]
The standalone program in Appendix~\ref{app:eighth-tame} checks these classes in the exact ring
\[
 (\mathbb Z/32)\oplus
 (\mathbb Z/16)\pi\oplus\cdots\oplus(\mathbb Z/16)\pi^4,
 \qquad\pi^5=2.
\]
There are twenty distinct eighth-power images. Their generated group
has $8192$ elements and ordinary additive diameter twenty-one.
The breadth-first calculation verifies every parent representation
and closure under every generator. Applying it to the target minus
one and appending $1^8$ gives a primitive bound twenty-two.
This exhausts the ramification cases and proves
Theorem~\ref{thm:local-eighth-thirtyfive}.


% Root-owned continuation lemma. All scale parameters are volumes.
\section{Restricted-height smooth moments over a number field}
\label{sec:restricted-field}

The following field version of Br\"udern--Wooley's restricted-height
mean theorem is needed by the eighth- and ninth-power routes.  The use of
ideals in its proof is particularly useful: an extracted ideal factor
need not be principal.  The quotient variables instead lie in one of
finitely many fixed integral ideals.

Write $V=K\otimes\mathbb R$, $n=[K:\mathbb Q]$, and $d_v=[K_v:\mathbb R]$.
For $T\ge1$ put
\[
 A_T(\beta)=\prod_{v\mid\infty}
                 \max(1,T^{1/n}|\beta_v|)^{d_v}.
\]
For a fixed integral ideal $C$, use the coefficient torus
$\mathbb T_C=V/(C^{-k}\mathfrak D^{-1})$ with probability Haar measure.
The conductor of $\gamma\in K/(C^{-k}\mathfrak D^{-1})$ is
\[
 \mathfrak a_C(\gamma)
       =\{a\in\mathcal O_K:a\gamma C^k\subseteq\mathfrak D^{-1}\}.
\]
Let $\mathfrak M_C(Z;Y)$ be the union of the product arcs
\[
 (N\mathfrak a_C(\gamma)) A_{Y^k}(\alpha-\gamma)\le Z.
\]
All lifts of a center are included.  A subscript $\mathfrak a$ on this
set restricts to centers of exact conductor $\mathfrak a$.
For $C=\mathcal O_K$, suppress the subscript $C$.

Let $g(\alpha;X,R)$ be a sum over the nonzero algebraic integers in a fixed
bounded box of scale $X^{1/n}$ whose prime ideal factors have norm at most
$R$.  Nonnegative bounded weights and smaller fixed boxes are allowed.

\begin{proposition}[Field restricted-height mean]
\label{prop:restricted-field}
Let $s\ge2$ be even.  Suppose $\lambda=s-k+\Delta$ is a permissible
full smooth-moment exponent, with
$\lambda\ge\max(s/2,s-k)$ and $\Delta\ge0$.  Then for every
$\epsilon>0$ a sufficiently small smoothness exponent gives
\[
 \int_{\mathfrak M(Z;X)}|g(\alpha;X,R)|^s\,d\alpha
       \ll_{K,k,s,\epsilon} X^{s-k+\epsilon}Z^{2\Delta/k}
       \qquad(1\le Z\le X^{k/2}).
\]
The same conclusion holds at intermediate real orders when their full
exponents are obtained by H\"older interpolation between such even orders.
\end{proposition}

\subsection{Fixed root lattices and multiplication of arcs}

Choose integral ideal representatives $C_1,\ldots,C_h$ for the ideal
classes.  For every nonzero integral ideal $\mathfrak b$ in the class of
$C_i$, choose $b=b(\mathfrak b)\in K^\times$ with
\[
 \mathfrak b=b C_i,\qquad
 |b_v|\asymp_K(N\mathfrak b)^{1/n}\quad(v\mid\infty).
\]
Dirichlet's unit theorem gives the second condition after choosing one
representative in a fundamental domain for the unit logarithms.  This
choice is made once for each ideal.  Multiplication by $b$ maps $C_i$
into $\mathcal O_K$.

If $(\mathfrak a,\mathfrak b)=1$, multiplication by $b^k$ maps centers
of exact conductor $\mathfrak a$ on $\mathbb T_{\mathcal O_K}$ bijectively
to those of exact conductor $\mathfrak a$ on $\mathbb T_{C_i}$.
Indeed the conductor calculation is local, and
$b^kC_i^k=\mathfrak b^k$ is a unit ideal at every prime dividing
$\mathfrak a$.  Equivalently, multiplication by $b^k$ is an isomorphism
on these finite conductor-character groups.
Its archimedean Jacobian, with the indicated probability measures, is
$(N\mathfrak b)^k$.  To check the normalization, ordinary Lebesgue measure
has Jacobian $|Nb|^k=(N\mathfrak b/NC_i)^k$, while the volume of
$\mathbb T_{C_i}$ is $(NC_i)^{-k}$ times that of $\mathbb T_{\mathcal O_K}$.

Consequently, whenever the image arcs are disjoint, for any nonnegative
integrand $F$ one has
\begin{equation}\label{eq:field-arc-scaling}
 \int_{\mathfrak M_{\mathfrak a}(Z;X)}F(b^k\alpha)\,d\alpha
 \ll_K (N\mathfrak b)^{-k}
 \int_{\mathfrak M_{C_i,\mathfrak a}(c_1Z;X/N\mathfrak b)}
                         F(\theta)\,d\theta .
\end{equation}
Here $c_1$ is a fixed constant, independent of the ideal and its class.
Balancing $b$ makes the archimedean scaling comparable at every place.
The requisite disjointness holds if
$c_1Z\le c_2Y^{k/2}$, with a sufficiently small $c_2=c_2(K,k)$,
where $Y=X/N\mathfrak b$.
For distinct lifted centers the nonzero norm of their difference is
$\gg_K(QQ')^{-1}$, while overlapping product arcs would give an upper
bound $\ll_KY^{-k}Z^2/(QQ')$.  This also excludes overlap of distinct
lifts of the same center.  There are only finitely many root lattices
in use, so one constant suffices for all of them.

\subsection{The ideal-factor decomposition}

Fix a total ordering of the prime ideals which refines their norm ordering.
Given a conductor ideal $\mathfrak a$, factor
\[
 (x)=\mathfrak v\mathfrak u,
 \qquad (\mathfrak v,\mathfrak a)=1,
 \qquad \mathfrak u\text{ supported on primes dividing }\mathfrak a.
\]
Call a root bad when $N\mathfrak v\le M$.  The supported-ideal counting
lemma from the fundamental differencing argument, followed by bounded-height
unit counting, gives at most $X^\epsilon M$ such roots, uniformly for
$N\mathfrak a\le X^{k/2}$.  This remains true when the variable lies in
a smaller box or has a nonnegative weight bounded by a fixed constant.

For a good root, multiply the prime factors of $\mathfrak v$ in decreasing
order until their product first has norm exceeding $M$.  Denote the
resulting ideal by $\mathfrak m$ and its last prime factor by $\mathfrak p$.
This decomposition is unique, and
\[
 M<N\mathfrak m\le MR,\qquad
 (\mathfrak m,\mathfrak a)=1.
\]
Every prime of the remaining ideal $\mathfrak w=\mathfrak v/\mathfrak m$
is at most $\mathfrak p$ in the fixed ordering.  The condition on
$\mathfrak m$ is exactly that its factors are at least $\mathfrak p$ and
$N(\mathfrak m/\mathfrak p)\le M<N\mathfrak m$.
Thus for each outer $\mathfrak m$, the remaining $\mathfrak w$ runs over
the complete set of $\mathfrak p$-smooth ideals prime to $\mathfrak a$.

Write $\mathfrak m=m C_i$ using the balanced choice above.  Then $x=mr$
with $r\in C_i$ and
$(r)C_i^{-1}=\mathfrak w\mathfrak u$.
Split $\mathfrak u$ into its prime factors at most $\mathfrak p$ and those
larger than $\mathfrak p$; call the latter product $\mathfrak z$.
The smaller part combines with $\mathfrak w$ to give every
$\mathfrak p$-smooth integral ideal, without a coprimality restriction.
Writing $\mathfrak z=z C_j$, one obtains $r=zy$ with
\[
 y\in C_iC_j,
 \qquad (y)(C_iC_j)^{-1}\text{ a }\mathfrak p\text{-smooth integral ideal}.
\]
The original box condition is simply its original weight evaluated at
$mzy$.  For fixed $\mathfrak m,\mathfrak z$ this inner sum is independent
of the conductor ideal.  All root lattices $C_i$ and $C_iC_j$ are fixed;
there are only finitely many of them.  This is the exact replacement for
the rational factor decomposition in \cite[\S\S3--4]{BW2023}.

\subsection{Integration and the exponent}

Choose a sufficiently small fixed $c>0$ and put
\[
                  M=cX Z^{-2/k}R^{-1}.
\]
If $M<R$, the desired assertion follows directly from the full moment:
$Z\gg_K X^{k/2}R^{-k}$, so the loss is at most $R^{2\Delta}$,
absorbed by choosing the smoothness exponent small enough.
We may therefore assume $M\ge R$.
There are $O_K(MR)$ choices of the ideal $\mathfrak m$.
H\"older's inequality applied to the preceding exact decomposition gives
a factor $(MR)^{s-1}$ and a sum of $s$th moments of the inner sums.
The bad part contributes at most
\begin{equation}\label{eq:restricted-bad}
                       X^{-k+\epsilon}Z^2M^s.
\end{equation}
Here the total measure of the product arcs is
$\ll_KX^{-k}Z^2(1+\log X)^{O_K(1)}$: a conductor of norm $q$ has at
most $q$ centers, each error region has volume
$\ll_K X^{-k}(Z/q)(1+\log Z)^{O_K(1)}$, and there are $O_K(Z)$ ideals
of norm at most $Z$.

For each fixed $\mathfrak m$, apply \eqref{eq:field-arc-scaling}.
Our choice of $c$ ensures disjointness on the quotient root scale because
$N\mathfrak m\le MR=cXZ^{-2/k}$.
There are only $X^\epsilon$ ideals $\mathfrak z$ supported on a given
conductor and of the relevant norm.  Another application of H\"older
therefore costs $X^\epsilon$.
Now extend the $\mathfrak z$ summation to all integral ideals in its
allowed size range.  For fixed $\mathfrak m,\mathfrak z$, the inner
integrand no longer depends on $\mathfrak a$.  The sum over the disjoint
conductor arcs is consequently at most a full torus integral.
Multiplication by $z^k$ maps $\mathbb T_{C_i}$ onto
$\mathbb T_{C_iC_j}$ and preserves probability Haar integrals.
The remaining full moment is bounded by
\[
 X^\epsilon\left(\frac{X}{N\mathfrak m N\mathfrak z}\right)^\lambda.
\]
Indeed its weights are nonnegative and bounded, its support lies in a
fixed box of that volume scale, and its roots lie in a fixed integral
ideal.  At even order, orthogonality bounds it by a constant times the
complete smooth moment on $\mathcal O_K$, after allowing the finitely
many fixed prime factors of $C_iC_j$.
This use of coefficient domination is only on a full torus integral.

The permissible bound is uniform for these smaller scales: for scales
$Y\ge X^\delta$, choose the original smoothness exponent at most
$\delta$ times a permissible one; for $Y<X^\delta$, use the trivial
bound and choose $\delta$ sufficiently small in terms of the allocated
$\epsilon$.  Bounded scales and the finitely many fixed primes cause
only a constant or $X^\epsilon$ factor.

Since $\lambda\ge s/2\ge1$, ideal counting gives
$\sum_{N\mathfrak z\le X}(N\mathfrak z)^{-\lambda}\ll_K X^\epsilon$,
including the logarithmic endpoint $\lambda=1$.
The good contribution is therefore at most
\[
 X^\epsilon(MR)^{s-1}
 \sum_{M<N\mathfrak m\le MR}(N\mathfrak m)^{-k}
             \left(\frac{X}{N\mathfrak m}\right)^\lambda
 \ll X^{s-k+\epsilon}(X/M)^\Delta.
\]
Inserting $M$ gives the asserted main bound.
The bad contribution \eqref{eq:restricted-bad} is
$\ll X^{s-k+\epsilon}Z^{2-2s/k}$, which is no larger because
$\Delta\ge k-s$ and $Z\ge1$.
H\"older interpolation proves the final assertion at real orders.
This completes the proof of Proposition~\ref{prop:restricted-field}.


\section{Quadratic fourth moments sampled over a number field}
\label{sec:quad-sampling-field}

Let $K$ be fixed, $n=[K:\mathbb Q]$, $d_v=[K_v:\mathbb R]$,
$V=K\otimes_{\mathbb Q}\mathbb R$, and
$\mathbb T=V/\mathfrak D^{-1}$. All size parameters in this section
are volumes. Write $\mathcal B(U)$ for a bounded archimedean box with
all radii $O(U^{1/n})$. Fixed dilations and translates are harmless.
For a vector $\mathbf T=(T_v)_{v\mid\infty}$ of positive real numbers put
\[
 A_{\mathbf T}(\beta)=\prod_v\max\{1,T_v|\beta_v|\}^{d_v},\qquad
 \mathcal H_{\mathbf T}(\alpha)=
       \min_{\gamma\in K} Q_\gamma A_{\mathbf T}(\alpha-\gamma),
 \qquad
 Q_\gamma=N\{a\in\mathcal O_K:a\gamma\in\mathfrak D^{-1}\}.
\]
When every $T_v=T^{1/n}$ we abbreviate the height to $\mathcal H_T$.
The minimum includes all lifts of a center. Product Dirichlet approximation
gives $\mathcal H_T(\alpha)\ll_K T^{1/2}$.

\begin{lemma}[Anisotropic bilinear and cluster estimates]
\label{lem:anisotropic-bilinear}
Let $\mathcal U,\mathcal W$ be archimedean boxes of radii $U_v,W_v$,
with volumes $J=\prod_vU_v^{d_v}\ge1$ and
$N=\prod_vW_v^{d_v}\ge1$. Put $T_v=U_vW_v$.
For any coefficients supported on $\mathcal W\cap\mathcal O_K$,
\[
 \sum_{z\in\mathcal U\cap\mathcal O_K}
  \left|\sum_m c_m e_K(\alpha zm)\right|^2
 \ll_K \left(J+N+\frac{JN}{\mathcal H_{\mathbf T}(\alpha)}\right)
               \sum_m|c_m|^2.
\]
Consequently the number of $z$ in $\mathcal U$ for which $\alpha z$
modulo $\mathfrak D^{-1}$ lies in a prescribed box of radii $O(W_v^{-1})$
is
\[
 \ll_K 1+\frac JN+\frac J{\mathcal H_{\mathbf T}(\alpha)}.
\]
The constants are independent of the eccentricities of the two boxes.
\end{lemma}
\begin{proof}
The proof of Lemma~\ref{lem:g789-bilinear} works
with these radii without alteration. We give the relevant uniformity details.
For $J\le N$, subdivide $\mathcal U$ by a fixed number of subdivisions
at each embedding. The product of the sample-box volume and a
frequency-error box volume is then sufficiently small for the affine
field-line argument: the determinant of any two difference pairs has
norm smaller than the fixed lower bound for a nonzero element of
$\mathfrak D^{-1}$. A cluster therefore lies on an affine $K$-line of
slope $\gamma$. Ideal-lattice packing bounds its size by
\[
 1+Q_\gamma^{-1}
       \prod_v\min\{U_v,W_v^{-1}/|\alpha_v-\gamma_v|\}^{d_v}
 =1+\frac J{Q_\gamma A_{\mathbf T}(\alpha-\gamma)}.
\]
The Gaussian large sieve then gives the stated bound, with
$N+JN/\mathcal H_{\mathbf T}(\alpha)$ on the right. When $J>N$,
transpose the bilinear matrix, interchanging the two boxes. Its operator
norm is unchanged and the product radii $T_v$ are unchanged. This proves
the displayed symmetric bound.

For the cluster assertion, put coefficient weight
$\exp(-c\sum_v|m_v/W_v|^2)$ on the dual box. Poisson summation gives a
nonnegative Gram kernel which is $\gg_K N$ for any pair of the selected
sample frequencies. If there are $M$ samples, the corresponding quadratic
form is therefore $\gg_K NM^2$. Its upper bound from the bilinear operator
estimate is $O_K((J+N+JN/\mathcal H_{\mathbf T})M)$.
Gaussian tails are summed by dyadic dilation and rapid decay. Division
gives the claim. The only lattice input throughout is
$\#(\mathfrak a\cap\mathcal R)\ll_K1+\operatorname{vol}(\mathcal R)/N\mathfrak a$
for an arbitrary product box, so no bounded-eccentricity hypothesis is used.
\end{proof}

\begin{proposition}[Distribution of short heights of multiples]
\label{prop:short-height-multiples}
Let $X,H\ge1$, put $V_0=HX^2$, and let
$L=\mathcal H_{V_0}(\alpha)$. For $1\le Y\ll_K X$ one has
\[
 \#\{0\ne h\in\mathcal B(H):\mathcal H_{X^2}(h\alpha)\le Y\}
 \ll_{K,\epsilon} V_0^\epsilon
       \left(\frac{HY^2}{X^2}+\frac{HY}{L}\right).
\]
\end{proposition}
\begin{proof}
Choose an approximation $h\alpha=\gamma+\beta$ of height at most $Y$.
Let $\mathfrak q$ be its conductor and write $Q=N\mathfrak q$.
Minkowski's theorem supplies $0\ne b\in\mathfrak q$ with
$|b_v|\ll_KQ^{1/n}$ at every embedding. Since $(b)\subset\mathfrak q$,
$|Nb|\ge Q$; consequently also $|b_v|\gg_KQ^{1/n}$ at every embedding.
Thus $|Nb|\asymp_K Q$ and $b\gamma\in\mathfrak D^{-1}$.
Put $m=bh$. It lies in a box of volume $O_K(HQ)$ and satisfies
\[
 \alpha m\equiv b\beta\pmod{\mathfrak D^{-1}}.
\]
Each nonzero $m$ has only $V_0^\epsilon$ possible bounded-height factors
$b$, by the ideal-divisor bound and the unit-lattice count.

Decompose $Q$ and
$\max(1,X^{2/n}|\beta_v|)$ into dyadic intervals. Denote the latter
upper endpoints by $A_v\ge1$ and their product by
$A=\prod_v A_v^{d_v}\ll_KY/Q$.
There are only logarithmically many such choices. For a fixed choice,
the residue error $b\beta$ is contained in a box of radii
$O_K(Q^{1/n}X^{-2/n}A_v)$. Apply
Lemma~\ref{lem:anisotropic-bilinear} with
\[
 U_v=(HQ)^{1/n},\qquad
 W_v=(X^2/Q)^{1/n}A_v^{-1},\qquad
 J=HQ,\quad N=\frac{X^2}{QA}.
\]
Here $N\gg_KX\ge1$, after harmless fixed rescaling, and
$T_v=V_0^{1/n}/A_v$.
For every approximation center,
$A_{V_0}(\alpha-\gamma)\le A\,A_{\mathbf T}(\alpha-\gamma)$.
Taking minima gives $L\le A\mathcal H_{\mathbf T}(\alpha)$.
The number of such $m$ is therefore
\[
 \ll_K1+\frac{HQ^2A}{X^2}+\frac{HQA}{L}
 \ll_K1+\frac{HY^2}{X^2}+\frac{HY}{L}.
\]
Finally, the approximation to $h\alpha$ lifts to one to $\alpha$
with center $\gamma/h$. Its conductor has norm at most $Q|Nh|$, and
\[
 |Nh|A_{V_0}(\beta/h)
   =\prod_v\max\{|h_v|,V_0^{1/n}|\beta_v|\}^{d_v}
   \ll_K H A_{X^2}(\beta).
\]
Thus the set being counted is empty unless $L\ll_K HY$.
On its nonempty range this absorbs the constant $1$ above.
Summing the dyadic choices and factor multiplicities proves the result.
\end{proof}

\begin{corollary}[Sampled fourth moment]
\label{cor:quadratic-fourth-field}
Let $Z_t(\beta)$ be a quadratic exponential sum on a box of root volume
$X$, with arbitrary linear modulation $t$. For arbitrary modulations $t_h$,
\[
 \sum_{0\ne h\in\mathcal B(H)}|Z_{t_h}(h\alpha)|^4
 \ll_{K,\epsilon}(HX^2)^\epsilon
 \left\{HX^2+X^4\min\left(\frac HL,\frac{H^2}{L^2}\right)\right\},
 \qquad L=\mathcal H_{HX^2}(\alpha).
\]
The assertion is uniform for fixed nonzero integral multipliers of $\alpha$.
\end{corollary}
\begin{proof}
One ordinary difference removes the linear modulation. The positive
linear-kernel estimate, proved by the bilinear Gram argument above, gives
\[
 |Z_t(\beta)|^2\ll X^\epsilon
       \left(X+\frac{X^2}{B}\right),\qquad
 B=\mathcal H_{X^2}(\beta)\ll_K X.
\]
Sharp fixed boxes can be completed with a logarithmic loss; smooth majorants
give the same estimate directly. In particular the bound is uniform in $t$.
Hence $|Z_t(\beta)|^4\ll X^\epsilon(X^2+X^4B^{-2})$.
Proposition~\ref{prop:short-height-multiples}, summed dyadically from
$Y\asymp_K\max(1,L/H)$ to $O_K(X)$, shows that
\[
 \sum_{0\ne h\in\mathcal B(H)}\mathcal H_{X^2}(h\alpha)^{-2}
 \ll (HX^2)^\epsilon
       \left\{\frac H{X^2}
        +\min\left(\frac HL,\frac{H^2}{L^2}\right)\right\}.
\]
Combining the last two bounds proves the assertion. A fixed multiplier
changes conductor heights only by fixed factors.
\end{proof}

In the eighth-power terminal, $H=X^h$ and $L\asymp X^\ell$.
Normalization by $HX^4$ gives
\[
 z_4(\ell)=\min\{2,\max(\ell,2\ell-h)\}.
\]
Both resources used by the exact policy,
$\min\{2,2\ell-h\}$ and $\min\{2,\ell,h\}$, are weaker and therefore valid.
Products of a fixed number of shifts may replace $h$, at the cost of the
same divisor loss. The quadratic second-moment resource remains the
previously proved bound with normalized saving $1$ when $H\ge X$ and
$L\gg X^{1-O(\eta)}$.




\section{Field seed moments for eighth and sixteenth powers}
\label{sec:field-seed-moments}

The seed curves needed below can be obtained over the field itself. In
particular, no entry of a rational-only moment table is used as a theorem.
We retain the volume convention $X=P^{[K:\mathbb Q]}$.
The principal-block efficient-differencing argument in Section~\ref{sec:g6-transfer}
has the same proof with $6$ replaced by any fixed $k\ge3$: the derivative of
$x^k$ is nonzero on the small ordinary supports, the class-group cost is
absorbed into a fixed power of the smoothness parameter, and every
congruence root count is a bounded-degree polynomial root count. We describe
all auxiliary counts required for the two seed curves.

\begin{lemma}[Sparse progression endpoint over a number field]
\label{lem:field-sparse-endpoint}
Let $p\ge2$ and $J\ge1$ be fixed integers. For $1\le s\le J+1$,
\[
 \#\left\{\mathbf x,\mathbf y\in\mathcal B(M)^s:
       \sum_{i=1}^s x_i^{pr}=\sum_{i=1}^s y_i^{pr}\ (1\le r\le J)\right\}
 \ll_{K,p,J,\epsilon} M^{s+\epsilon}.
\]
\end{lemma}
\begin{proof}
For $s\le J$, Newton's identities determine the multiset of $p$th powers
on one side from the other. Each powered coordinate has at most $p$ roots,
so there are $O_{p,J}(M^s)$ solutions. The case $J=1,s=2$ is the ordinary
fourth-moment bound for the degree-$p$ sum, supplied by Hua's differencing
lemma over $K$.

Suppose $J\ge2$ and $s=J+1$. Interchanging the first coordinates of the two
sides turns the equations into equality of signed power sums
\[
 S_r(\mathbf x)=\sum_{i=1}^J x_i^{pr}-x_0^{pr}
             =S_r(\mathbf y)\qquad(1\le r\le J).
\]
These equalities imply
\[
 D(\mathbf x)=\prod_{i=1}^J(x_0^p-x_i^p)
             =\prod_{i=1}^J(y_0^p-y_i^p)=D(\mathbf y).
\]
Here is a direct algebraic verification. If $x_0^p=y_0^p$, apply Newton's
identities to the other $J$ powered coordinates. Otherwise the two monic
polynomials with roots $(y_0^p,x_1^p,\ldots,x_J^p)$ and
$(x_0^p,y_1^p,\ldots,y_J^p)$ differ by a constant, because their first $J$
power sums agree. Evaluating their difference at $x_0^p$ and $y_0^p$ gives
$(x_0^p-y_0^p)D(\mathbf x)=(x_0^p-y_0^p)D(\mathbf y)$.

Fix $\mathbf y$ with $D(\mathbf y)\ne0$. Every difference
$d_i=x_0^p-x_i^p$ is a bounded-height divisor of this fixed nonzero algebraic
integer. There are $O(M^\epsilon)$ choices for the whole difference tuple.
Moreover
\[
 (J-1)x_0^p=S_1(\mathbf y)+\sum_i d_i,
\]
so the powers of all coordinates, and hence the coordinates up to a fixed
multiplicity, are determined. This gives $O(M^{J+1+\epsilon})$ solutions.
If $D(\mathbf y)=0$, choose a coordinate with $y_i^p=y_0^p$; likewise one
coordinate of $\mathbf x$ has the same power as $x_0$.
Cancel these in the signed sums. The remaining $J-1$ powered coordinates
are determined up to permutation by Newton's identities. There are
$O(M^J)$ choices for $\mathbf y$ and $O(M)$ choices for the free center
$x_0$, again with bounded root multiplicities. This proves the lemma.
\end{proof}

For efficient differences choose $0\le\phi_i<1/k$ and let
\[
 M_i=X^{\phi_i},\quad H_i=X M_i^{-k},\quad
 \Phi_j=\sum_{i=1}^j\phi_i,\quad
 U_j=\prod_{i=1}^jH_iM_i.
\]
Fixed small constants and fixed smoothness powers are suppressed in these
scale definitions. The centered polynomial is denoted by $\Psi_j(z;\mathbf
h,\mathbf m)$, and $F_j$ is its sum over the small ordinary box and the
independent shift and smooth-coefficient boxes. The original ordinary
vertices remain in a common small box away from zero at every embedding.
Thus $\Psi_j/\prod h_i$ and its required derivatives are nonzero wherever
they are used. This is also true after the finitely many ordinary differences
below, which have small shifts and coefficient divisor $1$.

\begin{lemma}[Auxiliary moments sufficient for the seed certificates]
\label{lem:field-general-seed-aux}
Let $1\le\ell\le3$ and $1\le j\le k-1-\ell$, and put
$J=\lfloor(k-j-\ell+1)/2\rfloor$.
Suppose either $j\le J+1$, or the chronological divisor inequalities
\begin{equation}\label{eq:field-chronological-prefix}
 \Phi_i+(k-2)\phi_i\le1\qquad(1\le i\le j)
\end{equation}
hold. Suppose also that for every $r=0,\ldots,\ell-1$, either
$k-j-r$ is odd or belongs to $\{2,4\}$, or $\Phi_j\le1/3$.
Then
\[
 \int_{\mathbb T_K}|F_j(\alpha)|^{2^\ell}\,d\alpha
   \ll X^{2^\ell-\ell+\epsilon}U_j^{2^\ell-1}.
\]
For $j=1$ the last condition is automatic, since $\phi_1\le1/k\le1/3$.
\end{lemma}
\begin{proof}
First consider a second moment, including $r$ appended ordinary shifts.
The nonzero phase-value divisor argument and Cauchy's inequality reduce
the count, at cost $X^\epsilon$, to the equation with a common shift tuple
on both sides. Polynomial coincidence implies the coefficient equations
\[
 \sum_{i=1}^j h_i^{2a}(m_i^{2ka}-n_i^{2ka})=0\qquad(1\le a\le J).
\]
The appended coefficient variables are all $1$ and cancel. The equivalence
between polynomial coincidence and these equations is triangular: the
centered-difference generating function is a constant multiple of
$e^{zt}\prod_i\sinh(h_im_i^kt)/(h_im_i^kt)$, whose logarithm has nonzero
even coefficients through the finitely many orders in question.

When $j\le J+1$, fix the shifts and apply H\"older to the $j$ coefficient
sums on the $J$-dimensional trace torus. Multiplication of its $a$th coordinate
by $h_i^{2a}$ preserves Haar measure. Lemma~\ref{lem:field-sparse-endpoint}
therefore gives $O(X^\epsilon\prod_iM_i)$ coefficient pairs. Summing shifts
gives $O(X^\epsilon\prod_iH_iM_i)$.
Alternatively, retain only the first coefficient equation and eliminate its
last block. If $m_i^{2k}=n_i^{2k}$, there are at most $2k$ choices for $n_i$
per $m_i$. Otherwise the equation
$h_i^2(m_i^{2k}-n_i^{2k})=b\ne0$ has $O(X^\epsilon)$ solutions: $h_i$ is a
bounded-height divisor of $b$, then $m_i-n_i$ divides the remaining fixed
integer, and a nonzero fixed-degree polynomial determines $n_i$ up to bounded
multiplicity. The resulting recurrence is
\[
 N_i\ll X^\epsilon\left(H_iM_i N_{i-1}
                      +\prod_{a<i}H_aM_a^2\right).
\]
The inequalities \eqref{eq:field-chronological-prefix} say precisely
$\prod_{a<i}M_a\le H_iM_i$ and yield the same coefficient bound.

For odd residual degree, the phase has a factor $z$. Divisor grouping in
$z$, followed by coefficient comparison, bounds the common-shift equation
by $X^{1+\epsilon}$ times the coefficient count, plus
$O(X^\epsilon\prod H_iM_i^2)$. The latter is absorbed since
$\prod M_i\le X$. Residual degree two is a difference of squares.
At degree four complete the two quadratics in $z^2,w^2$; a nonzero right
side is a divisor count, while a zero right side either gives polynomial
coincidence or a binary quadratic equation. The zero-level split lines
with a nontrivial root-of-unity slope miss the common tiny ordinary boxes.
The uniform quadratic norm-equation bound from Section~\ref{sec:g6-transfer} thus
handles both degrees with the same estimate.

For even residual degree $d\ge6$, write $\Psi_j/\prod h_i=P(z^2)$.
If the constant terms on the two sides agree, factor out $z^2,w^2$ and
use the same nonzero-value divisor and coefficient argument. If they differ,
fix the parameters and factor $P(x)-Q(y)$. Its constant is nonzero and its
highest form is a nonzero multiple of $x^{d/2}-y^{d/2}$.
Every irreducible factor of degree at least three has
$O(X^{2/3+\epsilon})$ integral points in the volume-$X^2$ box, by the
number-field affine-curve theorem of Binyamini--Cluckers--Kato
\cite[Corollary~12]{BCK2025}.
A linear factor has nonzero constant and coefficients of both variables;
after $x=z^2,y=w^2$ it is a nonzero-level binary quadratic equation.
A quadratic factor has nondegenerate highest form, since the slopes of
$x^{d/2}-y^{d/2}$ are distinct. Translation and diagonalization reduce it
to the uniform quadratic norm count; factors defined only over extensions
have at most boundedly many $K$-points unless defined over $K$.
All factor coefficients have polynomial height, so the bounds are uniform
in the parameters. We obtain
\[
 R_{j+r}^{(1)}\ll X^{r+1+\epsilon}U_j
          +X^{r+2/3+\epsilon}\Bigl(\prod_iH_i\Bigr)
                                   \Bigl(\prod_iM_i\Bigr)^2.
\]
The final term is absorbed when $\Phi_j\le1/3$.

Thus every appended second moment through $r=\ell-1$ is
$O(X^{r+1+\epsilon}U_j)$. Repeated ordinary differencing and Cauchy's
inequality, retaining the original parameter tuple until after differencing,
give
$\int|F_j|^{2^\ell}\ll X^{2^\ell-\ell+\epsilon}U_j^{2^\ell-1}$.
This is the same finite Weyl-moment induction as in the sixth-power proof,
now with at most two appended ordinary shifts. The exponent is obtained
by substituting the displayed appended second moments into that induction.
\end{proof}

\subsection{Details closing the auxiliary seed estimates}
\label{sec:seed-details}

No ordering of the exponents $\phi_i$ is needed for either the sparse
coefficient argument or the chronological divisor argument.  In particular,
the eighth-power row at half-order twenty is allowed to have
$\phi_2>\phi_1$.  Zero exponents are also allowed: $\phi_i=0$ means the
extraction threshold $M_i=1$, while the fixed smoothness powers in the
principal-block lemma contribute $X^{O(\eta)}$.  For any prescribed final
$\epsilon$, choose $\eta$ sufficiently small after the finite proof is fixed.
The degree-sixteen rows at half-orders twenty-one through twenty-three use
this convention for their eleventh extraction.  All actual divisors remain
nonzero.

\paragraph{The signed Newton endpoint.}
For completeness, the sparse $s=J+1$ endpoint follows from an invariant
of the first $J$ signed power sums.  With $a=x_0^p$ and $b_i=x_i^p$, put
$S_r=\sum_{i=1}^Jb_i^r-a^r$.  The first $J$ coefficients of
\[
 \frac{\prod_{i=1}^J(1-b_it)}{1-at}
     =\exp\left(-\sum_{r\ge1}\frac{S_r}{r}t^r\right)
\]
are determined by $S_1,\ldots,S_J$ in characteristic zero.  Its $t^J$
coefficient is $\prod_i(a-b_i)$.  Thus equality of the signed power sums
forces equality of this product.  If the product is nonzero, each difference
$a-b_i$ is a bounded-height algebraic-integer divisor of a fixed value.
The divisor and unit counts cost $M^\epsilon$.  Moreover
$(J-1)a=S_1+\sum_i(a-b_i)$ determines $a$ when $J\ge2$.
Every powered coordinate has only boundedly many $p$th roots.
If the product is zero, cancel one equal powered coordinate against the
center on each side.  Newton's identities determine the remaining $J-1$
powered coordinates up to permutation.  There are $O(M^J)$ choices for the
side with a coincidence and $O(M)$ choices for the other center.
The case $J=1,s=2$ is separately the fourth Hua moment.  This proves the
sparse estimate without a rational-only mean-value input.

\paragraph{From polynomial coincidence to coefficient equations.}
After fixing the common shift tuple, the residual centered polynomial has
generating function proportional to
\[
 e^{zt}\prod_i\frac{\sinh(h_im_i^kt)}{h_im_i^kt}.
\]
The coefficient comparison is triangular because the coefficients of
$\log(\sinh t/t)$ in degrees $2,4,\ldots,14$ are respectively
\[
 \frac16,\quad-\frac1{180},\quad\frac1{2835},\quad
 -\frac1{37800},\quad\frac1{467775},\quad
 -\frac{691}{3831077250},\quad\frac2{127702575},
\]
all nonzero.  These are all the orders used in the degree-eight and
sixteen certificates.  Consequently polynomial coincidence forces
\[
 \sum_i h_i^{2a}m_i^{2ka}=\sum_i h_i^{2a}n_i^{2ka}
                        \qquad(1\le a\le J).
\]
Appended ordinary differences have divisor variable one, so their terms
cancel on the two sides.  In the sparse case, H\"older on the $J$-fold
trace torus and measure preservation under multiplication of its $a$th
coordinate by $h_i^{2a}$ give the count $X^\epsilon\prod_iM_i$ for fixed
shifts.  Summing shifts gives $X^\epsilon U_j$.

For the chronological alternative, let $N_i$ count the first $i$ blocks
in the first coefficient equation.  A zero last summand contributes
$O(H_iM_iN_{i-1})$.  For a nonzero prescribed last summand $b$, factor
$h_i^2(m_i^{2k}-n_i^{2k})=b$.  First $h_i$ is a bounded-height divisor of
$b$; next $m_i-n_i$ is a divisor of the remaining nonzero value.  For every
such nonzero difference the resulting polynomial in $n_i$ has degree
$2k-1$ and nonzero leading coefficient.  Hence there are $O(X^\epsilon)$
possibilities.  This proves
\[
 N_i\ll X^\epsilon\left(H_iM_iN_{i-1}
                  +\prod_{a<i}H_aM_a^2\right).
\]
The exact condition for the second term is
$\prod_{a<i}M_a\le H_iM_i$, or equivalently
$\Phi_i+(k-2)\phi_i\le1$.  Thus induction gives $N_j\ll X^\epsilon U_j$
without any ordering assumption.

\paragraph{The low even residual degrees.}
At residual degree two the common-shift equation is a difference of
squares with fixed nonzero leading coefficient.  A nonzero right side
is a divisor count; a zero right side forces $z=w$ on the common tiny
box and gives coefficient coincidence.
At residual degree four write the two residual polynomials as
$az^4+b_mz^2+c_m$ and $aw^4+b_nw^2+c_n$, with $a\ne0$ fixed.
Set $u=2az^2+b_m$, $v=2aw^2+b_n$.  Then
\[
        u^2-v^2=b_m^2-b_n^2-4a(c_m-c_n).
\]
A nonzero right side is again a divisor count.  If it is zero, small
increments and the common nonzero center guarantee $u+v\ne0$ at every
embedding.  Thus $u=v$, reducing to a difference of squares in $z,w$.
Either that right side is nonzero, or $b_m=b_n$, $z=w$, and then
$c_m=c_n$.  This is precisely polynomial coincidence.  These arguments
work at complex embeddings because they use closeness to the same
nonzero center, rather than a sign condition.

At every higher even residual degree, the square-variable curve argument
in the seed lemma applies.  When the constant terms agree, the divisor
Cauchy step is performed after grouping by the common shift tuple, the
common constant term, and the remaining phase value.  That value is
nonzero on the retained small support.  When the constant terms differ,
the polynomial in $z^2,w^2$ has nonzero constant term and square-free
highest homogeneous part.  The linear and quadratic factors are handled
by the uniform norm-equation lemma; components of degree at least three
are handled by the coefficient-uniform number-field curve theorem.
With $r$ appended ordinary shifts this gives
\[
 R_{j,r}\ll X^{r+1+\epsilon}U_j+
 X^{r+2/3+\epsilon}\Bigl(\prod_iH_i\Bigr)\Bigl(\prod_iM_i\Bigr)^2.
\]
The second term is absorbed whenever $\Phi_j\le1/3$.
For odd residual degrees, the phase-value divisor argument in $z$ gives
the same first term, plus an error absorbed by $\prod M_i\le X$.
For every row this coefficient-product condition follows from
$0\le\phi_i<1/k$ and $j\le k-1-\ell$.

\paragraph{The padded ordinary-moment induction.}
Here is an explicit induction from the appended second moments to the
higher moments used by A and B.  Let $F_r$ denote the family with $r$
ordinary differences already appended, and set $U_r=X^rU_j$.
Cauchy in the parameter tuple and ordinary differencing give
\[
                     |F_r|^2\ll XU_r^2+U_r|F_{r+1}|.
\]
The zero ordinary shift is the first term; all nonzero shifts retain the
small-support restrictions.  Put
$B_{r,a}=\int|F_r|^{2^a}$.  For $a\ge2$, multiply the pointwise inequality
by $|F_r|^{2^a-2}$.  Interpolation in the first term and H\"older in the
second, followed by absorption, yield
\[
 B_{r,a}\ll (XU_r^2)^{2^{a-2}}B_{r,a-1}
                  +U_r^{2^{a-1}}B_{r+1,a-1}.
\]
Starting with $B_{r,1}\ll X^{1+\epsilon}U_r$, induction gives
\[
                  B_{r,a}\ll X^{2^a-a+\epsilon}U_r^{2^a-1}.
\]
Indeed, the two powers of $X$ obtained on substitution are
$3\cdot2^{a-2}-a+1$ and $2^a-a$, and the former is no larger than the
latter.  The restriction $j+\ell\le k-1$ means that every final residual
degree is at least two; the extra multiplicity from a residual linear
phase never occurs.  The proof retains the parameter tuple throughout the
ordinary differencing and consequently introduces no new independent
copies of the efficient-divisor variables.

\paragraph{Preserving the supports.}
All arguments above may be performed on the original supports.  If
Cartesian boxes are wanted, split the balanced extracted generators into
dyadic norm annuli and then small archimedean boxes.  On one piece their
conjugates have comparable fixed scales.  The shift ranges can be enlarged
by fixed factors while all ordinary vertices remain in a slightly larger
common tiny box.  The enlargement is made in nonnegative orthogonality
counts before an absolute-value bound is used.  The number of pieces is
$X^\epsilon$ and changes in scale are $X^{O(\eta)}$.
Choose the initial relative diameter small enough for the fixed total
number of differences.  The leading monomial of every required residual
polynomial, and of its relevant derivative, then dominates at every
embedding.  This justifies all exclusions of zero phase values above.

\paragraph{Exact row checks.}
An independent rational replay checks the analytic hypotheses separately
for every appended residual degree, not only through the shorthand
special-case flag.  It also checks every interpolation index and every
condition $M_i\le Q_{i-1}$ needed by divisor extraction.
The eighteen eighth-power rows have minimum strict supersolution slack
$1164945381/800000000000000$ and contraction constant at most
$976509627/1000000000$.
The twenty-eight sixteenth-power rows have minimum slack
$21545278901/80000000000000000$ and contraction constant at most
$9962585803/10000000000$.
All rows pass; no rational permissible-exponent table is an analytic
premise of this verification.

\subsection{The exact supersolution certificate}
Write $b_s$ for a permissible full exponent at moment $2s$ and $b(t)$ for
linear interpolation between adjacent integer indices. Each certificate
row specifies $s$, a depth $j$, an order $\ell$, an A or B terminal,
and rational $\phi_1,\ldots,\phi_j$. Put
\[
 q_i=1-\Phi_i,\qquad
 \kappa_\ell=2^\ell-\ell+(2^\ell-1)(j-(k-1)\Phi_j),\qquad u=s-j.
\]
For an A terminal set
\[
 C_j=2^{-\ell}\kappa_\ell+
        (1-2^{-\ell})q_j b\bigl(u/(1-2^{-\ell})\bigr).
\]
For a B terminal the additional integer $t$ specifies
$c=u/t$, $a=2(1-c)-2^{1-\ell}$ and
$b'=2^{1-\ell}-(1-c)$, and
\[
 C_j=a\kappa_{\ell-1}+b'\kappa_\ell+cq_jb(t).
\]
These are ordinary H\"older identities with nonnegative weights.
Proceeding backwards for $i=j-1,\ldots,1$, put
\[
 \begin{split}
 E_i&=\tfrac12\{i-(k-1)\Phi_i+(2(s-i-1)-1)\phi_{i+1}
                                      +q_i b(s-i+1)\},\\
 D_i&=i+1-(k-1)\Phi_i+\phi_{i+1}+q_{i+1}b(s-i-1),\\
 C_i&=E_i+\tfrac12\max(D_i,C_{i+1}).
 \end{split}
\]
The first-stage output is
\[
 T_s(b)=(2s-3)\phi_1+
             \max\{1+\phi_1+q_1b(s-1),C_1\}.
\]
The fundamental principal-block inequality proves each step over $K$.
All needed auxiliary hypotheses and interpolation ranges are checked exactly.

There are eighteen rows for $k=8$ (indices $3$ through $20$) and twenty-eight
rows for $k=16$ (indices $3$ through $30$).
The exact verifier gives $T_s(b)<b_s$ in every row, with minimum slacks
$1.45618172625\cdot10^{-6}$ and $2.693159862625\cdot10^{-7}$ respectively.
Every branch is monotone in the input moments and has coefficient sum at
most $q_1$. The two sup-norm contraction constants are at most
$0.976509627$ and $0.9962585803$. Clamp outputs below by the diagonal
exponent $s$ if necessary; this preserves both the supersolution and the
contraction bound. Starting from the elementary moments $b_0=0,b_1=1,b_2=2$
and the trivial $b_s=2s$ at larger indices, finite iteration approaches the
certified supersolution arbitrarily closely. Thus each printed exponent is
permissible with an arbitrary $X^\epsilon$ loss. Only finitely many stages
are used for a prescribed $\epsilon$, and the common smoothness parameter
is chosen after those stages.

Using these two field seed curves in the fixed twelve-round eighth-power
policy gives, by exact rational replay,
\[
 \Delta_{34}=\frac{5857852793}{50000000000}
              =0.117157055860<\frac2{17},
 \qquad
 4-34\Delta_{34}=0.016660100760>0.
\]
This last calculation uses the additional field resources established in
the other sections: the second coefficient moment, refined and cubic
coefficient sampling, sampled quadratic fourth moments, and restricted-height
smooth moments. It is not a consequence of the seed lemma alone.


\section{Averaged auxiliary major-arc bounds in residual degrees two and four}
\label{sec:aux-major-24}

All parameters are volumes.  Fix $K$, put $n=[K:\mathbb Q]$, and let
$d_v=1,2$ denote the dimensions of its archimedean places.  Write
$V=K\otimes\mathbb R$, $\mathbb T=V/\mathfrak D^{-1}$, with probability
Haar measure and the corresponding trace-dual archimedean normalization.
For $T\ge1$ use
\[
 A_T(\beta)=\prod_v\max\{1,T^{1/n}|\beta_v|\}^{d_v},\qquad
 Q_\gamma=N\{b\in\mathcal O_K:b\gamma\in\mathfrak D^{-1}\}.
\]
Let $k=8$, $j\in\{4,6\}$, and $d=k-j\in\{4,2\}$.  Put
\[
 M_i=X^{\phi_i},\quad H_i=XM_i^{-8},\quad
 M=\prod_{i\le j}M_i,\quad H=\prod_{i\le j}H_i,\quad
 Q_0=X/M,\quad T=Q_0^8=X^dH.
\]
The same proof also covers residual degree three, agreeing with the cubic
estimate proved earlier.

For each auxiliary tuple write $a=h_1\cdots h_j\ne0$ and
\[
 \Psi_j(z;\mathbf h;\mathbf m)=aP_{\mathbf h,\mathbf m}(z),
 \qquad P_{\mathbf h,\mathbf m}(z)=c_{8,j}z^d+\cdots,
 \qquad c_{8,j}\ne0.
\]
The leading coefficient is a fixed integer, and all other coefficients
are algebraic integers after clearing one fixed denominator if necessary.
The ordinary variable has a nonnegative smooth weight in the small box
retained by efficient differencing.  All scaled weight derivatives are
uniformly bounded.  Linear completion twists may be arbitrary, independently
for every auxiliary tuple.

\begin{proposition}\label{prop:aux-major-24}
Let $F_j$ be the resulting auxiliary sum.  On every broad arc
\[
 \alpha=\gamma+\beta,\qquad Q_\gamma A_T(\beta)\le C_KX,
\]
one has, for every $\epsilon>0$,
\begin{equation}\label{eq:aux-major-strong}
 |F_j(\gamma+\beta)|
 \ll_{K,\epsilon}X^{1+\epsilon}HM\,
            Q_\gamma^{-1/d}A_T(\beta)^{-1/2}.
\end{equation}
Consequently, with $L=Q_\gamma A_T(\beta)$,
\begin{equation}\label{eq:aux-major-height}
 |F_j(\gamma+\beta)|\ll X^{1+\epsilon}HM L^{-1/d}.
\end{equation}
The bound includes every Poisson mode; it has no omitted approximation
error.  In particular, when $d=2$ it does not depend on an absolute third
moment of the quadratic auxiliary sum.
\end{proposition}

\subsection{Uniform mixed complete sums}

The required complete-sum bound is slightly more general than the ordinary
monomial-plus-linear estimate.  If $P\in\mathcal O_K[z]$ has degree
$d\le4$ and fixed nonzero leading coefficient $c$, then for a primitive
character $\psi$ modulo $\mathfrak q$ and every integral $b$,
\begin{equation}\label{eq:aux-complete-general}
 (N\mathfrak q)^{-1}
 \left|\sum_{z\bmod\mathfrak q}\psi(P(z)+bz)\right|
 \ll_{K,c,d,\epsilon}(N\mathfrak q)^{-1/d+\epsilon},
\end{equation}
uniformly in the lower coefficients.  Those coefficients are allowed to
vary with $\mathbf h,\mathbf m$ and with the Poisson frequency.

Here is a local proof recording uniformity in the prime.  At a prime of
residue characteristic greater than $d$ and not dividing $c$, consider
an integral polynomial whose reduction has positive degree $r\le d$.
The normalized complete sum of conductor $\mathfrak p^E$ is
$O_d(q^{-E/r})$, where $q=N\mathfrak p$.
For $E=1$ this is the finite-field Weil bound (the linear case is zero).
For $E\ge2$, nonstationary residue classes contribute zero.
At a stationary residue class of multiplicity $m$ as a root of the reduced
derivative, translate to that class and substitute $z=z_0+\pi y$.
Let $s$ be the minimum valuation of the resulting nonconstant Taylor
coefficients.  Then $2\le s\le m+1\le r$; after dividing by $\pi^s$,
the new reduction degree $r'$ satisfies $1\le r'\le s$.
The class has measure $q^{-1}$, and induction supplies
\[
 q^{-1}O_d(q^{-(E-s)/r'})=O_d(q^{-E/r})
 \qquad(E>s).
\]
When $E\le s$, its trivial measure $q^{-1}$ is already at most $q^{-E/r}$.
A child with $r'=r$ can occur only when the reduced derivative has a single
root of multiplicity $r-1$.  Thus a same-degree recursion has only one
child and does not accumulate constants.  Every other branch lowers the
reduction degree; at most $r-1$ stationary residue classes occur.
Induction on reduction degree and conductor gives a constant depending
only on $d$ (one may take constants bounded by $(d-1)!$ for $d\le4$).
This proof applies over every finite residue extension and is unchanged
by ramification at such good primes.

At the finitely many remaining primes, the $d$th derivative is the fixed
nonzero constant $d!c$.  Cluckers' local-field van der Corput theorem
\cite[Proposition~3.3]{Cluckers2011} applies after putting
$f=(P+bz)/(d!c)$: then $f^{(d)}=1$, while the Gauss norm of
$f-f(0)$ is at most $|d!c|_\mathfrak p^{-1}$ uniformly in the integral
lower coefficients and in $b$.  Its permitted dependence on the fixed
local field, derivative order and Gauss norm consequently gives
$O_{K_\mathfrak p,c,d}(q^{-E/d})$ uniformly in these coefficients.
The fixed denominators from centering are absorbed into fixed fractional
lattices and fixed conductor factors.  Chinese remaindering and the
ideal-divisor bound absorb
products of the good-prime constants into $(N\mathfrak q)^\epsilon$.
This proves \eqref{eq:aux-complete-general}.

\subsection{Archimedean curvature and all Poisson modes}

For a fixed auxiliary tuple, put $P=X^{1/n}$ and write
$p_v(u)=P^{-d}P_{\mathbf h,\mathbf m,v}(Pu)$.
On the retained tiny support all increments $h_im_i^8$ are small relative
to the nonzero ordinary center.  The centered difference formula therefore
gives
\[
 p_v(u)=c_{8,j}u^d+O(\delta^2)
\]
in bounded $C^d$ norms, with $|u|\asymp_K1$.  Taking $\delta$ sufficiently
small ensures $|p_v''(u)|\gg_K1$ uniformly.  For $d=2$ this is immediate
from its fixed leading coefficient.
At a real place, the second-derivative test is uniform in any added linear
phase.  At a complex place, the Hessian determinant of
$\operatorname{Re}(t p_v(u))$ is a nonzero constant times
$-|t p_v''(u)|^2$, so two-dimensional stationary phase gives the corresponding
square decay.  After a fixed partition of the tiny support, the standard
uniform stationary-phase theorem \cite[\S7.7]{Hormander2003} and integration
by parts give
\begin{equation}\label{eq:aux-arch-tails}
 \left|\int w_v(u)e_{K_v}(t p_v(u)-\zeta u)\,du\right|
 \ll_{K,B}(1+|t|)^{-d_v/2}
       \left(1+\frac{|\zeta|}{1+|t|}\right)^{-B}
\end{equation}
for every prescribed $B$.  A linear completion twist simply translates
$\zeta$; the constants are independent of its size.  The same conclusion
holds for products of translated smooth vertex weights, since their
scaled derivatives remain uniformly bounded.  Sharp completion, if used,
only contributes an arbitrary $X^\epsilon$ factor.

Let $\mathfrak q$ be the conductor of $\gamma$ and
\[
 \mathfrak q_a=\mathfrak q/(\mathfrak q+(a)),\qquad q_a=N\mathfrak q_a.
\]
This is the conductor of $a\gamma$.
Poisson summation of the fixed-tuple sum modulo $\mathfrak q_a$, followed
by \eqref{eq:aux-complete-general}, bounds it by
\[
 Xq_a^{-1/d+\epsilon}A_{X^d}(a\beta)^{-1/2}
 \sum_{\xi\in\mathfrak q_a^{-1}\mathfrak D^{-1}}
 \prod_v\left(1+
   \frac{P|\xi_v-\xi_{0,v}|}{1+P^d|a_v\beta_v|}\right)^{-B},
\]
where $\xi_0$ is the arbitrary shift from completion.
The normalizing factor $1/q_a$ from Poisson has already been combined
with the complete-sum bound.

Ideal packing in a translated product box bounds the number of relevant
dual points by
\[
 O_K\left(1+\frac{q_a}{X}A_{X^d}(a\beta)\right).
\]
Indeed the radii are $(1+P^d|a_v\beta_v|)/P$ and their product volume is
$A_{X^d}(a\beta)/X$.  Since $|a_v|\ll H^{1/n}$,
\[
 q_aA_{X^d}(a\beta)\ll_K Q_\gamma A_T(\beta)\ll_KX.
\]
For completeness this packing also sums all tails.  Enlarge each radius
by an independent factor $2^{r_v}$; the number of dual points grows by at
most $\prod_v2^{d_vr_v}$, while \eqref{eq:aux-arch-tails} supplies
$\prod_v2^{-Br_v}$.  Choosing $B>n+2$ makes the resulting series convergent,
uniformly in $\xi_0$ and in the eccentricity of the gradient box.
Thus the whole dual sum is $O_K(1)$, proving the individual bound
\begin{equation}\label{eq:aux-fixed-tuple}
 |S_{\mathbf h,\mathbf m}(\gamma+\beta)|
 \ll X^{1+\epsilon}q_a^{-1/d}A_{X^d}(a\beta)^{-1/2}.
\end{equation}

\subsection{Averaging the conductor greatest common divisors}

There are $O(MX^\epsilon)$ choices of $\mathbf m$.
For prescribed $a\ne0$, the number of bounded-height tuples
$h_1\cdots h_j=a$ is $O(X^\epsilon)$, by ideal divisors and the unit
lattice.  We may therefore sum over $a\in\mathcal B(H)\setminus\{0\}$.
Write $Q=Q_\gamma$ and $A=A_T(\beta)$.  The elementary inequality
\[
 A_{X^d}(a\beta)\gg_K A|Na|/H
\]
gives
\[
 A_{X^d}(a\beta)^{-1/2}
 \ll_K\min\{1,(H/(A|Na|))^{1/2}\}.
\]
For an ideal $\mathfrak b$ of norm $b$, the small-norm ideal count
\[
 \#\{a\in\mathfrak b\cap\mathcal B(H):0<|Na|\le Y\}
       \ll H^\epsilon Y/b
\]
has no points for $Y<b$.  Dyadic summation consequently gives, including
thin scales,
\[
 \sum_{0\ne a\in\mathfrak b\cap\mathcal B(H)}
       \min\{1,(H/(A|Na|))^{1/2}\}
 \ll H^\epsilon\frac{H}{bA^{1/2}}.
\]
Split according to $\mathfrak b=\mathfrak q+(a)$, and enlarge each class
to all nonzero $a\in\mathfrak b$.  Equation~\eqref{eq:aux-fixed-tuple} then
yields
\begin{align*}
 |F_j(\gamma+\beta)|
 &\ll X^{1+\epsilon}M Q^{-1/d}
       \sum_{\mathfrak b\mid\mathfrak q}(N\mathfrak b)^{1/d}
         \frac{H}{(N\mathfrak b)A^{1/2}}\\
 &\ll X^{1+\epsilon}HM Q^{-1/d}A^{-1/2},
\end{align*}
since $\sum_{\mathfrak b\mid\mathfrak q}(N\mathfrak b)^{-1+1/d}
\ll Q^\epsilon$.  This proves \eqref{eq:aux-major-strong}.
Because $A\ge1$ and $d\ge2$, it implies \eqref{eq:aux-major-height}.

\subsection{The major terminal used by the certificate}

Suppose $v$ is an even smooth order with permissible excess $\Delta_v$.
The field restricted-height theorem bounds its mean on the union of arcs
of height at most $L$ by
\[
 Q_0^{v-8+\epsilon}L^{2\Delta_v/8}.
\]
On $1\le L\le C_KX$, sum dyadic shells using
\eqref{eq:aux-major-height}.  Provided $X\le Q_0^4$, this gives
\begin{equation}\label{eq:aux-major-terminal-8}
 \int_{\mathfrak M}|F_j(\alpha)g_{Q_0}(\alpha)^v|\,d\alpha
 \ll X^{1+\epsilon}HM Q_0^{v-8}
       X^{\max\{0,\Delta_v/4-1/(8-j)\}}.
\end{equation}
At equality in the exponent, the logarithm is absorbed in $X^\epsilon$.
For the quadratic policy $j=6,v=16$, the seed excess is less than two,
so the displayed maximum is zero.  This is precisely the required main
scale and proves the quadratic major contribution without an absolute
third-moment assertion.  The quartic policy uses $j=4$ and the same formula
with $1/(8-j)=1/4$.


% Complete terminal assembly for the eighth-power transfer.
\section{The quadratic terminal for eighth powers over a number field}
\label{sec:eighth-quadratic-terminal}

All parameters below are volumes. Fix a number field $K$ of degree $n$,
use $e_K(t)=\exp(2\pi i\operatorname{Tr}_{K/\mathbb Q}t)$, and give
$\mathbb T_K=(K\otimes\mathbb R)/\mathfrak D^{-1}$ probability Haar measure.
We state some formulas for a fixed degree $k$, but apply them only with
$k=8$ and efficient-differencing depths $j=4,5,6$.
Let
\[
 M_i=X^{\phi_i},\quad H_i=XM_i^{-k},\quad
 \Phi_i=\sum_{a\le i}\phi_a,\quad q_i=1-\Phi_i,
 \quad M=\prod_{i\le j}M_i,\quad H=\prod_{i\le j}H_i,
 \quad Q=X/M.
\]
Put $d=k-j$, $r=d-2$, $T=Q^k=X^dH$, and
\[
 A=XHM=X^{a_0},\qquad a_0=1+j-(k-1)\Phi_j,\qquad
 H_*=X^rH=X^{h_*},\quad h_*=k-2-k\Phi_j.
\]
The notation $\mathcal H_T(\alpha)$ denotes the conductor--product height
from the preceding sections. Product Dirichlet approximation gives
$\mathcal H_T(\alpha)\ll_K T^{1/2}$.
Fixed constants and factors $X^{O_K(\eta)}$ from smoothness are absorbed
into $X^\epsilon$ after the finite certificate is fixed.

\subsection{Cartesian supports before the terminal estimate}

The coefficient factors must be separated only after the ordinary
differences. We first justify the independent ranges that this requires.
At each counted efficient-differencing step, partition the extracted
balanced principal divisors into dyadic norm annuli and then finitely many
small archimedean boxes. Within one piece, every $|m_{i,v}|$ is bounded
above and below by fixed multiples of a prescribed radius. This costs only
a power of $\log X$, and changes the volume scales $M_i$ by at most
$X^{O_K(\eta)}$.

The original ordinary variables start in a sufficiently small fixed box
around a nonzero center at each embedding. In one divisor piece the actual
constraint $|h_{i,v}m_{i,v}^k|\ll\delta X^{1/n}$ can be enlarged to a
box for $h_i$ independent of $m_i$. Because the ratio of the largest and
smallest divisor radii on this piece is bounded, this enlargement increases
every vertex displacement by a fixed factor. The centered ordinary variable
can similarly be enlarged to one fixed small box independent of the
parameters. The constants grow only with the fixed differencing depth;
choosing the initial $\delta$ sufficiently small ensures that every enlarged
vertex, and every derivative support subsequently needed, stays in a small
box away from zero. Centering congruences modulo fixed powers of $2$ may be
discarded at this stage, after clearing the same fixed powers in the equation.

These enlargements occur in nonnegative orthogonality counts, before a
terminal absolute value is introduced. Consequently they give valid upper
bounds for the counted recurrence. They are not pointwise comparisons of
exponential sums. The resulting terminal sums have genuinely Cartesian
ranges: the box for $z$ is independent of every $m_i,h_i$, and each
coefficient set $\mathcal C_i$ is independent of the shifts. Working in a
fixed integral basis, one may take the ordinary box to be a parallelepiped;
this fits inside the prescribed tiny archimedean support. After ordinary
differencing, its intersections with translates depend only on the new
ordinary shifts. This is the support property used below.

\subsection{Ordinary differences and common completion}

For the centered efficient-difference polynomial one has the generating
expansion
\[
 \Psi_j(z;\mathbf h,\mathbf m)
 =2^jk!\Bigl(\prod_{i\le j}h_i\Bigr)
 \left\{\frac{z^d}{d!}
   +\frac{z^{d-2}}{6(d-2)!}\sum_{i\le j}h_i^2m_i^{2k}
   +\text{terms of degree at most }d-4\right\}.
\]
Apply $r=d-2$ ordinary differences with nonzero shifts $u_1,\ldots,u_r$.
With
\[
 \xi=2z+u_1+\cdots+u_r,\qquad
 \Gamma=\Bigl(\prod_{i\le j}h_i\Bigr)\Bigl(\prod_{a\le r}u_a\Bigr),
\]
the resulting polynomial is exactly
\begin{equation}\label{eq:eighth-exact-quadratic}
 \frac{2^jk!}{24}\Gamma
 \left(3\xi^2+\sum_{a\le r}u_a^2
                         +4\sum_{i\le j}h_i^2m_i^{2k}\right).
\end{equation}
This formula also covers $r=0$, in which case $\xi=2z$.
All constants are fixed; denominators can be cleared once.

Let $F_j$ be the original auxiliary sum and $\mathcal G_r$ its sum after
these ordinary differences, including the ordinary shifts. Cauchy's
inequality over the existing parameter tuple, before separating the
coefficient sums, gives
\begin{equation}\label{eq:eighth-unseparated}
 |F_j|^{2^r}\ll
       X^{2^r-1}(HM)^{2^r}
       +X^{2^r-r-1}(HM)^{2^r-1}|\mathcal G_r|.
\end{equation}
Only $r=0,1,2$ is needed. For $r=1$, the zero ordinary shift contributes
$X(HM)^2$ and the remaining term is $HM|\mathcal G_1|$.
For $r=2$, Cauchy on the original tuple and the first shift gives
$|\mathcal G_1|^2\ll X^3(HM)^2+XHM|\mathcal G_2|$; substitution gives
\eqref{eq:eighth-unseparated}. For $r=0$, it is the harmless bound
$|F_j|\le HM+|\mathcal G_0|$.

For each ordinary shift tuple, the remaining $z$-domain is an intersection
of translates of the same parallelepiped. Fourier completion in the fixed
integral basis gives an integrable kernel dominated by a common majorant
of $L^1$ norm $O((\log X)^n)$. The kernel may depend on $\mathbf u$, but
is independent of $\mathbf h,\mathbf m$. The condition
$\xi\equiv\sum_a u_a\pmod{2\mathcal O_K}$ is imposed by finitely many
additive characters, again independent of $\mathbf m$. Thus after one
common integration in a linear modulation $t$, the absolute value of
$\mathcal G_r$ is bounded by $X^\epsilon$ times
\begin{equation}\label{eq:eighth-separated-product}
 \sup_t\sum_{\mathbf h,\mathbf u}
       |Z_t(c_0\alpha\Gamma)|
       \prod_{i=1}^j|L_i(c_i\alpha\Gamma h_i^2)|,
\end{equation}
uniformly in $t$, where
\[
 Z_t(\beta)=\sum_{z\in\mathcal B(X)}e_K(\beta z^2+tz),\qquad
 L_i(\beta)=\sum_{m\in\mathcal C_i}e_K(\beta m^{2k}),
\]
and $c_0,c_i$ are fixed nonzero integers. Smooth completion gives the same
formula with bounded weights. No independent copies of the coefficient
tuple are introduced by \eqref{eq:eighth-unseparated}.

\subsection{The quadratic and coefficient resources}

On a minor height shell write
\[
 \mathcal H_T(\alpha)\asymp X^\ell,
 \qquad 1\le\ell\le\frac{kq_j}{2}.
\]
The fixed-factor endpoint changes and smoothness perturbations are
$O(1/\log X)+O_K(\eta)$ and are absorbed in the final error.
All products of a fixed number of nonzero algebraic integers have only
$X^\epsilon$ representations in the relevant boxes, by ideal divisors
and bounded-height unit counting. Thus $\Gamma$ may be grouped as a single
variable in $\mathcal B(H_*)$ for nonnegative moment estimates.

The quadratic second resource has normalized saving $1$ when $H_*\ge X$.
For completeness, expand
$\sum_{a\in\mathcal B(H_*)}|Z_t(c\alpha a)|^2$ after majorizing the
$a$-box by a positive Gaussian. The zero difference $z^2-w^2=0$ contributes
$O(XH_*)$. Each nonzero value of $(z-w)(z+w)$ has $X^\epsilon$ preimages.
The positive linear-kernel estimate is
\[
 \sum_{b\in\mathcal B(N)}D_{H_*}(\alpha b)
 \ll_K H_*+N+\frac{H_*N}{\mathcal H_{H_*N}(\alpha)};
\]
it follows from the Gaussian bilinear Gram bound and its cluster estimate.
At $N=X^2$, so $H_*X^2=T$, this gives
\[
 \sum_a|Z_t(c\alpha a)|^2
 \ll X^\epsilon\left(XH_*+X^2+H_*+\frac{T}{\mathcal H_T(\alpha)}\right)
 \ll X^{1+\epsilon}H_*.
\]
The modulation disappears on taking absolute values, so the bound is uniform
in the common $t$ from completion. The previously proved sampled fourth
moment gives normalized saving
$\min\{2,\max(\ell,2\ell-h_*)\}$.
The two weaker resources used by the certificate are consequently
\[
 z_2=1,\qquad z_4=\min\{2,\ell,h_*\},\qquad
 z_4^{\rm ref}=\min\{2,2\ell-h_*\}.
\]

For coefficient block $i$, group all shifts except $h_i$ into a variable
of volume
\[
 J_i=X^r\prod_{a\ne i}H_a,
 \qquad \log_X J_i=k-3-k\Phi_j+k\phi_i.
\]
The conductor identity is
$J_iH_i^3M_i^{2k}=T$.
For an even coefficient order $p$, raise $L_i$ to the integer power $p/2$.
Its Fourier coefficients have support in a box of volume
$N_i\ll X^\epsilon M_i^{2k}$, and their squared norm is at most
\[
 X^\epsilon M_i^{p-2k+\Delta_p^{(2k)}}.
\]
This follows by orthogonality and coefficient domination by the full
$2k$th-power smooth moment. The exponent $\Delta_p^{(2k)}$ is obtained
from the independently proved field seed curve. The refined coefficient
sampling estimate, with
$Y\le c_K\min\{N_i,J_i,(J_iN_iH_i^{-3})^{1/2}\}$, gives saving
\begin{align}
 b_{i,p}^{\rm ref}(\ell)
 &=\min\left\{2k\phi_i,\ k-3-k\Phi_j+k\phi_i,
       \frac{k-6-k\Phi_j+6k\phi_i}{2},\ \ell-2+2k\phi_i\right\}
       -\Delta_p^{(2k)}\phi_i.\label{eq:eighth-refined-resource}
\end{align}
The strengthened cubic-height sampling estimate gives the alternative
\begin{align}
 b_{i,p}^{\rm cube}(\ell)
 &=\min\left\{2k\phi_i,\ k-3-k\Phi_j+k\phi_i,
       \frac{k-6-k\Phi_j+6k\phi_i}{2},\ \ell/3\right\}
       -\Delta_p^{(2k)}\phi_i.\label{eq:eighth-cube-resource}
\end{align}
The coefficient normalization is exact: sampling produces
$J_iH_i M_i^{p+\Delta_p^{(2k)}}$ times the reciprocal saving, whereas the
trivial normalized size is $J_iH_iM_i^p$.

The basic second coefficient moment is supplied everywhere by the proved
field version of Vaughan--Wooley Lemma~4.3, with its parameter $d=k$.
Its allowable $U=X^b$ gives
\begin{equation}\label{eq:eighth-basic-resource}
 b_{i,2}^{\rm basic}=
 \min\left\{\phi_i,\ k-6-k\Phi_j+4k\phi_i,
               q_j/4,\ \frac{kq_j-(4k-1)\phi_i}{6}\right\}.
\end{equation}
On the present minor shells, the lemma's other term
$\mathcal H_T(\alpha)^{-1/k}$ is no larger than $X^{-b}$, since
$b\le\phi_i<1/k$ and $\ell\ge1$.
We use this full minimum even at depth four. Therefore no separate transfer
of the shallow coefficient corollary is required.

Each displayed coefficient resource bounds the sum of $|L_i|^p$, divided
by the common shift volume $H_*=X^rH$, over the \emph{same} shift tuple in \eqref{eq:eighth-separated-product}
by $X^\epsilon M_i^pX^{-b_{i,p}(\ell)}$.
The quadratic resources likewise bound these normalized sums of $|Z_t|^p$ by
$X^\epsilon X^pX^{-z_p(\ell)}$.
Grouping products only enlarges nonnegative sums, so no independence between
these factors is needed.

\subsection{One H\"older allocation, followed by the Weyl root}

Choose nonnegative resource weights $w_R$ on a fixed height interval.
If a resource has order $p_R$ and belongs to factor $i$, impose
\[
 \sum_{R\text{ on factor }i}p_Rw_R\le1
       \quad(i=0,1,\ldots,j),\qquad \sum_R w_R\le1.
\]
Factor $0$ denotes the quadratic sum. Divide each factor by its trivial
size, $X$ for $Z_t$ and $M_i$ for $L_i$. Its unused power is bounded
trivially, and the unused H\"older mass is assigned to the constant
function $1$. Generalized H\"older on the common shift space then gives
\[
 |\mathcal G_r|\ll X^{r+1+\epsilon}HM\,X^{-S(\ell)},\qquad
 S(\ell)=\sum_Rw_Rb_R(\ell).
\]
Combining this with \eqref{eq:eighth-unseparated} proves
\begin{equation}\label{eq:eighth-pointwise-terminal}
 |F_j(\alpha)|\ll A X^{\epsilon-
               2^{-(k-j-2)}\min\{1,S(\ell)\}}.
\end{equation}
In particular, the Weyl root is applied to the entire combined saving
$S(\ell)$, after all resources have been coupled by one H\"older inequality.

Let the remaining smooth half-order in the terminal integral be $u$.
Normally take $v=2u$, $a=0$, $c=1$, $\gamma=1$.
For the depth-six policy take instead
\[
 u=8,\quad v=20,\quad c=2u/v=4/5,\quad
 a=1-c=1/5,\quad\gamma=1-2a=3/5.
\]
The chronological divisor conditions
$\Phi_i+(k-2)\phi_i\le1$ supply the auxiliary bound
$\int|F_6|^2\ll X^\epsilon A$.
H\"older on a height shell gives
\[
 \int_{\rm shell}|F_j|\,|g_Q|^{2u}
 \le (\sup_{\rm shell}|F_j|)^\gamma
      \left(\int|F_j|^2\right)^a
      \left(\int_{\rm shell}|g_Q|^v\right)^c,
\]
where the middle factor is absent when $a=0$.
The identities $a+c=1$, $\gamma+2a=1$, and $cv=2u$ verify all powers.
The field restricted-height mean theorem and
\eqref{eq:eighth-pointwise-terminal} give the shell exponent
\begin{equation}\label{eq:eighth-shell-exponent}
 C(\ell)=(\gamma+a)a_0+c(v-k)q_j+
             \frac{2c\Delta_v^{(k)}}{k}\ell
       -\gamma 2^{-(k-j-2)}\min\{1,S(\ell)\}.
\end{equation}
Every resource is a minimum of affine functions of $\ell$ and is concave.
Hence $C(\ell)$ is convex on each interval with fixed resource weights.
Its maximum is attained at an endpoint, justifying the finite checks of
all height intervals.

On the broad major arcs $\mathcal H_T(\alpha)\le C_KX$, the averaged
Poisson estimate proved separately is
$|F_j(\alpha)|\ll X^\epsilon A\mathcal H_T(\alpha)^{-1/(k-j)}$.
Use the original smooth order $2u$, without the auxiliary redistribution.
Dyadic height summation with the field restricted-height mean gives
\begin{equation}\label{eq:eighth-major-exponent}
 C_{\rm maj}=a_0+(2u-k)q_j+
               \max\{0,2\Delta_{2u}^{(k)}/k-1/(k-j)\}.
\end{equation}
A logarithmic equality case is absorbed into $X^\epsilon$.
No absolute third moment for the quadratic auxiliary sum is asserted or used.

\subsection{The four policies and their counted recurrence}

The fixed policies have the following terminal data:
\[
\begin{array}{c|c|c|c|c|c}
 2s&j&(r_1,\ldots,r_j)&v&h_*&2^{-(k-j-2)}\\\hline
28&6&(13,12,11,10,9,8)&20&1.69428296&1\\
30&5&(14,13,12,11,11)&22&1.83213176&1/2\\
32&5&(15,14,13,12,11)&22&1.78352592&1/2\\
34&4&(16,15,14,14)&28&2.39658688&1/4
\end{array}
\]
Thus the quadratic second resource is available in every policy.
Only quadratic orders two and four are used. Coefficient orders are even
and at most $60$; the degree-sixteen curve through half-order $30$ supplies
every requested moment. The sole cube resource occurs at depth six and
coefficient order $60$. All other higher coefficient resources are refined
ones. The exact verifier checks the two quadratic fourth resources,
the full basic minimum \eqref{eq:eighth-basic-resource}, and every factor
load and H\"older mass.

For clarity, let $\lambda_t$ denote the current full exponent at smooth
moment order $2t$. Start $C$ with the maximum of
\eqref{eq:eighth-major-exponent} and every endpoint value of
\eqref{eq:eighth-shell-exponent} in the assigned intervals.
The principal-block counted recurrence gives, for $i=j-1,\ldots,1$,
\[
 \begin{split}
 D_i&=1+i-k\Phi_i+\Phi_{i+1}+q_{i+1}\lambda_{r_{i+1}},\\
 B_i&=\tfrac12\{q_i\lambda_{2r_i-r_{i+1}}
                   +i-(k-1)\Phi_i+(2r_{i+1}-1)\phi_{i+1}\},\\
 C&\longleftarrow B_i+\tfrac12\max(D_i,C).
 \end{split}
\]
The first-stage exponent is then
\[
 \lambda_s^{\rm new}=(2s-3)\phi_1+
           \max\{1+\phi_1+q_1\lambda_{s-1},C\}.
\]
Every subsisting half-order is integral, every backward stage satisfies
$0<2r_i-r_{i+1}<2r_i$, and each round reads only the exponents available
at its start. More precisely, fix a sufficiently small ordinary root box
$\mathcal B$ away from zero at every archimedean place, with a contained
ideal-smooth root set $\mathcal A$. Write $f_{\mathcal B}$ and
$g_{\mathcal A}$ for their unweighted exponential sums. The first-stage
formula bounds the localized joint moment
$\int|f_{\mathcal B}|^2|g_{\mathcal A}|^{2s-2}$ itself, with the support
enlargements described at the start of this section. Coefficient domination
in its nonnegative orthogonality count gives
\[
 \int|g_{\mathcal A}|^{2s}
 \le \int|f_{\mathcal B}|^2|g_{\mathcal A}|^{2s-2}.
\]
A fixed compact root annulus can be covered by finitely many such boxes;
$|\sum_{a=1}^N g_a|^{2s}\le N^{2s-1}\sum_{a=1}^N|g_a|^{2s}$ then gives
the required pure smooth moment on that annulus. Bounded nonnegative
weights are handled by the same coefficient domination, with constants
depending on their fixed bounds. These formulas coincide term by term with
the standalone exact verifier.
The support preparation, analytic estimates, and parameter choices above
therefore supply its complete field terminal interface.


\section{The eighth-power certificate and the uniform global conclusion}
\label{sec:eighth-global}

The three obligations left open in the preceding version are now supplied
by Sections~\ref{sec:field-seed-moments}--\ref{sec:eighth-quadratic-terminal}:
the degree-eight and degree-sixteen seed moments, the averaged auxiliary
major estimates, and the quadratic terminal with a common completion.
All additional ingredients are already proved in this manuscript. In
particular, the restricted-height mean is Proposition~\ref{prop:restricted-field},
the sampled quadratic fourth moment is
Corollary~\ref{cor:quadratic-fourth-field}, and the basic coefficient
second moment is Proposition~\ref{prop:k5-coeff2} with $d=8$.

\begin{theorem}[The field moment certified here]
\label{thm:eighth-moment}
For every fixed number field $K$ and every $\epsilon>0$, there is
$\eta>0$ such that the degree-eight ideal-smooth sum satisfies
\[
 \int_{\mathbb T_K}|g_8(\alpha;X,R)|^{34}\,d\alpha
 \ll_{K,\epsilon}X^{26+\Delta_{34}+\epsilon}
 \qquad(1\le R\le X^\eta),
 \qquad
 \Delta_{34}=\frac{5857852793}{50000000000}.
\]
The same assertion holds for the fixed boxes and bounded nonnegative
weights used in the representation argument.
\end{theorem}

\begin{proof}
The exact program in Appendix~\ref{app:eighth-field} contains all eighteen
degree-eight seed rows, all twenty-eight degree-sixteen seed rows, and
the four terminal policies, including every height interval and resource
weight. Its first part checks every seed supersolution inequality and
all side conditions used in the seed proof. The contraction argument in
Section~\ref{sec:field-seed-moments} starts these curves from the elementary
second and fourth moments and the trivial higher bounds.

The second part makes twelve simultaneous rounds at half-orders
$14,15,16,17$. Each round reads the exponents available at its beginning.
Before and after a round, H\"older interpolation over every pair of
available integer half-orders may improve an entry. This is a valid
operation on permissible field exponents. Each policy checks all
H\"older masses and factor loads, every coefficient-moment index, each
height interval endpoint, and the backwards counted recurrence.
Section~\ref{sec:eighth-quadratic-terminal} proves that these are the
exponents of the actual field estimates. Convexity of the shell exponent
justifies checking only the two endpoints of each of the eight assigned
height intervals.

After each policy application the exponent is rounded \emph{up} to a
multiple of $10^{-12}$, and the better of that value and the previously
available interpolated value is retained. All calculations use rational
arithmetic. The final outputs are:
\[
\begin{array}{c|c|c}
 \text{moment order}&\text{full permissible exponent}&\text{excess over }2s-8\\\hline
 28&20.360308744711&0.360308744711\\
 30&22.258467366407&0.258467366407\\
 32&24.178873819111&0.178873819111\\
 34&26.117157055860&0.117157055860
\end{array}
\]
These are certified upper values, not optimal-exponent assertions.
The last entry is the exponent in the theorem. Only finitely many
applications are needed for any prescribed final $\epsilon$; choose
the smoothness parameter after those applications and after the fixed
rational parameters. This absorbs the smoothness and divisor losses
without consuming the strict final margin.
\end{proof}

\subsection{The strict global margin}

The exact gate calculation is
\[
 \frac2{17}-\Delta_{34}
   =\frac{416502519}{850000000000}>0,
 \qquad
 D=4-34\Delta_{34}
   =\frac{416502519}{25000000000}
   =0.016660100760>0.
\]
For a simpler printed exponent one may weaken the moment to
$\Delta_{34}=0.11716=2929/25000$. Then
\[
 D=\frac{207}{12500},\qquad \frac{D}{72}=\frac{23}{100000}.
\]
Thus the transfer is not based on an equality or a numerical rounding
at the admissibility threshold.

\begin{theorem}[The additional uniform transfer]
\label{thm:eighth-transfer}
For every number field $K$,
\[
                         \boxed{G_{\mathbb K}(8)\le38.}
\]
The meaning of $G_{\mathbb K}$ and of uniformity is the one in
Section~\ref{sec:conclusions}.
\end{theorem}

\begin{proof}
Apply the epsilon-removal bootstrap in Section~\ref{sec:product-endgame}
with
\[
             (k,u,S,w)=(8,34,38,72).
\]
Its strict condition is exactly $D=4-34\Delta_{34}>0$, so it gives
$U_{38}\ll X^{30}$. The representation integral uses four ordinary
eighth-power roots and thirty-four product roots. Each product is one
algebraic integer and therefore contributes one Waring summand.
The ordinary localized seventy-second moment, the certified thirty-fourth
smooth moment, and the bootstrap give
\[
 \int_{\mathfrak m_*}|f_8|^4|h_8|^{34}
 \ll X^{30-23/100000+\epsilon},
\]
using the simpler weakened excess above. The broad-major annulus has
a fixed power saving by the ordinary thirty-eighth moment and the
product-root moment domination. The core Poisson formula gives a
positive main term of order $X^{30}$.

The local eighth-power Theorem~\ref{thm:local-eighth-thirtyfive} supplies a
primitive representation in at most thirty-five variables for every
element of the local additive eighth-power order. Padding with zero
gives the required thirty-eight-variable local solution, including at
the primes above $2$. The resulting singular series is positive.
The chosen positive real boxes and the complex boxes give a positive
singular integral after balancing the target by a totally positive
unit's eighth power. Undoing that balancing preserves the real signs
and the root norm bounds. This proves the stated Waring--Siegel assertion
for every sufficiently large eligible target.
\end{proof}

The bound $38$ is the project's stronger rational target, rather than
the published comparison bound $39$ in the cited literature. The field
proof uses slightly enlarged seed values and gives the excess printed
above, not an unchanged rational numerical exponent. No historical
priority or optimality assertion is part of this theorem.



\section{The ninth-power task that remains}
\label{sec:remaining-continuation}

The additional successful exponent in this continuation is eight.
Research on new exponents stops at that success, in accordance with
the requested stopping rule. The previously proved fifth-, sixth-,
and seventh-power transfers are retained in this expanded manuscript.

For ninth powers the local primitive bound $30$ and the forty-six-variable
product-root endgame are already available. The remaining sufficient
field moment condition is
\[
 \int_{\mathbb T_K}|g_9|^{40}
       \ll X^{31+\Delta_{40}+\epsilon},
 \qquad \Delta_{40}<\frac3{22}.
\]
The project's rational certificate crosses that threshold. The simpler
field route explored earlier, using auxiliary second and fourth moments
with a cubic terminal, reached an excess of about $0.14384$ after the
tested changes of integer path. This is above $3/22\approx0.13636$.
That finite exploration is neither a lower bound on attainable exponents
nor evidence of impossibility.

The present work removes several shared obstacles. The conductor sampling,
restricted-height theorem, common-completion argument, and the general
mixed complete-sum estimate are reusable over $K$. The seed arguments
also give methods for checking the degree-nine and degree-eighteen
auxiliary curves. What remains is their complete field verification
at the required parameter rows and the assembly and exact replay of the
ninth-power policy. No conclusion $G_{\mathbb K}(9)\le46$ is asserted here.

\subsection{Verification and scope}

The new exact program includes all numerical data and uses only Python's
standard library. Its seed checks were also replayed independently,
including every appended residual degree, each sparse or chronological
condition, the extraction ranges, and all interpolation indices.
The completion, Poisson, ideal-divisor, curve, and H\"older arguments
were checked separately from the arithmetic. Exact arithmetic establishes
the consequences of those mathematical estimates; it does not supply
the estimates themselves.

Only the numbers of summands are uniform over number fields. The
sufficiently-large threshold, the implicit constants, the fixed boxes,
and the allowable smoothness exponent may depend on $K$. The results
apply to the actual additive order generated by $k$th powers, with the
positivity and root norm conventions stated at the start. They do not
assert an exact value of a Waring function or a general logical
equivalence between every formulation of the number-field circle method.


\appendix
\section{Exact verification of the eighth-power field certificate}
\label{app:eighth-field}
This standalone program contains the eighteen degree-eight and twenty-eight
degree-sixteen seed rows, followed by all four fixed terminal policies.
It checks the seed side conditions and supersolutions, the coefficient
resource ranges, every H\"older mass and factor load, the height endpoints,
and twelve simultaneous rounds with upward rounding. All numerical
operations use exact fractions. The mathematical premises are proved
in Sections~\ref{sec:field-seed-moments}--\ref{sec:eighth-global}.
The only runtime requirement is Python's standard library.
\begin{lstlisting}
#!/usr/bin/env python3
"""Exact certificate; all analytic premises are proved in the manuscript."""
from fractions import Fraction as F
import json

SEEDS={
8:([0, 1000000, 2000000, 3046989, 4216499, 5491581, 6856667, 8310610, 9842873, 11452921, 13128318, 14866494, 16656149, 18490145, 20362923, 22269027, 24200341, 26150576, 28115091, 30089898, 32072146],[
(3, 'A', 1, 0, [234903730], False),
(4, 'A', 2, 0, [574019066], False),
(5, 'A', 1, 0, [727052915, 221429275], False),
(6, 'A', 2, 0, [809783242, 562805614], False),
(7, 'B', 2, 7, [882581976, 734440434, 263392407], False),
(8, 'A', 2, 0, [935533228, 827985009, 558042871], False),
(9, 'A', 2, 0, [990797142, 905540566, 749313767], False),
(10, 'A', 2, 0, [1031597344, 961403162, 836579341, 548998511], True),
(11, 'A', 2, 0, [1074221665, 1012236938, 922288755, 740466992], True),
(12, 'A', 2, 0, [1106949463, 1055547399, 977035374, 841990311, 527947234], True),
(13, 'A', 2, 0, [1135604602, 1090153632, 1024821569, 912379896, 684837846], True),
(14, 'A', 2, 0, [1162117406, 1123542368, 1070247412, 949156116, 813562385], True),
(15, 'A', 2, 0, [1186363867, 1156892994, 1093820448, 937560384, 803623186], True),
(16, 'A', 2, 0, [1204529686, 1185105398, 1087194987, 931881418, 798755501], True),
(17, 'A', 2, 0, [1218149633, 1206087867, 1082251785, 927644387, 795123760], True),
(18, 'A', 2, 0, [1228585446, 1222118985, 1078470795, 924403538, 792345890], True),
(19, 'A', 2, 0, [1236077788, 1234109194, 1075687573, 922017920, 790301074], True),
(20, 'A', 2, 0, [1241511020, 1242754281, 1073676385, 920294044, 788823466], True),
]),
16:([0, 1000000, 2000000, 3007492, 4063843, 5178749, 6341235, 7545733, 8790255, 10076926, 11401974, 12767190, 14170227, 15611941, 17090638, 18659938, 20265171, 21900470, 23564877, 25256429, 26973004, 28712662, 30474472, 32257964, 34062610, 35888365, 37735550, 39605675, 41503017, 43442618, 45476950],[
(3, 'A', 1, 0, [37414197], False),
(4, 'A', 1, 0, [188298668], False),
(5, 'A', 2, 0, [291922104], False),
(6, 'A', 1, 0, [337019148, 82529262], False),
(7, 'B', 2, 8, [361381975, 240453276], False),
(8, 'A', 2, 0, [378851332, 307163452], False),
(9, 'B', 2, 9, [397615430, 359291438, 199227299], False),
(10, 'A', 2, 0, [410254477, 382831518, 288575973], False),
(11, 'B', 2, 10, [424766731, 404882933, 361222671, 184037872], False),
(12, 'A', 2, 0, [436525663, 419528978, 385915026, 279229557], False),
(13, 'B', 2, 12, [449362623, 434343819, 413375685, 363298780, 181075949], False),
(14, 'A', 2, 0, [460814412, 446935396, 427851512, 391200512, 275915456], False),
(15, 'B', 2, 13, [488601671, 459517246, 444545990, 421601419, 370044564, 192789497], True),
(16, 'A', 2, 0, [488983581, 486345145, 456811263, 437557944, 399379876, 289341394], True),
(17, 'B', 2, 14, [494240208, 487641647, 484032399, 453965752, 430129669, 375646281, 203018466], True),
(18, 'B', 3, 14, [501982135, 493143959, 485876194, 480497859, 447400164, 410595812, 309910596], True),
(19, 'B', 3, 14, [510036657, 501140788, 492112716, 484313475, 477582628, 442812982, 397540410], True),
(20, 'A', 1, 0, [517846379, 509348577, 500481232, 491595321, 484221489, 479012032, 451728176, 428542946, 376483961], True),
(21, 'A', 1, 0, [525299753, 517193147, 508723316, 500003214, 491590441, 485595778, 464772956, 433788092, 404868886, 318822420, 0], True),
(22, 'A', 1, 0, [532660084, 524710546, 516677372, 508465683, 500513796, 494301898, 461511374, 430743949, 402027686, 318917324, 0], True),
(23, 'B', 2, 19, [539892466, 532194248, 524483527, 517083988, 510597893, 491716524, 458935423, 428339728, 399783746, 318600313, 0], True),
(24, 'B', 2, 19, [546801742, 539548829, 532258365, 525570370, 520805035, 489001043, 456400974, 425974242, 397575959, 369598125, 130565745], True),
(25, 'B', 2, 19, [553499949, 546686983, 540142819, 534592848, 521671826, 486893704, 454434124, 424138516, 395862615, 369471774, 218626657], True),
(26, 'A', 2, 0, [559913658, 553542379, 547645930, 543288863, 519707277, 485060126, 452722784, 422541265, 394371847, 368080391, 272451775], True),
(27, 'A', 2, 0, [566242010, 560057990, 554744632, 551317389, 517842531, 483319696, 451098383, 421025157, 392956813, 366759692, 324088483], True),
(28, 'B', 3, 22, [573083193, 566589653, 561778491, 553236577, 516354138, 481930529, 449801827, 419815039, 391827369, 365705545, 341325175], True),
(29, 'B', 3, 22, [582922068, 573528102, 568700817, 551656600, 514879494, 480554194, 448517248, 418616098, 390708358, 364661134, 340350392], True),
(30, 'A', 2, 0, [604502390, 583157999, 575827996, 549100774, 512494055, 478327785, 446439266, 416676648, 388898205, 362971658, 338773547, 289139661], True),
]),
}

report={}
for k,(ints,rows) in SEEDS.items():
 b=[F(x,10**6)for x in ints];bounds=[];minphi=F(1)
 def at(t):
  t=F(t);assert 0<=t<=len(b)-1
  i=t.numerator//t.denominator
  return b[i] if t==i else (i+1-t)*b[i]+(t-i)*b[i+1]
 for S,kind,ell,t,nums,prefix in rows:
  phi=[F(x,10**10)for x in nums];j=len(phi);p=2**ell;u=S-j
  assert 1 <= ell <= 3 and 1 <= j <= k-1-ell
  assert all(0<=f<F(1,k) for f in phi)
  assert all(phi[i] <= 1-sum(phi[:i]) for i in range(j))
  minphi=min(minphi,phi[0]);total=sum(phi)
  def special(l):return j==1 or (l==1 and ((k-j)%2 or k-j in (2,4))) or (l==2 and 3<=k-j<=5) or (l==3 and k-j in (4,5))
  used=[ell] if kind=='A' else [ell-1,ell]
  for order in used:
   assert 1 <= order <= 3
   for appended in range(order):
    residual = k-j-appended
    assert residual >= 2
    if residual % 2 == 0 and residual >= 6:
     assert total <= F(1,3)
  assert all(special(l) or total<=F(1,3) for l in used)
  J=(k-j-ell+1)//2
  if not prefix:assert j<=J+1
  else:
   for i in range(j):assert sum(phi[:i+1])+(k-2)*phi[i]<=1
  def kap(l):p=2**l;return p-l+(p-1)*(j+(1-k)*total)
  if kind=='A':
   c=1-F(1,p);C=kap(ell)/p+c*(1-total)*at(F(u)/c)
  else:
   c=F(u,t);a=2*(1-c)-F(2,p);bb=F(2,p)-(1-c)
   assert min(a,bb,c)>=0 and a+bb+c==1 and a*(p//2)+bb*p==1
   C=a*kap(ell-1)+bb*kap(ell)+c*(1-total)*at(t)
  for i in range(j-1,0,-1):
   pre=sum(phi[:i]);ui=S-i-1;ti=S-i+1
   E=(i+(1-k)*pre+(2*ui-1)*phi[i]+(1-pre)*at(ti))/2
   D=i+1+(1-k)*pre+phi[i]+(1-pre-phi[i])*at(ui)
   C=E+max(D,C)/2
  out=(2*S-3)*phi[0]+max(1+phi[0]+(1-phi[0])*at(S-1),C)
  assert S<=out<b[S],(k,S,float(out),float(b[S]),float(b[S]-out))
  bounds.append({'s':S,'b':str(b[S]),'out':str(out),'slack':str(b[S]-out)})
 print(k,'all',len(bounds),'rows passed; minimum slack',float(min(F(x['slack'])for x in bounds)),'contraction <=',float(1-minphi))
 report[k]={'rows':bounds,'contraction':str(1-minphi)}

# All frozen policies, in literal rational data.
POLICIES=json.loads(r'''
{
"14":{
"phi":[
"5793479/50000000","11191889/100000000","10604341/100000000",
"9516681/100000000","3846939/50000000","806929/25000000"],
"path":[13, 12, 11, 10, 9, 8],"v":20,
"pieces":[
["0","1/8",[
[["z2",-1,2],"99999/200000"],
[["refined",0,44],"22727/1000000"],
[["refined",1,46],"10869/500000"],
[["refined",2,48],"20833/1000000"],
[["basic",3,2],"1659/4000"],
[["refined",3,52],"1639/500000"],
[["cube",4,60],"8333/500000"]
]],
["1/8","1/4",[
[["z2",-1,2],"99999/200000"],
[["refined",0,40],"24999/1000000"],
[["refined",1,42],"23809/1000000"],
[["refined",2,44],"22727/1000000"],
[["basic",3,2],"203269/500000"],
[["refined",3,46],"4063/1000000"],
[["refined",4,56],"279/15625"]
]],
["1/4","3/8",[
[["z2",-1,2],"99999/200000"],
[["refined",0,38],"5263/200000"],
[["refined",1,38],"5263/200000"],
[["refined",2,40],"24999/1000000"],
[["basic",3,2],"198741/500000"],
[["refined",3,42],"4881/1000000"],
[["refined",4,50],"19999/1000000"]
]],
["3/8","1/2",[
[["z2",-1,2],"99999/200000"],
[["refined",0,34],"29411/1000000"],
[["refined",1,34],"29411/1000000"],
[["refined",2,36],"27777/1000000"],
[["basic",3,2],"384593/1000000"],
[["refined",3,38],"6073/1000000"],
[["refined",4,44],"22727/1000000"]
]],
["1/2","5/8",[
[["z2",-1,2],"99999/200000"],
[["refined",0,30],"8333/250000"],
[["refined",1,32],"31249/1000000"],
[["refined",2,32],"31249/1000000"],
[["basic",3,2],"2983/8000"],
[["refined",3,34],"7477/1000000"],
[["refined",4,42],"23809/1000000"]
]],
["5/8","3/4",[
[["z2",-1,2],"99999/200000"],
[["refined",0,28],"35713/1000000"],
[["refined",1,28],"35713/1000000"],
[["refined",2,28],"35713/1000000"],
[["refined",3,32],"31249/1000000"],
[["basic",4,2],"354683/1000000"],
[["refined",4,42],"6919/1000000"]
]],
["3/4","7/8",[
[["z2",-1,2],"99999/200000"],
[["refined",0,24],"20833/500000"],
[["refined",1,24],"20833/500000"],
[["refined",2,26],"38461/1000000"],
[["refined",3,32],"31249/1000000"],
[["basic",4,2],"339299/1000000"],
[["refined",4,42],"1913/250000"]
]],
["7/8","1",[
[["z2",-1,2],"135579/1000000"],
[["z4_refined",-1,4],"182207/1000000"],
[["refined",0,22],"22727/500000"],
[["refined",1,24],"20833/500000"],
[["refined",2,26],"38461/1000000"],
[["basic",3,2],"6767/250000"],
[["refined",3,32],"29557/1000000"],
[["basic",4,2],"99999/200000"]
]]
]}
,
"15":{
"phi":[
"5878239/50000000","2280537/20000000","10847441/100000000",
"1980311/20000000","4095097/50000000"],
"path":[14, 13, 12, 11, 11],"v":null,
"pieces":[
["0","1/8",[
[["z2",-1,2],"99999/200000"],
[["refined",0,44],"22727/1000000"],
[["refined",1,44],"22727/1000000"],
[["refined",2,46],"10869/500000"],
[["basic",3,2],"206021/500000"],
[["refined",3,50],"1759/500000"],
[["refined",4,58],"17241/1000000"]
]],
["1/8","1/4",[
[["z2",-1,2],"99999/200000"],
[["refined",0,40],"24999/1000000"],
[["refined",1,40],"24999/1000000"],
[["refined",2,42],"23809/1000000"],
[["basic",3,2],"16101/40000"],
[["refined",3,44],"443/100000"],
[["refined",4,52],"1923/100000"]
]],
["1/4","3/8",[
[["z2",-1,2],"99999/200000"],
[["refined",0,36],"27777/1000000"],
[["refined",1,38],"5263/200000"],
[["refined",2,38],"5263/200000"],
[["basic",3,2],"392471/1000000"],
[["refined",3,40],"84/15625"],
[["refined",4,46],"10869/500000"]
]],
["3/8","1/2",[
[["z2",-1,2],"99999/200000"],
[["refined",0,32],"31249/1000000"],
[["refined",1,34],"29411/1000000"],
[["refined",2,34],"29411/1000000"],
[["basic",3,2],"378153/1000000"],
[["refined",3,36],"423/62500"],
[["refined",4,40],"24999/1000000"]
]],
["1/2","5/8",[
[["z2",-1,2],"99999/200000"],
[["refined",0,28],"35713/1000000"],
[["refined",1,30],"8333/250000"],
[["refined",2,30],"8333/250000"],
[["basic",3,2],"9029/25000"],
[["refined",3,32],"8677/1000000"],
[["refined",4,36],"27777/1000000"]
]],
["5/8","3/4",[
[["z2",-1,2],"99999/200000"],
[["refined",0,24],"20833/500000"],
[["refined",1,26],"38461/1000000"],
[["refined",2,26],"38461/1000000"],
[["refined",3,28],"35713/1000000"],
[["basic",4,2],"67323/200000"],
[["refined",4,36],"2269/250000"]
]],
["3/4","7/8",[
[["z2",-1,2],"99999/200000"],
[["refined",0,20],"49999/1000000"],
[["refined",1,20],"49999/1000000"],
[["refined",2,22],"22727/500000"],
[["refined",3,26],"38461/1000000"],
[["basic",4,2],"152631/500000"],
[["refined",4,36],"5409/500000"]
]],
["7/8","1",[
[["z2",-1,2],"841/8000"],
[["z4_refined",-1,4],"98717/500000"],
[["refined",0,18],"27777/500000"],
[["refined",1,20],"49999/1000000"],
[["refined",2,22],"22727/500000"],
[["basic",3,2],"2157/250000"],
[["refined",3,26],"37797/1000000"],
[["basic",4,2],"249997/500000"]
]]
]}
,
"16":{
"phi":[
"5944601/50000000","5776809/50000000","1376757/12500000",
"2514103/25000000","4096319/50000000"],
"path":[15, 14, 13, 12, 11],"v":null,
"pieces":[
["0","1/8",[
[["z2",-1,2],"99999/200000"],
[["refined",0,44],"22727/1000000"],
[["refined",1,44],"22727/1000000"],
[["refined",2,46],"10869/500000"],
[["basic",3,2],"411889/1000000"],
[["refined",3,48],"3671/1000000"],
[["refined",4,58],"17241/1000000"]
]],
["1/8","1/4",[
[["z2",-1,2],"99999/200000"],
[["refined",0,40],"24999/1000000"],
[["refined",1,40],"24999/1000000"],
[["refined",2,42],"23809/1000000"],
[["basic",3,2],"16101/40000"],
[["refined",3,44],"443/100000"],
[["refined",4,52],"1923/100000"]
]],
["1/4","3/8",[
[["z2",-1,2],"99999/200000"],
[["refined",0,36],"27777/1000000"],
[["refined",1,38],"5263/200000"],
[["refined",2,38],"5263/200000"],
[["basic",3,2],"392471/1000000"],
[["refined",3,40],"84/15625"],
[["refined",4,46],"10869/500000"]
]],
["3/8","1/2",[
[["z2",-1,2],"99999/200000"],
[["refined",0,32],"31249/1000000"],
[["refined",1,34],"29411/1000000"],
[["refined",2,34],"29411/1000000"],
[["basic",3,2],"189707/500000"],
[["refined",3,36],"3349/500000"],
[["refined",4,42],"23809/1000000"]
]],
["1/2","5/8",[
[["z2",-1,2],"99999/200000"],
[["refined",0,30],"8333/250000"],
[["refined",1,30],"8333/250000"],
[["refined",2,30],"8333/250000"],
[["basic",3,2],"3637/10000"],
[["refined",3,32],"4259/500000"],
[["refined",4,36],"27777/1000000"]
]],
["5/8","3/4",[
[["z2",-1,2],"99999/200000"],
[["refined",0,26],"38461/1000000"],
[["refined",1,26],"38461/1000000"],
[["refined",2,26],"38461/1000000"],
[["refined",3,28],"35713/1000000"],
[["basic",4,2],"340009/1000000"],
[["refined",4,36],"1111/125000"]
]],
["3/4","7/8",[
[["z2",-1,2],"99999/200000"],
[["refined",0,22],"22727/500000"],
[["refined",1,22],"22727/500000"],
[["refined",2,22],"22727/500000"],
[["refined",3,26],"38461/1000000"],
[["basic",4,2],"314887/1000000"],
[["refined",4,36],"10283/1000000"]
]],
["7/8","1",[
[["z2",-1,2],"106167/1000000"],
[["z4_refined",-1,4],"196913/1000000"],
[["refined",0,20],"49999/1000000"],
[["refined",1,20],"49999/1000000"],
[["refined",2,22],"22727/500000"],
[["basic",3,2],"7041/500000"],
[["refined",3,26],"37377/1000000"],
[["basic",4,2],"249997/500000"]
]]
]}
,
"17":{
"phi":[
"5998369/50000000","11673903/100000000","1392129/12500000",
"10234991/100000000"],
"path":[16, 15, 14, 14],"v":null,
"pieces":[
["0","1/8",[
[["z2",-1,2],"99999/200000"],
[["basic",0,2],"85661/200000"],
[["refined",0,42],"3413/1000000"],
[["refined",1,42],"23809/1000000"],
[["refined",2,44],"22727/1000000"],
[["refined",3,46],"10869/500000"]
]],
["1/8","1/4",[
[["z2",-1,2],"99999/200000"],
[["basic",0,2],"52587/125000"],
[["refined",0,38],"4173/1000000"],
[["refined",1,38],"5263/200000"],
[["refined",2,40],"24999/1000000"],
[["refined",3,42],"23809/1000000"]
]],
["1/4","3/8",[
[["z2",-1,2],"99999/200000"],
[["basic",0,2],"407621/1000000"],
[["refined",0,32],"5773/1000000"],
[["refined",1,34],"29411/1000000"],
[["refined",2,34],"29411/1000000"],
[["refined",3,36],"27777/1000000"]
]],
["3/8","1/2",[
[["z2",-1,2],"99999/200000"],
[["basic",0,2],"389739/1000000"],
[["refined",0,28],"63/8000"],
[["refined",1,28],"35713/1000000"],
[["refined",2,30],"8333/250000"],
[["refined",3,30],"8333/250000"]
]],
["1/2","5/8",[
[["z2",-1,2],"99999/200000"],
[["basic",0,2],"358329/1000000"],
[["refined",0,22],"6439/500000"],
[["refined",1,22],"22727/500000"],
[["refined",2,24],"20833/500000"],
[["refined",3,24],"20833/500000"]
]],
["5/8","3/4",[
[["z2",-1,2],"99999/200000"],
[["basic",0,2],"130207/500000"],
[["refined",0,14],"1369/40000"],
[["refined",1,14],"71427/1000000"],
[["refined",2,14],"71427/1000000"],
[["refined",3,16],"62499/1000000"]
]],
["3/4","7/8",[
[["z2",-1,2],"233391/500000"],
[["z4",-1,4],"8303/500000"],
[["refined",0,6],"2657/40000"],
[["refined",0,8],"75179/1000000"],
[["refined",1,8],"62499/500000"],
[["refined",2,8],"62499/500000"],
[["refined",3,8],"62499/500000"]
]],
["7/8","1",[
[["z4",-1,4],"69863/1000000"],
[["z4_refined",-1,4],"180133/1000000"],
[["refined",0,4],"249997/1000000"],
[["refined",1,6],"20833/125000"],
[["refined",2,6],"20833/125000"],
[["refined",3,6],"20833/125000"]
]]
]}

}
''' )
def interpolate(raw):
    values, witnesses = [], []
    for i in range(len(raw)):
        candidates = [(raw[i], i, i)]
        candidates += [((b-i)*raw[a]/(b-a)+(i-a)*raw[b]/(b-a), a, b)
                       for a in range(i+1) for b in range(i, len(raw)) if a < b]
        value, a, b = min(candidates)
        values.append(value)
        witnesses.append((a, b))
    return values, witnesses


GRID=10**12
def resource(k, phi, height, coeff, res):
    (kind, i, order), _ = res
    total = sum(phi)
    length = 1-total
    h = k-2-k*total
    if kind == 'z2':
        assert i == -1 and order == 2 and h >= 1
        return F(1)
    if kind == 'z4':
        assert i == -1 and order == 4
        return min(height, F(2), h)
    if kind == 'z4_refined':
        assert i == -1 and order == 4
        return min(2*height-h, F(2))
    assert 0 <= i < len(phi) and order % 2 == 0
    assert 2 <= order <= 2*(len(coeff)-1)
    assert kind in ('basic','refined','cube')
    f = phi[i]
    if kind == 'basic':
        assert order == 2
        caps = [f, k-6-k*total+4*k*f, length/4, (k*length-(4*k-1)*f)/6]
        assert min(caps) >= 0
        return min(caps)
    delta = max(F(0), coeff[order//2]-order+2*k)
    cap0 = min(2*k*f, k-3-k*total+k*f, (k-6-k*total+6*k*f)/2)
    assert cap0 >= 0
    cap = height-2+2*k*f if kind == 'refined' else height/3
    return min(cap0, cap)-delta*f

def evaluate(k, curve, coeff, s, policy):
    phi = list(map(F, policy['phi']))
    path = policy['path']
    j = len(phi)
    assert len(path) == j and path[0] == s-1 and 1 <= j <= k-2
    assert all(0 <= f <= F(1,k) for f in phi)
    totals = [F(0)]
    for f in phi: totals.append(totals[-1]+f)
    lengths = [1-t for t in totals]
    shifts = [i-k*t for i,t in enumerate(totals)]
    top = F(k,2)*lengths[j]
    assert top > 1
    amplitude = 1+shifts[j]+totals[j]
    r = path[-1]
    v = policy['v'] or 2*r
    c = F(2*r,v) if policy['v'] is not None else F(1)
    a = 1-c
    gamma = 1-2*a
    assert v % 2 == 0 and 0 <= a <= F(1,2) and 0 < c <= 1
    assert gamma >= 0 and a+c == 1 and gamma+2*a == 1
    assert all(phi[i] <= lengths[i] for i in range(j))
    assert curve[v//2] >= max(F(v,2),v-k)
    assert curve[r] >= max(F(r),2*r-k)
    prefix = [1-totals[i]-(k-2)*phi[i-1] for i in range(2,j+1)] if policy['v'] is not None else []
    assert all(t >= 0 for t in prefix)
    base = (gamma+a)*amplitude+c*(v-k)*lengths[j]
    dv = max(F(0), curve[v//2]-v+k)
    dr = max(F(0), curve[r]-2*r+k)
    cn = F(2,k)*c*dv
    root = F(1,2**(k-j-2))
    major = amplitude+(2*r-k)*lengths[j]+max(F(0),F(2,k)*dr-F(1,k-j))
    ends, pieces, last = [major], [], F(0)
    for lo,hi,resources in policy['pieces']:
        lo,hi = F(lo),F(hi)
        assert lo == last and lo < hi <= 1
        last = hi
        mass = sum((F(weight) for _,weight in resources),F(0))
        loads = [F(0)]*(j+1)
        for (_,i,order),weight in resources:
            assert F(weight) > 0
            loads[i+1] += order*F(weight)
        assert mass <= 1 and all(0 <= load <= 1 for load in loads)
        endpoints = []
        for t in (lo,hi):
            ell = 1+t*(top-1)
            saving = sum((F(res[1])*resource(k,phi,ell,coeff,res) for res in resources),F(0))
            exponent = base+cn*ell-gamma*root*min(F(1),saving)
            endpoints.append(dict(ell=ell,saving=saving,exponent=exponent))
            ends.append(exponent)
        pieces.append(dict(mass=mass,loads=loads,endpoints=endpoints))
    assert last == 1
    current, backwards = max(ends), []
    for i in range(j-1,0,-1):
        rr,u = path[i-1],path[i]
        h = 2*rr-u
        assert 0 < h < 2*rr
        diagonal = 1+shifts[i]+totals[i+1]+curve[u]*lengths[i+1]
        prefactor = (curve[h]*lengths[i]+shifts[i]+totals[i]+(2*u-1)*phi[i])/2
        current = prefactor+max(diagonal,current)/2
        backwards.append(dict(i=i,diagonal=diagonal,prefix=prefactor,output=current))
    first = 1+phi[0]+curve[s-1]*lengths[1]
    raw = (2*s-3)*phi[0]+max(first,current)
    rounded = F((raw.numerator*GRID+raw.denominator-1)//raw.denominator,GRID)
    return dict(major=major,prefix_margins=prefix,pieces=pieces,backward=backwards,raw_lambda=raw,rounded_up=rounded)


def run(source, coefficient, parameters, rounds):
    coeff_interp = interpolate(coefficient)
    previous, history = source[:], []
    for r in range(1,rounds+1):
        inter = interpolate(previous)
        entries = [evaluate(8,inter[0],coeff_interp[0],int(s),p)
                   for s,p in sorted(parameters.items(),key=lambda kv:int(kv[0]))]
        outputs = previous[:]
        for (s,_),entry in zip(sorted(parameters.items(),key=lambda kv:int(kv[0])),entries):
            s = int(s)
            outputs[s] = min(inter[0][s],entry['rounded_up'])
        history.append(dict(iteration=r,inputs=previous,interpolation=inter,entries=entries,outputs=outputs))
        previous = outputs
    return coeff_interp,history,interpolate(previous)


source=[F(x,10**6)for x in SEEDS[8][0]]
coefficient=[F(x,10**6)for x in SEEDS[16][0]]
_,history,final=run(source,coefficient,POLICIES,12)
delta=final[0][17]-26
assert delta==F(5857852793,50000000000)
assert delta<F(2,17)
print('VERIFIED eighth-power field certificate: delta34=',delta,'gate margin=',4-34*delta)

assert delta < F(2929,25000) < F(2,17)
assert (4-34*F(2929,25000))/72 == F(23,100000)
for s in range(14,18):
    print('order',2*s,'exponent',final[0][s],'excess',final[0][s]-(2*s-8))
\end{lstlisting}

\section{Exact verification of the fifth-power table}
\label{app:quintic}
The program contains every parameter in the six-row supersolution.
All inequalities, including the seventeen-variable saving, use exact
rational arithmetic.  It requires no external data or optimizer.

\begin{lstlisting}
#!/usr/bin/env python3
"""Exact six-row fifth-power smooth-moment supersolution.
Analytic prerequisites: fundamental field differencing, auxiliary L2/L4,
sharp cubic coefficient dissection, and degree-ten mu44 <= 34.227.
Only fractions.Fraction is used for verification.
"""
from fractions import Fraction as F
import json

B=[F(0),F(1),F(2)]+list(map(F,['3.1364','4.4389','5.9254','7.5422','9.2726','11.0764']))
ROWS=[
 (3,'A',1,1,[68130]),
 (4,'A',1,2,[105600]),
 (5,'A',2,2,[136590,72270]),
 (6,'A',2,2,[151340,113110]),
 (7,'E',2,6,[163790,139365,139360]),
 (8,'E',2,6,[169960,138337,138330]),
]
MU=F('34.227')

def at(vals,x):
 x=F(x)
 i=x.numerator//x.denominator
 if i==len(vals)-1:return vals[i]
 assert 0<=i<len(vals)-1
 return vals[i]+(x-i)*(vals[i+1]-vals[i])

def row_value(vals,row):
 S,kind,j,ell,raw=row
 ph=[F(x,10**6)for x in raw]
 if kind=='A':
  assert all(0<x<=F(1,5)for x in ph)
  assert all(ph[i]>=ph[i+1]for i in range(j-1))
  cum=[F(0)]
  for x in ph:cum.append(cum[-1]+x)
  p=2**ell;b=F(p-1,p);r=S-j
  C=F(p-ell+(p-1)*j,p)+F((p-1)*(1-5),p)*cum[j]+b*at(vals,F(r)/b)*(1-cum[j])
  for i in range(j-1,0,-1):
   u=S-i-1;t=S-i+1
   D=1+i-5*cum[i]+cum[i+1]+at(vals,u)*(1-cum[i+1])
   P=(at(vals,t)*(1-cum[i])+i-4*cum[i]+(2*u-1)*ph[i])/2
   C=P+max(D,C)/2
  return max(1+ph[0]+at(vals,S-1)*(1-ph[0]),C)+(2*S-3)*ph[0]
 theta,phi,u=ph;r=ell
 assert F(1,15)<=theta<=F(1,5) and 0<phi<=theta
 assert phi<5*theta-F(1,3)
 assert 0<u<phi and u<-1-5*theta+15*phi
 assert 4*u+theta+phi<1 and 6*u+5*theta+24*phi<5
 rho=(1+F(21,22)*u+(F(1,3)+(34-MU)*theta)/22)/4
 assert rho<F(1,3)
 q=1-theta-phi;t=2*(S-1)-r
 P=(at(vals,t)*(1-theta)+1-4*theta+(2*r-1)*phi)/2
 D=2-4*theta+phi+at(vals,r)*q
 C=3-4*(theta+phi)+max(-rho+at(vals,r)*q,(-F(5,4)+F(3,4)*at(vals,F(4*r,3)))*q)
 return max(1+at(vals,S-1)+(2*S-2-at(vals,S-1))*theta,(2*S-3)*theta+P+max(D,C)/2)

def verify():
 outs=[row_value(B,row)for row in ROWS]
 for row,out in zip(ROWS,outs):assert out<B[row[0]],(row[0],float(out),float(B[row[0]]))
 contraction=F(94,100)
 assert all(1-F(row[4][0],10**6)<contraction for row in ROWS)
 # Every branch has total variable coefficient <=1-theta;
 # see the accompanying proof. All coefficients are nonnegative.
 Delta=B[8]-11
 assert Delta<F(1,13)
 saving=(1-13*Delta)/14
 assert saving>0 and 2-12*Delta>0
 print(json.dumps(dict(status='VERIFIED',rows=[dict(s=row[0],kind=row[1],parameters=row[4],bound=str(B[row[0]]),output=str(out),slack=str(B[row[0]]-out))for row,out in zip(ROWS,outs)],contraction=str(contraction),joint16=str(B[8]),Delta16=str(Delta),minor_saving=str(saving),bootstrap18_margin=str(2-12*Delta)),indent=2))
if __name__=='__main__':verify()
\end{lstlisting}

\section{Exact verification of the auxiliary degree-ten table}
\label{app:decic}
The degree-ten program checks all twenty supersolution rows, their
interpolations, and their parameter conditions using exact fractions.
Its contraction argument is proved in Section~\ref{sec:tenth-moment}.
The field estimates used by the program are established in that section;
the program does not presume rational-only smooth moments over $K$.

\begin{lstlisting}
#!/usr/bin/env python3
"""Exact degree-ten smooth-moment certificate; no external data or optimizer.
Checks a supersolution for the specified analytic recurrences.
The contraction proof, not this program, permits initialization from Hua.
"""
from fractions import Fraction as Q
k=10
D=1000000
b=[Q(x,1000) for x in [0, 1000, 2000, 3059, 4187, 5396, 6677, 8027, 9445, 10925, 12468, 14068, 15722, 17425, 19174, 20962, 22786, 24642, 26522, 28424, 30344, 32279, 34227]]
rows=[(3, 'A', 1, 0, [14055]), (4, 'A', 2, 0, [43306]), (5, 'A', 1, 0, [54419, 4144]), (6, 'B', 2, 6, [60561, 37484]), (7, 'A', 2, 0, [65402, 52459]), (8, 'B', 2, 7, [69569, 61371, 36047]), (9, 'A', 2, 0, [72921, 66690, 51503]), (10, 'A', 2, 0, [76424, 71431, 62513, 36792]), (11, 'A', 2, 0, [79297, 75107, 68187, 52487]), (12, 'A', 2, 0, [82168, 78422, 73311, 63489, 37614]), (13, 'A', 2, 0, [84622, 81317, 76802, 69403, 52369]), (14, 'B', 2, 11, [86978, 83979, 80265, 74636, 63970, 36608]), (15, 'A', 2, 0, [88956, 86310, 82900, 78089, 69494, 50134]), (16, 'A', 2, 0, [90880, 88512, 85044, 80465, 75297, 62716]), (17, 'M', 0, 0, [92576, 90641, 88425, 85972, 83083, 79668, 75000]), (18, 'M', 0, 0, [93766, 92184, 90188, 87784, 84903, 81007, 74999]), (19, 'M', 0, 0, [94865, 93377, 91689, 89425, 86442, 82205, 75000]), (20, 'M', 0, 0, [95784, 94467, 92853, 90807, 87806, 83166, 74999]), (21, 'M', 0, 0, [96546, 95377, 93905, 91877, 88947, 84028, 74999]), (22, 'M', 0, 0, [97173, 96131, 94778, 92829, 89806, 84712, 74999])]
def at(t):
 t=Q(t);assert 0<=t<=22;i=t.numerator//t.denominator
 return b[i] if t==i else b[i]+(t-i)*(b[i+1]-b[i])
slacks=[]
for s,kind,ell,t,nums in rows:
 ph=[Q(x,D)for x in nums];j=len(ph);Phi=[Q(0)]
 for x in ph:Phi.append(Phi[-1]+x)
 q=[1-x for x in Phi];u=s-j
 assert all(0<x<Q(1,k)for x in ph)
 assert all(ph[i]>=ph[i+1]for i in range(j-1))
 assert ph[0]>Q(1,100)
 def kap(l):
  pp=2**l;return pp-l+(pp-1)*(j+(1-k)*Phi[j])
 if kind=='M':
  assert j==k-3 and k*q[j]>2
  C=1+j+(1-k)*Phi[j]+max(-Q(1,4)+q[j]*at(u),
       q[j]*(-Q(k,4)+Q(3,4)*at(Q(4,3)*u)))
 else:
  assert ell in(1,2) and 1<=j<=6
  for I in range(3,j+1):assert Phi[I]+k*(ph[I-2]+ph[I-1])<2
  # BCK applied in the squared variables for each residual even degree>=6.
  if (k-j)%2==0 and k-j>=6:assert Phi[j]<Q(1,3)
  if ell==2 and (k-j-1)%2==0 and k-j-1>=6:assert Phi[j]<Q(1,3)
  pp=2**ell
  if kind=='A':
   c=1-Q(1,pp);C=kap(ell)/pp+c*q[j]*at(Q(u)/c)
  else:
   c=Q(u,t);aa=2*(1-c)-Q(2,pp);bb=Q(2,pp)-(1-c)
   assert min(aa,bb,c)>=0 and aa+bb+c==1 and aa*(pp//2)+bb*pp==1
   C=aa*kap(ell-1)+bb*kap(ell)+c*q[j]*at(t)
 for i in range(j-1,0,-1):
  ui=s-i-1;ti=s-i+1
  Di=1+i-k*Phi[i]+Phi[i+1]+q[i+1]*at(ui)
  Ei=(q[i]*at(ti)+i+(1-k)*Phi[i]+(2*ui-1)*ph[i])/2
  C=Ei+max(Di,C)/2
 out=(2*s-3)*ph[0]+max(1+ph[0]+q[1]*at(s-1),C)
 assert s<out<b[s]<2*s,(s,out,b[s])
 slacks.append(b[s]-out)
 print(s,kind,'output=',float(out),'upper=',float(b[s]),'exact slack=',b[s]-out)
assert b[22]<34+Q(1,3)
print('VERIFIED: all20rows; lambda22<=',b[22],'; contraction<=99/100; minslack=',min(slacks))
\end{lstlisting}
\section{Exact verification of the sixth-power table}
\label{app:sixth}
The following Python 3 program is self-contained. Its assertions use exact
rational arithmetic; decimal conversion is used only for display. The
contraction estimate is proved in Section~\ref{sec:g6-certificate}.

\begin{lstlisting}
"""Exact replay of the eight sixth-power exponent inequalities.

This verifies finite rational arithmetic only.  It does not establish the
number-field efficient-differencing or auxiliary-moment estimates.
"""
from fractions import Fraction as F
from pathlib import Path
import json

ROWS = [
    (3, [], (1, 1), F(8), [4545]),
    (4, [], (2, 2), F(8), [8333]),
    (5, [5], (1, 2), F(10), [10673, 5080]),
    (6, [6], (2, 2), F(32, 3), [11855, 8751]),
    (7, [7, 6], (1, 2), F(12), [12981, 10763, 5551]),
    (8, [8, 7], (2, 2), F(40, 3), [13784, 12076, 8562]),
    (9, [9, 8, 8], (1, 1), F(16), [14593, 13194, 10858, 1059]),
    (10, [10, 9, 9], (1, 2), F(18), [15303, 14283, 12101, 6908]),
]
B = {r: F(v) for r, v in zip(range(6, 21, 2),
     ['3.0915','4.334','5.726','7.233','8.852','10.562','12.356','14.220'])}

def mu(x, values):
    if x == 4:
        return F(2)
    if x == F(32, 3):
        return (2*values[10] + values[12])/3
    if x == F(40, 3):
        return (values[12] + 2*values[14])/3
    return values[int(x)]

def update(row, values):
    S, path, (l1, l2), u, integers = row
    phi = [F(0)] + [F(z, 100000) for z in integers]
    d = len(integers)
    Phi = [sum(phi[:j+1]) for j in range(d+1)]
    q = [1-v for v in Phi]
    s = [None, S-1]
    for t in path:
        s.append(2*s[-1]-t)
    c = 2*s[d]/u
    if l1 == l2:
        a, b = F(1, 2**l1), F(0)
    else:
        b = (1-2**l1*(1-c))/(2**l2-2**l1)
        a = 1-b-c
    assert a+b+c == 1 and 2**l1*a + 2**l2*b == 1
    assert min(a,b,c) >= 0
    def A(l):
        return (2**l-l+(2**l-1)*d-5*(2**l-1)*Phi[d]
                +(Phi[d] if d+l == 6 else 0))
    T = a*A(l1)+b*A(l2)+c*q[d]*mu(u, values)
    for j in range(d-1, 0, -1):
        t = path[j-1]
        p = 2*s[j]-t
        diag = 1+j-6*Phi[j]+Phi[j+1]+mu(2*p,values)*q[j+1]
        T = (mu(2*t,values)*q[j]+j-5*Phi[j]
             +(2*p-1)*phi[j+1]+max(diag,T))/2
    outer_diag = 1+2*(S-1)*phi[1]+mu(2*(S-1),values)*q[1]
    candidate = max(outer_diag,(2*S-3)*phi[1]+T)
    side = [Phi[I]+6*(phi[I-1]+phi[I]) for I in range(3,d+1)]
    assert all(z < 2 for z in side)
    return candidate, side, (a,b,c)

def main():
    ledger=[]
    largest=F(0)
    for row in ROWS:
        r=2*row[0]
        result, side, weights=update(row,B)
        slack=B[r]-result
        assert slack > 0
        largest=max([largest]+side)
        ledger.append({'order':r,'b':str(B[r]),'T_b':str(result),
                       'T_b_decimal':format(float(result),'.12f'),
                       'slack':str(slack),'holder_weights':list(map(str,weights)),
                       'terminal_side_conditions':list(map(str,side))})
        print(r, format(float(result),'.12f'), 'slack', str(slack))
    initial={r:F(v) for r,v in zip(range(6,21,2),
        ['3.5','5','6.75','8.5','10.25','12','13.875','15.75'])}
    assert max(initial[r]-B[r] for r in B) == F('1.53')
    # The sup-norm contraction is proved analytically branch by branch;
    # here verify the advertised finite bound after 79 simultaneous rounds.
    bound=B[20]+F('1.53')*F(23,25)**79
    assert bound < F(14)+F(2,9)
    report={'scope':'Finite exponent arithmetic, not an analytic transfer proof',
            'rows':ledger,'largest_terminal_condition':str(largest),
            'contraction_upper_bound':'23/25',
            '79_round_order20_bound':str(bound),
            '79_round_order20_decimal':float(bound),
            'target':'14+2/9'}
    Path(__file__).with_suffix('.json').write_text(json.dumps(report,indent=2)+'\n')
    print('largest condition',float(largest))
    print('79-round upper bound',float(bound))

if __name__ == '__main__':
    main()
\end{lstlisting}

\section{Exact verification of the seventh-power table}
\label{app:seventh}
This second program contains every parameter and verifies the twelve strict
supersolution inequalities and the two final analytic margins. It requires
only Python's standard library and no external data.

\begin{lstlisting}
#!/usr/bin/env python3
"""Self-contained exact verification of the twelve rows printed in the TeX.
All arithmetic is rational; no optimization library or external data is used.
This checks the numerical certificate, conditional on the stated analytic lemmas.
"""
from fractions import Fraction as Q
D=100000
b=list(map(lambda x:Q(x,D),[0,100000,200000,306493,426513,559014,
 701542,854214,1015371,1184806,1360673,1542553,1729695,1919477,2111765]))
# s, process, ell, t, phi numerators, certified slack numerator /10^8.
rows=[
(3,'A',1,0,[3196],52177),
(4,'A',2,0,[6818],8692),
(5,'B',2,5,[8699,3541],10185),
(6,'A',2,0,[9641,6937],12539),
(7,'B',2,6,[10564,8803,4059],14896),
(8,'A',2,0,[11203,9889,6902],12594),
(9,'B',2,8,[11874,10797,8946,4151],16152),
(10,'A',2,0,[12330,11453,9898,6609],13579),
(11,'A',2,0,[12805,12028,10870,8586],14177),
(12,'A',2,0,[13252,12659,10684,10188],17123),
(13,'M',0,0,[13392,12953,12116,10714],14754),
(14,'M',0,0,[13559,13070,12326,10846],15886)]
def at(t):
 t=Q(t);assert 0<=t<=14;i=t.numerator//t.denominator
 return b[i] if t==i else b[i]+(t-i)*(b[i+1]-b[i])
for s,kind,ell,t,nums,claimed in rows:
 ph=[Q(x,D)for x in nums];j=len(ph);Phi=[Q(0)]
 for x in ph:Phi.append(Phi[-1]+x)
 q=[1-x for x in Phi];u=s-j
 assert all(0<=x<=Q(1,7)for x in ph)
 assert all(ph[i]>=ph[i+1]for i in range(j-1))
 assert ph[0]>=Q(3,100)
 def kappa(l):
  p=2**l;return p-l+(p-1)*(j-6*Phi[j])
 if kind=='M':
  assert j==4 and 7*q[j]>2
  C=5-6*Phi[4]+max(-Q(1,4)+q[4]*at(u),
                     q[4]*(-Q(7,4)+Q(3,4)*at(Q(4,3)*u)))
 else:
  for I in range(3,j+1):assert Phi[I]+7*(ph[I-2]+ph[I-1])<=2
  p=2**ell
  if kind=='A':
   c=1-Q(1,p);C=kappa(ell)/p+c*q[j]*at(Q(u)/c)
  else:
   c=Q(u,t);a=2*(1-c)-Q(2,p);bb=Q(2,p)-(1-c)
   assert min(a,bb,c)>=0 and a+bb+c==1 and a*(p//2)+bb*p==1
   C=a*kappa(ell-1)+bb*kappa(ell)+c*q[j]*at(t)
 for i in range(j-1,0,-1):
  ui=s-i-1;ti=s-i+1
  Di=1+i-7*Phi[i]+Phi[i+1]+q[i+1]*at(ui)
  Ei=(q[i]*at(ti)+i-6*Phi[i]+(2*ui-1)*ph[i])/2
  C=Ei+max(Di,C)/2
 out=(2*s-3)*ph[0]+max(1+ph[0]+(1-ph[0])*at(s-1),C)
 assert s<out<b[s]<2*s
 slack=b[s]-out
 assert slack>=Q(claimed,10**8)>0
 print(s,'exact row slack =',slack)
assert b[14]==Q(422353,20000)<21+Q(59,500)
assert 24-(Q(6,56)*48+Q(25,28)*(21+Q(59,500)))==Q(1,560)
assert 4-24*Q(59,500)==Q(146,125)
print('VERIFIED: all 12 rows, contraction <=97/100, lambda14 <=21.11765.')
\end{lstlisting}

\section{Finite eighth-power certificate at high ramification}
\label{app:eighth-ghost}
The program constructs all 8192 normal forms, checks that they are distinct,
checks the eight carry relations, and verifies the image of every binary root
polynomial modulo $T^6$. The positive coefficient budget is exactly 26.
It can emit the full list of 64 coefficient witnesses, so the certificate
requires no sampling of local fields and no omitted input file.

\begin{lstlisting}
#!/usr/bin/env python3
"""Finite algebra certificate for the seven-root eighth-power chart.

Checks all 64 binary polynomials modulo T^6 and eight carry relations in
the injective combined ghost model for W4(F2[[T]])/I_(2,4,6,6).
No local-field sampling or floating-point arithmetic is used.
"""
from itertools import product
from pathlib import Path
import hashlib,json

N=(2,4,6,6)
ROOTS=((1,0,0,0,0,0),(1,0,0,0,0,1),(1,0,0,0,1,0),
       (1,0,0,1,0,0),(1,0,1,0,0,0),(1,1,0,0,0,0),
       (0,1,0,0,0,0),(0,1,1,0,0,0))
ORDERS=(16,2,2,2,2,4,4,2)

def mul(a,b):
    z=[0]*6
    for i,x in enumerate(a):
        for j,y in enumerate(b[:6-i]):z[i+j]=(z[i+j]+x*y)%16
    return z

def ghost(root):
    p=list(root);out=[]
    for j,n in enumerate(N):
        out.extend((x*(1<<(3-j)))%16 for x in p[:n])
        p=mul(p,p)
    return tuple(out)

def add(x,y):return tuple((a+b)%16 for a,b in zip(x,y))
def scale(c,x):return tuple(c*a%16 for a in x)

def main():
    gens=[ghost(x) for x in ROOTS]
    normal={tuple([0]*sum(N)):()};carries=[]
    for i,(g,o) in enumerate(zip(gens,ORDERS)):
        carry=scale(o,g)
        assert carry in normal, ('carry failure',i,carry)
        carries.append({'generator_index':i,'multiple':o,'earlier_coefficients':normal[carry]})
        expanded={}
        for v,coeffs in normal.items():
            for c in range(o):
                w=add(v,scale(c,g))
                assert w not in expanded,('normal form collision',i)
                expanded[w]=coeffs+(c,)
        normal=expanded
    witnesses=[]
    for root in product(range(2),repeat=6):
        v=ghost(root);assert v in normal,('power image missing',root)
        c=normal[v]
        check=tuple([0]*sum(N))
        for n,g in zip(c,gens):check=add(check,scale(n,g))
        assert check==v
        witnesses.append({'root_coefficients':root,'positive_generator_coefficients':c})
    out={'weights':N,'generator_roots':ROOTS,'generator_vectors':gens,
         'cyclic_quotient_orders':ORDERS,'carry_relations':carries,
         'normal_form_count':len(normal),'positive_coefficient_budget':sum(o-1 for o in ORDERS),
         'all_binary_roots_checked':len(witnesses),'witnesses':witnesses,
         'status':'all exact identities and carry relations passed'}
    path=Path(__file__).with_suffix('.json')
    path.write_text(json.dumps(out,indent=2)+'\n')
    print(json.dumps({k:out[k] for k in ('normal_form_count','positive_coefficient_budget','all_binary_roots_checked','status')}))
    print('carry relations',carries)
    print('sha256',hashlib.sha256(path.read_bytes()).hexdigest())

if __name__=='__main__':main()
\end{lstlisting}

\section{Finite eighth-power certificate in the tame quintic case}
\label{app:eighth-tame}
The exact quotient ring, all 64 root classes, and the multiplication relation
$\pi^5=2$ are specified in the program. The positive sumset has 8192 classes
and diameter 21. Subtracting one from the target and appending a unit root
therefore gives primitive width 22. The exhaustive breadth-first search
also verifies parent representations and closure under all generators.

\begin{lstlisting}
#!/usr/bin/env python3
"""Exact primitive eighth-power certificate in Z2[pi], pi^5=2.

The target quotient is A/(16*pi) = Z/32 + (Z/16)pi + ... + (Z/16)pi^4.
Eighth powers depend only on roots modulo2*pi, so all64 root classes
are checked. Ordinary width21 upgrades to primitive width22 by applying
the result to target minus1 and appending the unit root1.
"""
from collections import Counter,deque
from itertools import product
from pathlib import Path
import hashlib,json

MOD=(32,16,16,16,16)

def mul(x,y):
    z=[0]*9
    for i,a in enumerate(x):
        for j,b in enumerate(y):z[i+j]+=a*b
    for i in range(8,4,-1):z[i-5]+=2*z[i]
    return tuple(z[i]%MOD[i] for i in range(5))

def add(x,y):return tuple((a+b)%m for a,b,m in zip(x,y,MOD))

def main():
    roots=[(a,)+b for a in range(4) for b in product(range(2),repeat=4)]
    image={}
    for r in roots:
        x=r
        for _ in range(3):x=mul(x,x)
        image.setdefault(x,r)
    generators=sorted(image)
    zero=(0,)*5;dist={zero:0};parent={};queue=deque([zero])
    while queue:
        x=queue.popleft()
        for j,g in enumerate(generators):
            y=add(x,g)
            if y not in dist:
                dist[y]=dist[x]+1;parent[y]=(x,j);queue.append(y)
    for y,(x,j) in parent.items():
        assert add(x,generators[j])==y and dist[y]==dist[x]+1
    for x in dist:
        for g in generators:assert add(x,g) in dist
    assert len(dist)==8192 and max(dist.values())==21
    serialized=json.dumps(sorted((list(x),d) for x,d in dist.items()),separators=(',',':')).encode()
    result={'minimal_polynomial':'X^5-2','root_modulus':'2*pi','target_modulus':'16*pi',
            'root_classes_checked':len(roots),'distinct_power_images':len(generators),
            'generators':[{'value':g,'root':image[g]} for g in generators],
            'generated_group_order':len(dist),'ordinary_width':max(dist.values()),
            'primitive_bound':max(dist.values())+1,'distance_histogram':dict(sorted(Counter(dist.values()).items())),
            'distance_table_sha256':hashlib.sha256(serialized).hexdigest(),
            'checks':['all64rootclasses','exactringarithmetic','allparentwitnesses','closureunderallgenerators'],
            'status':'passed'}
    path=Path(__file__).with_suffix('.json');path.write_text(json.dumps(result,indent=2)+'\n')
    print(json.dumps({k:v for k,v in result.items() if k!='generators'},indent=2))

if __name__=='__main__':main()
\end{lstlisting}


\begin{thebibliography}{99}
\bibitem{BCK2025}
G. Binyamini, R. Cluckers and F. Kato,
\emph{Sharp bounds for the number of rational points on algebraic curves and
dimension growth, over all global fields}, Proc. London Math. Soc. (3)
\textbf{130} (2025), e70016. Corollary 12 and the height definition preceding (28).
\href{https://doi.org/10.1112/plms.70016}{doi:10.1112/plms.70016}.

\bibitem{Birch1961}
B. J. Birch, \emph{Waring's problem in algebraic number fields},
Proc. Cambridge Philos. Soc. \textbf{57} (1961), 449--459.
\href{https://doi.org/10.1017/S0305004100035490}{doi:10.1017/S0305004100035490}.

\bibitem{BW2023}
J. Br\"udern and T. D. Wooley, \emph{On Waring's problem for larger powers},
J. reine angew. Math. \textbf{805} (2023), 115--142. Theorem 4.2.
\href{https://doi.org/10.1515/crelle-2023-0072}{doi:10.1515/crelle-2023-0072}.

\bibitem{Cluckers2011}
R. Cluckers, \emph{Analytic van der Corput lemma for $p$-adic and
$\mathbb F_q((t))$ oscillatory integrals, singular Fourier transforms,
and restriction theorems}, Exposition. Math. \textbf{29} (2011), 371--386.
Proposition 3.3.
\href{https://doi.org/10.1016/j.exmath.2011.06.004}{doi:10.1016/j.exmath.2011.06.004}.

\bibitem{Davidson1999}
M. Davidson, \emph{On Waring's problem in number fields},
J. London Math. Soc. (2) \textbf{59} (1999), 435--447.
\href{https://doi.org/10.1112/S0024610799007309}{doi:10.1112/S0024610799007309}.

\bibitem{Hormander2003}
L. H\"ormander, \emph{The Analysis of Linear Partial Differential Operators I},
2nd ed., Classics in Mathematics, Springer, Berlin, 2003, Section 7.7.
\href{https://doi.org/10.1007/978-3-642-61497-2}{doi:10.1007/978-3-642-61497-2}.

\bibitem{Joly1970}
J.-R. Joly, \emph{Sommes de puissances $d$-i\`emes dans un anneau commutatif},
Acta Arith. \textbf{17} (1970), 37--114. Proposition 7.16 and
Theorem 7.34, pp. 95 and 105--106.
\href{https://doi.org/10.4064/aa-17-1-37-114}{doi:10.4064/aa-17-1-37-114}.

\bibitem{Ramanujam1963}
C. P. Ramanujam, \emph{Sums of $m$-th powers in $p$-adic rings},
Mathematika \textbf{10} (1963), 137--146. Lemma 1.
\href{https://doi.org/10.1112/S002557930000406X}{doi:10.1112/S002557930000406X}.

\bibitem{VW1994}
R. C. Vaughan and T. D. Wooley,
\emph{Further improvements in Waring's problem, II: Sixth powers},
Duke Math. J. \textbf{76} (1994), 683--710.
\href{https://doi.org/10.1215/S0012-7094-94-07626-6}{doi:10.1215/S0012-7094-94-07626-6}.

\bibitem{VW1995}
R. C. Vaughan and T. D. Wooley, \emph{Further improvements in Waring's problem},
Acta Math. \textbf{174} (1995), 147--240.
\href{https://doi.org/10.1007/BF02392467}{doi:10.1007/BF02392467}.

\bibitem{VW2000}
R. C. Vaughan and T. D. Wooley,
\emph{Further improvements in Waring's problem, IV: Higher powers},
Acta Arith. \textbf{94} (2000), 203--285.

\bibitem{Weil1948}
A. Weil, \emph{On some exponential sums},
Proc. Natl. Acad. Sci. USA \textbf{34} (1948), 204--207.
\href{https://doi.org/10.1073/pnas.34.5.204}{doi:10.1073/pnas.34.5.204}.

\bibitem{WooleyThesis}
T. D. Wooley, \emph{On simultaneous additive equations and Waring's problem},
Ph.D. thesis, University of London, 1990, Chapter 6, especially pp. 230--243.
\href{https://www.math.purdue.edu/~twooley/publ/1990\%20phd\%20thesis.pdf}{Author's thesis PDF}.

\bibitem{Wooley1992}
T. D. Wooley, \emph{Large improvements in Waring's problem},
Ann. of Math. (2) \textbf{135} (1992), 131--164.
\href{https://doi.org/10.2307/2946566}{doi:10.2307/2946566}.

\bibitem{Wooley2012}
T. D. Wooley, \emph{The asymptotic formula in Waring's problem},
Int. Math. Res. Not. IMRN \textbf{2012}, no. 7, 1485--1504.
\href{https://www.math.purdue.edu/~twooley/publ/2012asympwp.pdf}{Corrected author version}.

\bibitem{WooleyCorr}
T. D. Wooley, \emph{Corrigendum to ``The asymptotic formula in Waring's problem''},
Int. Math. Res. Not. IMRN (2015).
\href{https://doi.org/10.1093/imrn/rnu264}{doi:10.1093/imrn/rnu264}.

\bibitem{Wooley2016}
T. D. Wooley, \emph{On Waring's problem for intermediate powers},
Acta Arith. \textbf{176} (2016), 241--247.
\href{https://doi.org/10.4064/aa8439-8-2016}{doi:10.4064/aa8439-8-2016}.
\href{https://www.math.purdue.edu/~twooley/publ/wpip.pdf}{Author version}.

\bibitem{Wooley2019}
T. D. Wooley, \emph{Nested efficient congruencing and relatives of Vinogradov's
mean value theorem}, Proc. London Math. Soc. (3) \textbf{118} (2019), 942--1016.
Theorem 15.1 and Corollary 15.3.
\href{https://doi.org/10.1112/plms.12204}{doi:10.1112/plms.12204}.
\href{https://www.math.purdue.edu/~twooley/publ/2017nec.pdf}{Author version}.

\bibitem{ProjectAnalytic}
Waring--Siegel project, \emph{Product-shaped major arcs and a degree-uniform
Waring--Siegel principle}, project research manuscript, 22 November 2026;
and \emph{Sixteen fourth powers over number fields}, subsequent project note.
The ordinary analytic inputs required here are reproduced in
Section~\ref{sec:ordinary-bundle}.

\bibitem{ProjectG56}
Waring--Siegel project, \emph{Waring $G(5)$ and $G(6)$ checkpoint},
supplied project research document, consulted 24 November 2026.

\bibitem{ProjectG89}
Waring--Siegel project, \emph{$G(k)$ target selection working manuscript}
and \emph{Waring $G(8)$/$G(9)$ finite-certificate checkpoint},
supplied project documents and certificate sources, consulted 24 November 2026.
The present note uses their rational numerical bounds as target data only;
it does not certify their unproved number-field transfers.
\end{thebibliography}

\end{document}
